\documentclass[11pt]{article}
\usepackage[localtheorems]{raeez-math-template}
\usepackage{amscd}

\theoremstyle{plain}
\newtheorem{theorem}{Theorem}[section]
\newtheorem{proposition}[theorem]{Proposition}
\newtheorem{lemma}[theorem]{Lemma}
\newtheorem{corollary}[theorem]{Corollary}
\theoremstyle{definition}
\newtheorem{definition}[theorem]{Definition}
\newtheorem{example}[theorem]{Example}
\newtheorem{openproblem}[theorem]{Open problem}
\theoremstyle{remark}
\newtheorem{remark}[theorem]{Remark}

\newcommand{\GammaR}{\Gamma_{\mathbb R}}
\newcommand{\Adele}{\mathbb A_{\mathbb Q}}
\newcommand{\vmLambda}{\Lambda_{\mathrm{vM}}}
\newcommand{\psiGamma}{\psi_{\Gamma}}
\newcommand{\one}{\mathbf 1}
\newcommand{\Hp}{\mathbb H}
\newcommand{\ord}{\operatorname{ord}}
\newcommand{\dd}{\,ds}
\newcommand{\im}{\operatorname{im}}
\newcommand{\bcA}{\mathcal A_{\mathrm{BC}}}
\newcommand{\bcZhat}{\widehat{\mathbb Z}}
\newcommand{\bcUnits}{\widehat{\mathbb Z}^{\times}}
\newcommand{\bcN}{\mathbb N^{\times}}
\newcommand{\bcphase}{\mathrm e}
\newcommand{\bcProb}{\operatorname{Prob}}
\newcommand{\bcTr}{\operatorname{Tr}}
\newcommand{\bcGal}{\operatorname{Gal}}


\title{Euler Factors, Determinant Covers, and Arithmetic Boundary Classes}
\author{Raeez Lorgat}
\date{}
\hypersetup{pdftitle={Euler Factors, Determinant Covers, and Arithmetic Boundary Classes},
 pdfauthor={Raeez Lorgat},
 pdfsubject={Arithmetic Euler factors, regularized determinants, and finite arithmetic cohomology}}

\begin{document}
\maketitle
\begin{abstract}
An Euler factor can be a normalized local integral, an occupation trace,
or the reciprocal of a determinant. Recovering an arithmetic complex
requires comparisons that also retain multiplication, boundary maps and
the chosen state. We construct these scalar and operator realizations
with their measures, domains and normalizations, and determine the
minimal continuation surfaces of the regularized prime determinants.
For finitely many regularization orders on a smaller positive half-plane,
the simultaneous monodromy is the finite subgroup generated by their
rational exponent vectors. Its lattice presentation determines the exact
degree and a connected normalized fibre-product realization; orders two
and three on $\operatorname{Re}s>1/5$ require twenty-four sheets.
An integral divisor bar and a torsion Hall exchange retain algebraic
operations invisible in the scalar factor. Independently, finite
arithmetic cochains give an oriented degree-three trace and compatible
boundary torsors, including coefficient primes. The archimedean
oscillator fixes the constant in the conditional Deninger determinant
formula. Positive narrow support sectors retain the full Weil functional,
while arithmetic twisting, a full positive representation and a
cohomological spectral realization remain separate construction problems.
\end{abstract}

\tableofcontents
\newpage
\section{One Euler factor and five kinds of structure}
\label{sec:opening}

For a prime number $p$, put
\[
 E_p(s)=\frac1{1-p^{-s}},\qquad
 p^{-s}=\exp(-s\log p),\qquad \operatorname{Re}s>0,
\]
where $\log p$ is real. The geometric series defining $E_p$ counts
arbitrary occupation of one mode of energy $\log p$. It also occurs in
integration over the rational local field. The question is what, beyond
this scalar value, either construction determines. An arithmetic
cohomological explanation of the zeta function must retain operations
that the value alone does not record: multiplication, local restriction,
duality, and the choice of a state.

The local integral has complex coefficients and the $p$-adic field
$\mathbb Q_p$ as its domain. Its test space consists of locally constant,
compactly supported functions, with the inductive-limit topology of
finite-dimensional spaces of functions of bounded support and fixed
local constancy. These functions have degree zero and pointwise product.
Let $\mathbb Z_p$ be the valuation ring, take $|p|_p=p^{-1}$, and
normalize multiplicative Haar measure $d^\times x$ by
$\operatorname{vol}(\mathbb Z_p^\times)=1$. The chosen test is
$\one_{\mathbb Z_p}$. Its support and this measure are essential:
changing to additive measure changes the normalization.

The occupation space is the complex Hilbert space
$\mathscr H_p=\ell^2(\mathbb N_0)$ with orthonormal vectors $e_r$.
The number operator in energy units is
\[
 N_pe_r=r\log p\,e_r,\qquad
 \mathcal D(N_p)=
 \left\{a:\sum_{r\ge0}r^2(\log p)^2|a_r|^2<\infty\right\}.
\]
Here $r$ is occupation number, not cohomological degree. The finite
sequences form a common algebraic domain. Bounded operators multiply by
composition; the Hilbert space itself has no specified multiplication.
The trace of $e^{-sN_p}$ exists exactly for $\operatorname{Re}s>0$.
The basis vector $e_0$ is the vacuum. At real inverse temperature
$\beta>0$, the equilibrium density is the normalized positive operator
$(1-p^{-\beta})e^{-\beta N_p}$.

A third carrier is the one-dimensional complex Hilbert space
$V_p=\mathbb C e_p$, with bounded operator $K_p(s)e_p=p^{-s}e_p$.
It is concentrated in degree zero. Its ordinary determinant, normalized
to one as real $s\to+\infty$, is $\det(1-K_p(s))=1-p^{-s}$.
Thus the comparison of these three constructions is the scalar identity
\begin{equation}
 \int_{\mathbb Q_p^\times}\one_{\mathbb Z_p}(x)|x|_p^s\,d^\times x
 =\operatorname{Tr}(e^{-sN_p})
 =\det(1-K_p(s))^{-1}=E_p(s).
 \label{eq:opening-euler}
\end{equation}
The first equality is proved with the local-field and measure
constructions in Proposition~\ref{prop:local-tate}; the occupation sum
is the same convergent geometric series. No arrow between algebras or
complexes has been inferred from this equality.

Multiplication enters separately through effective divisors. Over a
commutative coefficient ring $R$, let $A_p=R[u_p]$, with unit $1$,
augmentation $\epsilon(u_p)=0$, discrete coefficients and algebraic
tensor products. Its elements have homological degree zero and the
weight grading $\operatorname{wt}(u_p^r)=r$; multiplication adds weights.
The bar construction is a complex built from this multiplication and
augmentation. Over $\mathbb C$, completing the monomials to an
orthonormal basis gives the occupation space above, and multiplication
by $u_p$ gives the unilateral shift. The square-zero algebra
$\mathbb C\oplus\bigoplus_{r>0}\mathbb C e_r$, with the same unit,
augmentation, weights and energy operator, has the same heat trace
but a different bar homology (Proposition~\ref{prop:two-bars}).
Consequently the scalar factor does not determine the multiplication
needed for arithmetic factorization.

An arithmetic cohomology class has a fifth carrier. For a totally
imaginary number field $F$, set $X=\operatorname{Spec}\mathcal O_F$.
For a positive integer $n$, assume $F$ contains a chosen primitive
$n$th root of unity. Ordinary sheaf cohomology on the small \'etale
site of $X$ has finite discrete coefficient ring $\mathbb Z/n$,
cohomological grading and cup product. A finite-group torsor on $X$
pulls a group-cohomology class back to degree three; the chosen root
normalizes its trace in $\tfrac1n\mathbb Z/\mathbb Z$. Removing a
nonempty finite set of places containing those above $n$ gives local
cochains and finite gluing maps. These are arithmetic operations, but
their coefficients and grading differ from those of the complex
Euler factor. A determinant comparison would require a realization
functor, an operator and a trace-compatible map; none is specified by
the existence of the degree-three class.

The analytic comparison can nevertheless be made exact. On the Hilbert
space with one basis vector for each prime, the diagonal $K(s)$ has
singular values $p^{-\operatorname{Re}s}$. Its $m$th regularized
determinant is defined precisely when $m\operatorname{Re}s>1$.
Analytic continuation preserves the normalized germ but may make it
multivalued. The logarithmic derivative defines a period character;
its kernel determines the smallest pointed surface on which the germ
is single-valued. For several orders the character is vector-valued.
On $\operatorname{Re}s>1/5$, orders two and three have individual
degrees six and twelve, but their joint surface has degree twenty-four.
This finite example distinguishes simultaneous continuation from
replacing several degree conditions by their least common multiple.

Write $\zeta(s)=\sum_{n\ge1}n^{-s}$ on $\operatorname{Re}s>1$
and $\GammaR(s)=\pi^{-s/2}\Gamma(s/2)$, where
$\Gamma(z)=\int_0^\infty e^{-u}u^{z-1}\,du$ for $\operatorname{Re}z>0$.
The real Gaussian factor has a different operator explanation.
An equally spaced additive spectrum gives a reciprocal Gamma
determinant, with spacing and scale determining its constant. At
spacing $2$ and scale $2\pi$ the value is $\sqrt2/\GammaR(s)$.
Together with the order-one prime determinant $1/\zeta(s)$ this
recovers $2^{-1/2}\GammaR(s)\zeta(s)$. This is an analytic identity
on specified carriers. Deninger's stronger question asks for
arithmetically defined cohomology and an operator whose alternating
determinant has this value. Its degree-one determinant is determined
conditionally in Theorem~\ref{thm:deninger}; the theorem does not
construct that operator or its eigenvalue multiplicities.

Two positive constructions clarify what an arithmetic comparison
could preserve. On a smooth projective curve over a finite field,
Hall multiplication and its adjoint give an operator exchange with
the curve zeta function as coefficient, before any trace is taken.
For the Riemann zeta function, the complete Weil functional has
infinite-dimensional positive subspaces on sufficiently narrow finite
unions of intervals. These subspaces retain all polar, prime-power
and archimedean terms. Their supports are not stable under convolution,
so they do not supply the full positive representation required by
Weil's criterion. Arithmetic twisting, full positivity and
cohomological spectral realization are therefore distinct targets,
with the comparison maps stated where their carriers are constructed.

The inner-product convention is part of every operator comparison.
The occupation and oscillator spaces use an inner product linear in
the first argument. The Hall pairing and the positive Weil pairing use
one linear in the second.
If $\langle x,y\rangle_1$ is the first convention, put
$\langle x,y\rangle_2=\langle y,x\rangle_1$ to pass to the second.
The adjoint operator is unchanged: conjugating
$\langle Ax,y\rangle_1=\langle x,A^*y\rangle_1$ gives the
corresponding identity in the second convention. Rank-one formulas
must also be transported: $h\mapsto\langle h,g\rangle_1 f$ becomes
$h\mapsto\langle g,h\rangle_2 f$. This fixes the convention before
the equilibrium, Hall and finite-rank trace calculations are compared.

\section{The scalar Tate calculation}
\label{sec:arithmetic-tate}

A local zeta integral separates into valuation shells, each weighted by
a power of the prime. The resulting geometric series contributes to a
single integral on the rational ideles. Its continuation comes from
Fourier summation on a rational lattice. The finite unit balls
select the integers from that lattice, and the real Gaussian supplies the
theta transformation. We construct the fields, measures, and lattice before
making this argument.

All vector spaces of test functions below have complex coefficients,
pointwise multiplication, and cohomological degree zero. At a finite place,
the locally constant functions on a fixed compact support with a fixed
open subgroup of constancy form a finite-dimensional space; their filtered
union has the locally convex inductive-limit topology. At the real place
we use the Schwartz topology, defined by the seminorms
\(\sup_x|x^a f^{(b)}(x)|\), \(a,b\geq0\).
The adelic test space below is the inductive limit of its finite restricted
tensor products with these topologies. These test algebras on the
noncompact fields and on the adeles have no pointwise multiplicative unit. The displayed integrals are scalar
linear functionals on their stated convergence domains; they are not
asserted to preserve pointwise products.
Integrals use positive Borel measures. A complex power of a positive real
number means \(a^s=\exp(s\log a)\), with the real logarithm.
We use countable product measures, Tonelli's theorem, dominated convergence,
and elementary real and complex analysis. The gamma function initially means
only Euler's integral
\[
 \Gamma(z)=\int_0^\infty e^{-u}u^{z-1}\,du,
 \qquad \operatorname{Re}z>0.
\]
No meromorphic continuation of \(\Gamma\) is needed in the proof.

\subsection{The local field and its unit ball}
\label{sec:adeles-local-field}

Fix a prime integer \(p\). Define the ring
\[
 \mathbb Z_p=\varprojlim_{n\geq1}\mathbb Z/p^n\mathbb Z
 =\{(a_n)_{n\geq1}:a_{n+1}\equiv a_n\pmod{p^n}\}.
\]
Addition and multiplication act on each coordinate. An integer maps to its
compatible residue classes. This map is injective, since an integer divisible
by every \(p^n\) is zero. Each element has unique digits
\[
 a=\sum_{j\geq0}a_jp^j,\qquad a_j\in\{0,\ldots,p-1\}.
\]
This notation specifies residues modulo \(p^n\). It does not assert
convergence in the real numbers.

Give each finite residue ring its discrete topology and \(\mathbb Z_p\) the induced
product topology. Its basic neighborhoods are the residue classes
\(a+p^n\mathbb Z_p\). Indeed, multiplication by \(p^n\) identifies \(\mathbb Z_p\)
with the elements whose first \(n\) digits vanish.
A diagonal subsequence argument shows compactness: successively fix each
digit, and then take the diagonal subsequence. The resulting subsequence
converges in every residue coordinate. This topology is metrizable by the
distance defined below, so sequential compactness implies compactness.
Each residue class is both open and closed.

An element \(u\in\mathbb Z_p\) is a unit precisely when its residue modulo \(p\)
is nonzero. In that case its inverse in each ring \(\mathbb Z/p^n\mathbb Z\)
exists by Bezout's identity. Uniqueness of the inverse makes these inverses
compatible, so they define \(u^{-1}\in\mathbb Z_p\).
For \(a\neq0\), its first nonzero digit therefore gives a unique expression
\[
 a=p^r u,\qquad r\in\mathbb N_0,\quad u\in\mathbb Z_p^\times.
\]
The product of two such expressions has exponent equal to the sum of their
exponents. In particular, \(\mathbb Z_p\) has no zero divisors.

Define \(\mathbb Q_p=\mathbb Z_p[1/p]\), the localization of this integral domain.
Every nonzero element now has a unique expression
\begin{equation}
 x=p^r u,\qquad r\in\mathbb Z,\quad u\in\mathbb Z_p^\times.
 \label{eq:adeles-valuation-decomposition}
\end{equation}
It is invertible. Thus \(\mathbb Q_p\) is a field and contains \(\mathbb Q\).
Put \(v_p(x)=r\), \(|x|_p=p^{-r}\), \(v_p(0)=+\infty\), and
\(|0|_p=0\). Divisibility gives
\[
 |xy|_p=|x|_p|y|_p,\qquad
 |x+y|_p\leq\max\{|x|_p,|y|_p\}.
\]
Consequently \(d_p(x,y)=|x-y|_p\) is a metric. On \(\mathbb Z_p\) it gives the
residue topology above. Addition and multiplication are continuous.
Inversion is continuous on \(\mathbb Q_p^\times\), since
\[
 |x^{-1}-y^{-1}|_p=\frac{|x-y|_p}{|x|_p|y|_p}
\]
and \(|y|_p=|x|_p\) whenever \(|y-x|_p<|x|_p\).

\begin{proposition}[Completion and local compactness]
\label{prop:adeles-completion}
The field \(\mathbb Q_p\) is complete for \(d_p\), and \(\mathbb Q\) is dense
in it. The sets \(a+p^r\mathbb Z_p\), with \(a\in\mathbb Q_p\) and \(r\in\mathbb Z\),
are compact open balls and form a neighborhood basis.
Thus \(\mathbb Q_p\) is a locally compact topological field.
\end{proposition}

\begin{proof}
A Cauchy sequence is bounded: its tail lies in a ball of radius one around
one of its terms, and its remaining terms form a finite set.
Multiplication by a sufficiently large power of \(p\) places the entire
sequence in \(\mathbb Z_p\). Each residue coordinate is eventually constant.
These constants are compatible and define its limit. This proves
completeness after scaling back. Truncated digit expansions approximate
every element of \(\mathbb Z_p\) by integers. Scaling proves density of
\(\mathbb Q\) in \(\mathbb Q_p\). The balls are translates and dilates of the
compact open set \(\mathbb Z_p\). The metric inequalities identify them with the
stated neighborhood basis.
\end{proof}

The notation \(\mathbb Q_p\) therefore agrees with the metric completion
of \(\mathbb Q\) for the usual rational \(p\)-adic absolute value.
No algebraic extension or choice of a root is involved.

\subsection{Measures, characters, and one Fourier calculation}
\label{sec:adeles-local-fourier}

On \(\mathbb Z_p\), give the digits independent uniform probability distributions.
The resulting measure assigns mass \(p^{-n}\) to every residue class modulo
\(p^n\). Translation and multiplication by a unit permute these residue
classes. They therefore preserve the measure, first on cylinder sets and
then on Borel sets by uniqueness of the product measure.
Also, multiplication by \(p\) scales this measure by \(p^{-1}\).

Extend the measure to \(\mathbb Q_p=\bigcup_{m\geq0}p^{-m}\mathbb Z_p\) by
\[
 \mu_p(E)=p^m\mu_p(p^mE),\qquad E\subset p^{-m}\mathbb Z_p.
\]
On the right, \(\mu_p\) denotes the probability measure on \(\mathbb Z_p\).
The scaling identity makes these prescriptions compatible as \(m\)
increases. They define a Borel measure on the union. Write it as \(dx_p\).
Translations preserve it, and
\begin{equation}
 \operatorname{vol}(a+p^r\mathbb Z_p)=p^{-r},\qquad
 d(bx)_p=|b|_p\,dx_p\quad(b\neq0).
 \label{eq:adeles-additive-scaling}
\end{equation}
The second identity follows from the unit and power-of-\(p\) cases.
This is an additive Haar measure: it is translation invariant, finite on
compact sets, and nonzero on open sets. Regularity follows from the
countable basis of compact open balls and the regular probability measure
on each \(p^{-m}\mathbb Z_p\). Its normalization determines it uniquely, since
the masses of all balls are then forced.

For \(x\in\mathbb Q_p\), its Laurent digit expansion has only finitely many
negative powers. Define
\[
 \{x\}_p=\sum_{r\leq j<0}a_jp^j
 \in\mathbb Z[1/p]\cap[0,1),\qquad
 \psi_p(x)=\exp(-2\pi i\{x\}_p).
\]
The fractional part is the unique element of this set congruent to \(x\)
modulo \(\mathbb Z_p\). To see uniqueness, note that
\(\mathbb Z[1/p]\cap\mathbb Z_p=\mathbb Z\). Two representatives in \([0,1)\)
cannot differ by a nonzero integer.
In particular, \(\{x+y\}_p-\{x\}_p-\{y\}_p\) is an integer.
Thus \(\psi_p:(\mathbb Q_p,+)\to\{z\in\mathbb C:|z|=1\}\) is a continuous
additive character, with kernel exactly \(\mathbb Z_p\).

For a locally constant, compactly supported function \(f:\mathbb Q_p\to\mathbb C\),
define its Fourier transform by
\[
 \widehat f(y)=\int_{\mathbb Q_p}f(x)\overline{\psi_p(xy)}\,dx_p.
\]
Such functions form the Schwartz--Bruhat space \(\mathcal S(\mathbb Q_p)\).

At \(p=2\) and \(y=1/2\), the even and odd residue classes each have
mass \(1/2\). The Fourier integral of \(\one_{\mathbb Z_2}\) is therefore
\[
 \frac12\cdot1+\frac12\cdot(-1)=0.
\]
At \(y\in\mathbb Z_2\), the same integral is one. The following identity
gives the full calculation, including every scale and translate.

\begin{proposition}[Fourier transform of a ball]
\label{prop:adeles-ball-fourier}
For \(a,y\in\mathbb Q_p\) and \(r\in\mathbb Z\),
\begin{equation}
 \widehat{\one_{a+p^r\mathbb Z_p}}(y)
 =p^{-r}\overline{\psi_p(ay)}\one_{p^{-r}\mathbb Z_p}(y).
 \label{eq:adeles-ball-fourier}
\end{equation}
In particular, \(\widehat{\one_{\mathbb Z_p}}=\one_{\mathbb Z_p}\).
Moreover, \(\widehat{\widehat f}(x)=f(-x)\) for every
\(f\in\mathcal S(\mathbb Q_p)\). Thus the additive measure of unit-ball mass one
is self-dual for this Fourier convention.
\end{proposition}

\begin{proof}
For \(y\in p^{-r}\mathbb Z_p\),
the character \(z\mapsto\psi_p(zy)\) is trivial on \(p^r\mathbb Z_p\).
Its integral over that ball is \(p^{-r}\).
If \(y\notin p^{-r}\mathbb Z_p\), then \(p^ry\notin\mathbb Z_p\).
Translation by \(p^r\) multiplies the integral by the nontrivial scalar
\(\overline{\psi_p(p^ry)}\). Translation invariance forces the integral
to vanish. Translation by \(a\) proves
\eqref{eq:adeles-ball-fourier}.

Every locally constant compactly supported function is a finite linear
combination of indicators of balls of one sufficiently small radius.
Indeed, finitely many neighborhoods of constancy cover its support.
Their finite refinement by balls of a common radius has this property.
Applying the Fourier transform twice to a ball indicator gives
\[
 p^{-r}\int_{p^{-r}\mathbb Z_p}
       \overline{\psi_p((a+x)y)}\,dy_p
 =\one_{a+p^r\mathbb Z_p}(-x).
\]
Linearity proves the assertion for the entire test space.
\end{proof}

Define a measure on \(\mathbb Q_p^\times\) by
\begin{equation}
 d^{\times} x_p=(1-p^{-1})^{-1}\frac{dx_p}{|x|_p}.
 \label{eq:adeles-multiplicative-measure}
\end{equation}
The scaling formula \eqref{eq:adeles-additive-scaling} proves multiplicative
invariance. The shell \(\mathbb Z_p^\times=\mathbb Z_p\setminus p\mathbb Z_p\) has additive
mass \(1-p^{-1}\). It therefore has multiplicative mass one.
Every shell \(p^r\mathbb Z_p^\times\) also has multiplicative mass one.
This is a locally finite regular Haar measure on the multiplicative group.

A quasi-character of a topological group is a continuous homomorphism to
\(\mathbb C^\times\). For every \(s\in\mathbb C\), the map
\(x\mapsto|x|_p^s\) is a quasi-character of \(\mathbb Q_p^\times\).

\begin{proposition}[The finite local integral]
\label{prop:adeles-finite-factor}
The integral
\[
 Z_p(s)=\int_{\mathbb Q_p^\times}\one_{\mathbb Z_p}(x)|x|_p^s\,d^{\times} x_p
\]
is absolutely convergent exactly when \(\operatorname{Re}s>0\). In that
half-plane,
\[
 Z_p(s)=\sum_{r\geq0}p^{-rs}=(1-p^{-s})^{-1}.
\]
With additive measure instead,
\[
 \int_{\mathbb Z_p}|x|_p^{s-1}\,dx_p
 =\frac{1-p^{-1}}{1-p^{-s}},\qquad\operatorname{Re}s>0.
\]
The value assigned to the additive integrand at zero has no effect.
\end{proposition}

\begin{proof}
The disjoint valuation decomposition
\eqref{eq:adeles-valuation-decomposition} reduces the first integral to the
geometric series. Its absolute integral is \(\sum_{r\geq0}p^{-r\operatorname{Re}s}\).
The additive shell \(p^r\mathbb Z_p^\times\) has mass
\((1-p^{-1})p^{-r}\), which gives the second formula.
The singleton \(\{0\}\) has measure zero since it lies in balls of mass
\(p^{-n}\) for every \(n\).
\end{proof}

These measure normalizations are the rational unramified case of Tate's
conventions \cite{Tate1950}*{\S2.2, pp.~309--311, and \S2.3, pp.~312--313}.
Tate's additive character is the complex conjugate of the one fixed here.
Conjugating the character leaves these self-dual measures unchanged.
The computations above establish the formulas directly, including the prime two.

\subsection{The real Gaussian and Poisson summation}
\label{sec:adeles-real-poisson}

At the real place, put
\[
 \psi_\infty(x)=e^{2\pi ix},\qquad dx_\infty=dx,
 \qquad d^{\times} x_\infty=\frac{dx}{|x|}.
\]
Here \(dx\) is Lebesgue measure, and \(d^{\times} x_\infty\) is a measure on
\(\mathbb R^\times\). Both signs occur. The real Schwartz space
\(\mathcal S(\mathbb R)\) consists of smooth functions whose derivatives
decay faster than every inverse power. Its Fourier transform is
\[
 \widehat f(y)=\int_{\mathbb R}f(x)e^{-2\pi ixy}\,dx.
\]
Differentiation under the integral and integration by parts show that this
map preserves \(\mathcal S(\mathbb R)\).

\begin{lemma}[Gaussian Fourier transform and Mellin integral]
\label{lem:adeles-gaussian}
For \(t>0\), define \(g_t(x)=e^{-\pi t^2x^2}\). Then
\[
 \widehat g_t(y)=t^{-1}e^{-\pi y^2/t^2}.
\]
For \(\operatorname{Re}s>0\),
\[
 Z_\infty(s)=\int_{\mathbb R^\times}e^{-\pi x^2}|x|^s\,d^{\times} x_\infty
 =\pi^{-s/2}\Gamma(s/2)=:\GammaR(s).
\]
The latter integral is absolutely convergent exactly in this half-plane.
\end{lemma}

\begin{proof}
The positive Gaussian integral has value one. Indeed, its square is the
integral of \(e^{-\pi(x^2+y^2)}\) over the plane. Polar coordinates give
\(2\pi\int_0^\infty e^{-\pi r^2}r\,dr=1\).
Set \(G(y)=\widehat g_1(y)\). Differentiation and integration by parts give
\(G'(y)=-2\pi yG(y)\), with \(G(0)=1\).
The boundary term vanishes by Gaussian decay. Thus
\(G(y)=e^{-\pi y^2}\). Substitution \(x\mapsto x/t\) gives the scaled
identity.

Evenness and the substitution \(u=\pi x^2\) give
\[
 2\int_0^\infty e^{-\pi x^2}x^{s-1}\,dx
 =\pi^{-s/2}\int_0^\infty e^{-u}u^{s/2-1}\,du.
\]
Gaussian decay controls infinity. At zero, absolute convergence is
equivalent to \(\operatorname{Re}s>0\).
\end{proof}

\begin{proposition}[The local factors]
\label{prop:adeles-local-factors}
\label{prop:local-tate}
With the specified additive character and multiplicative measures, for
\(\operatorname{Re}s>0\),
\[
 Z_p(\one_{\mathbb Z_p},s)=(1-p^{-s})^{-1},\qquad
 Z_\infty(e^{-\pi x^2},s)=\GammaR(s)=\pi^{-s/2}\Gamma(s/2).
\]
Here \(Z_v(f,s)=\int_{\mathbb Q_v^\times}f(x)|x|_v^s\,d^{\times} x_v\),
where \(\mathbb Q_\infty=\mathbb R\).
\end{proposition}

\begin{proof}
These are Proposition~\ref{prop:adeles-finite-factor} and
Lemma~\ref{lem:adeles-gaussian} with their test functions displayed.
\end{proof}

For completeness, the Gaussian also proves real Fourier inversion with
this measure. If \(f\in\mathcal S(\mathbb R)\) and \(\epsilon>0\), Fubini
and Lemma~\ref{lem:adeles-gaussian} give
\[
 \int_{\mathbb R}\widehat f(y)e^{2\pi ixy}e^{-\pi\epsilon y^2}\,dy
 =\int_{\mathbb R}f(v)\epsilon^{-1/2}
       e^{-\pi(x-v)^2/\epsilon}\,dv.
\]
The right side tends to \(f(x)\), since the positive kernels have mass one
and their mass outside every fixed neighborhood of zero tends to zero.
The left side tends to the unregularized integral by dominated convergence,
because \(\widehat f\) is integrable. Hence
\(\widehat{\widehat f}(x)=f(-x)\), which fixes the self-dual normalization.

\begin{proposition}[Real Poisson summation]
\label{prop:adeles-real-poisson}
For \(f\in\mathcal S(\mathbb R)\), \(d>0\), and \(r\in\mathbb R\),
\begin{equation}
 \sum_{n\in\mathbb Z}f(r+dn)
 =\frac1d\sum_{m\in\mathbb Z}
       \widehat f(m/d)e^{2\pi irm/d}.
 \label{eq:adeles-real-poisson}
\end{equation}
Both series converge absolutely. For fixed \(d\), they converge uniformly
as \(r\) ranges over a compact set.
\end{proposition}

\begin{proof}
Periodize \(f\) by \(P(x)=\sum_n f(x+dn)\). Rapid decay of all derivatives
gives uniform convergence on \([0,d]\), including each differentiated
series. Thus \(P\) is smooth and \(d\)-periodic.
Its \(m\)th Fourier coefficient is
\[
 \frac1d\int_0^dP(x)e^{-2\pi imx/d}\,dx
 =\frac1d\int_{\mathbb R}f(x)e^{-2\pi imx/d}\,dx
 =\frac1d\widehat f(m/d).
\]
Absolute integrability justifies the interchange and the partition into
intervals. The coefficients decay faster than every inverse power, so
their Fourier series converges absolutely and uniformly.

To verify equality with \(P\), one may use the elementary Fejer kernels
\[
 K_N(u)=\frac1N\left|\sum_{j=0}^{N-1}e^{2\pi iju/d}\right|^2.
\]
They are nonnegative and satisfy \(d^{-1}\int_0^dK_N(u)\,du=1\).
Outside a fixed neighborhood of \(d\mathbb Z\), the geometric-series
formula bounds them by a constant times \(N^{-1}\).
Uniform continuity therefore makes convolution with \(K_N\), using
\(du/d\), converge uniformly to any continuous periodic function.
A function with zero Fourier coefficients has all these convolutions zero.
Apply this observation to the difference between \(P\) and its uniformly
convergent Fourier series, then evaluate at \(r\).
\end{proof}

\subsection{The adele ring and the rational fundamental domain}
\label{sec:adeles-fundamental-domain}

Define the finite adele ring and the full adele ring by
\[
 \mathbb A_f=\prod_p'(\mathbb Q_p,\mathbb Z_p)
 =\{(x_p)_p:x_p\in\mathbb Z_p\text{ for all but finitely many }p\},
 \qquad \Adele=\mathbb R\times\mathbb A_f.
\]
Ring operations are componentwise. The prime on the product imposes the
displayed integrality condition. Put \(\widehat{\mathbb Z}=\prod_p\mathbb Z_p\).
The restricted product topology has basic open sets
\[
 \prod_{p\in S}U_p\times\prod_{p\notin S}\mathbb Z_p,
\]
where \(S\) is finite and \(U_p\subset\mathbb Q_p\) is open.
These sets cover \(\mathbb A_f\), and their translates form the same basis.
In particular, \(\widehat{\mathbb Z}\) is compact and open. Compactness follows from
the countable product of the compact metrizable spaces \(\mathbb Z_p\), or from
the same diagonal subsequence argument used for their digits.
Addition and multiplication are continuous: at the finitely many
exceptional primes this follows from local continuity, and outside them
\(\mathbb Z_p\) is closed under both operations. Thus \(\Adele\) is a locally
compact Hausdorff topological ring.

Every rational number belongs to \(\mathbb Z_p\) at all primes except those
dividing its denominator. Hence the diagonal embedding
\(\mathbb Q\hookrightarrow\Adele\) is defined.

Give \(\widehat{\mathbb Z}\) its product probability measure. On each open subring
\[
 A_S=\mathbb R\times\prod_{p\in S}\mathbb Q_p\times\prod_{p\notin S}\mathbb Z_p
\]
use the product of Lebesgue measure, the finitely many local additive
measures, and the remaining product probability measure.
These measures agree on overlaps because \(\operatorname{vol}(\mathbb Z_p)=1\).
An increasing exhaustion by finite sets of primes defines a measure
\(dx=dx_\infty\prod_pdx_p\) on \(\Adele\).
Translation by an adele preserves it: each relevant set and translation
lies in one sufficiently large \(A_S\), where product invariance applies.
It is regular and locally finite, as follows on each open \(A_S\) from
the corresponding product measures. This constructs the additive Haar
measure without an unspecified constant.

\begin{proposition}[The rational lattice]
\label{prop:adeles-rational-lattice}
The diagonal subgroup \(\mathbb Q\subset(\Adele,+)\) is discrete and closed.
The set
\begin{equation}
 \mathcal F=[0,1)\times\widehat{\mathbb Z}
 \label{eq:adeles-fundamental-domain}
\end{equation}
contains exactly one representative of each coset in \(\Adele/\mathbb Q\).
It has additive measure one, and its closure is compact.
For every nonnegative Borel function \(H\),
\[
 \int_{\Adele}H(x)\,dx
 =\int_{\mathcal F}\sum_{q\in\mathbb Q}H(x+q)\,dx.
\]
The same formula holds for an absolutely integrable complex function.
\end{proposition}

\begin{proof}
For \(x\in\Adele\), the sum \(r=\sum_p\{x_p\}_p\) is a finite rational
sum. At a fixed prime \(p\), each term with another prime in its denominator
belongs to \(\mathbb Z_p\). Consequently \(x_p-r\in\mathbb Z_p\) for every \(p\).
Choose the integer \(n=\lfloor x_\infty-r\rfloor\). Then
\(x-(r+n)\in\mathcal F\).

A rational number integral at every prime is an integer: in a reduced
fraction, any prime dividing the denominator would violate integrality.
Two representatives in \(\mathcal F\) that differ by a rational number
therefore differ by an integer. Their real coordinates force that integer
to be zero. This proves uniqueness.

The neighborhood \((-1,1)\times\widehat{\mathbb Z}\) meets \(\mathbb Q\) only at zero,
so the subgroup is discrete. It is also closed. Indeed, choose a smaller
neighborhood \(V\) with \(V-V\subset(-1,1)\times\widehat{\mathbb Z}\).
Each translate of \(V\) contains at most one rational point. A point in
the closure of \(\mathbb Q\) must then equal that point, by the Hausdorff
property. The product measure gives \(\operatorname{vol}(\mathcal F)=1\),
and \([0,1]\times\widehat{\mathbb Z}\) is compact. Finally, the countably many sets
\(\mathcal F+q\) form a disjoint Borel partition. Translation invariance
and Tonelli's theorem prove the integral formula.
\end{proof}

Define
\begin{equation}
 \psi(x)=\psi_\infty(x_\infty)\prod_p\psi_p(x_p)
 =\exp\!\left(2\pi i\left(x_\infty-\sum_p\{x_p\}_p\right)\right).
 \label{eq:adeles-character}
\end{equation}
Only finitely many factors differ from one. On each \(A_S\), the character
is a finite product of continuous characters. It is therefore continuous
on \(\Adele\). For \(q\in\mathbb Q\), the rational number
\(q-\sum_p\{q\}_p\) is integral at every prime, hence belongs to
\(\mathbb Z\). Thus \(\psi(q)=1\), and \(\psi\) descends to
\(\Adele/\mathbb Q\).

The annihilator of the diagonal \(\mathbb Q\) under the pairing
\((x,y)\mapsto\psi(xy)\) is exactly the diagonal \(\mathbb Q\).
Here is a direct verification. Subtract a rational number from an element
\(y\) of the annihilator to put it in \(\mathcal F\).
Testing against integers first gives \(y_\infty=0\), since
\(\psi(ny)=e^{2\pi iny_\infty}\). Testing against \(p^{-k}\), for each
prime \(p\) and \(k\geq1\), then gives \(y_p\in p^k\mathbb Z_p\).
All other finite components contribute one. Hence every \(y_p\) is zero.
Conversely, \(\psi(qr)=1\) for rational \(q,r\).
The pairing therefore uses the rational lattice with its precise dual and
covolume. These are the rational cases of Tate's additive global data
\cite{Tate1950}*{Theorems~4.1.3--4.1.4, pp.~330--331}.

\subsection{Fourier summation on the adeles}
\label{sec:adeles-poisson}

Define \(\mathcal S(\Adele)\) to be the finite linear span of functions
\[
 F(x)=f_\infty(x_\infty)\prod_p f_p(x_p),
 \qquad f_\infty\in\mathcal S(\mathbb R),\quad
 f_p\in\mathcal S(\mathbb Q_p),
\]
where \(f_p=\one_{\mathbb Z_p}\) for all but finitely many primes.
This is the Schwartz--Bruhat test space used here.
Every such function is integrable for the additive Haar measure.
Fubini and the local ball calculation give the product Fourier transform
\[
 \widehat F(y)=\int_{\Adele}F(x)\overline{\psi(xy)}\,dx.
\]
For a fixed \(y\), only finitely many local integrals differ from the
unit-ball integral. The latter vanishes when its frequency is nonintegral.
The transform preserves the displayed test space, and its square is
\(F(x)\mapsto F(-x)\), by the proved local inversions.

\begin{lemma}[Elementary finite-adele test functions]
\label{lem:adeles-elementary-tests}
Every \(F\in\mathcal S(\Adele)\) is a finite sum of functions
\[
 f(x_\infty)\one_{r+d\widehat{\mathbb Z}}(x_f),
 \qquad f\in\mathcal S(\mathbb R),\quad r\in\mathbb Q,\quad
 d\in\mathbb Q_{>0}.
\]
\end{lemma}

\begin{proof}
Decompose each of the finitely many exceptional local factors into ball
indicators, using the proof of Proposition~\ref{prop:adeles-ball-fourier}.
Each resulting finite factor is a constant times an indicator of a coset
\(a_f+\prod_p p^{r_p}\mathbb Z_p\), with only finitely many \(r_p\neq0\).
Set \(d=\prod_p p^{r_p}>0\). Multiplication by the other prime factors
acts by units on \(\mathbb Z_p\), so this subgroup is \(d\widehat{\mathbb Z}\).
The proof of Proposition~\ref{prop:adeles-rational-lattice}, applied to
\(a_f/d\), supplies \(b\in\mathbb Q\) with \(a_f/d-b\in\widehat{\mathbb Z}\).
Then \(r=db\) is the required rational representative.
\end{proof}

For a rational \(q\neq0\), prime factorization proves the product formula
\begin{equation}
 |q|_\infty\prod_p|q|_p=1.
 \label{eq:adeles-product-formula}
\end{equation}
In particular, the additive finite measure of \(d\widehat{\mathbb Z}\) is \(d^{-1}\)
for \(d\in\mathbb Q_{>0}\).

\begin{theorem}[Adelic Poisson summation over the rationals]
\label{thm:adeles-poisson}
For every \(F\in\mathcal S(\Adele)\),
\begin{equation}
 \sum_{q\in\mathbb Q}F(q)=\sum_{q\in\mathbb Q}\widehat F(q).
 \label{eq:adeles-poisson}
\end{equation}
Both series converge absolutely. The measure and character are exactly
\eqref{eq:adeles-additive-scaling} and \eqref{eq:adeles-character}.
\end{theorem}

\begin{proof}
By Lemma~\ref{lem:adeles-elementary-tests}, it suffices to take
\(F(x)=f(x_\infty)\one_{r+d\widehat{\mathbb Z}}(x_f)\).
A rational \(q\) satisfies \(q-r\in d\widehat{\mathbb Z}\) precisely when
\((q-r)/d\in\mathbb Z\). Thus
\[
 \sum_{q\in\mathbb Q}F(q)=\sum_{n\in\mathbb Z}f(r+dn).
\]
Put \(\psi_f(x_f)=\prod_p\psi_p(x_p)\).
The local Fourier formula and the product formula give
\[
 \widehat F(y)
 =d^{-1}\widehat f(y_\infty)
       \overline{\psi_f(ry_f)}\one_{d^{-1}\widehat{\mathbb Z}}(y_f).
\]
For rational \(q\), the finite support condition is \(q=m/d\), with
\(m\in\mathbb Z\). Since \(rq\) is rational, \(\psi(rq)=1\), whence
\(\overline{\psi_f(rq)}=e^{2\pi irq}\). Therefore
\[
 \sum_{q\in\mathbb Q}\widehat F(q)
 =d^{-1}\sum_{m\in\mathbb Z}\widehat f(m/d)e^{2\pi irm/d}.
\]
Proposition~\ref{prop:adeles-real-poisson} identifies the two sides and
proves absolute convergence. A finite sum completes the proof.
\end{proof}

In particular, put
\[
 \phi(x)=e^{-\pi x_\infty^2}\one_{\widehat{\mathbb Z}}(x_f).
\]
The local formulas prove \(\widehat\phi=\phi\).
For \(t>0\), let \(a_t=(t,(1)_p)\) be the adele with real coordinate
\(t\) and all finite coordinates one. Applying
Theorem~\ref{thm:adeles-poisson} to \(x\mapsto\phi(a_tx)\) gives
\begin{equation}
 \Theta(t)=t^{-1}\Theta(t^{-1}),\qquad
 \Theta(t):=\sum_{n\in\mathbb Z}e^{-\pi n^2t^2}.
 \label{eq:adeles-theta}
\end{equation}
Here the finite unit balls impose \(q\in\mathbb Z\), and the Fourier
transform is \(t^{-1}\phi(a_{t^{-1}}x)\).
This proves the needed summation identity with its scaling constant.

\subsection{Ideles and absolute convergence of the Euler integral}
\label{sec:adeles-ideles}

The idele group is
\[
 \Adele^\times=\mathbb R^\times\times\prod_p'(\mathbb Q_p^\times,\mathbb Z_p^\times).
\]
It is the group of units of the adele ring as a set: both \(x\) and
\(x^{-1}\) are adeles precisely when the finite coordinates are units
at almost every prime. Its topology is the multiplicative restricted
product topology. In particular,
\(U=\prod_p\mathbb Z_p^\times\) is a compact open subgroup of the finite ideles.
This topology matters. It is not being inferred from the subspace topology
of the adele ring.

On each open product with finitely many exceptional primes, take the
product of the local multiplicative Haar measures. The remaining unit
groups carry their probability measures. These products agree on overlaps
because each \(\mathbb Z_p^\times\) has mass one. They define
\[
 d^{\times} x=\frac{dx_\infty}{|x_\infty|}\prod_pd^{\times} x_p.
\]
The same exhaustion argument used for additive measure proves local
finiteness, regularity, and multiplicative invariance.
For an idele define
\[
 |x|_{\Adele}=|x_\infty|\prod_p|x_p|_p\in\mathbb R_{>0}.
\]
This finite product is a continuous homomorphism. It is trivial on
\(\mathbb Q^\times\) by \eqref{eq:adeles-product-formula}.
Consequently \(x\mapsto|x|_{\Adele}^s\) is a quasi-character of the idele
group, trivial on \(\mathbb Q^\times\), for every \(s\in\mathbb C\).

Consider the scalar integral
\begin{equation}
 Z(\phi,s)=\int_{\Adele^\times}\phi(x)|x|_{\Adele}^s\,d^{\times} x.
 \label{eq:adeles-global-integral}
\end{equation}
It is an integral over the idele group, with \(\phi\) restricted from
the adele ring.

\begin{proposition}[Absolute Euler product]
\label{prop:adeles-euler-product}
The integral \eqref{eq:adeles-global-integral} is absolutely convergent
exactly when \(\operatorname{Re}s>1\). In this half-plane,
\begin{equation}
 Z(\phi,s)=\GammaR(s)\prod_p(1-p^{-s})^{-1}
 =\GammaR(s)\sum_{n\geq1}n^{-s}.
 \label{eq:adeles-euler}
\end{equation}
The product and series converge locally uniformly there.
\end{proposition}

\begin{proof}
Let \(\sigma=\operatorname{Re}s\). On the support of \(\phi\) in the
ideles, the finite valuations are nonnegative and only finitely many are
nonzero. For a finite set \(S\) of primes, restrict the valuations outside
\(S\) to zero. For \(\sigma>0\), Tonelli and the finite local calculation
give the absolute integral on this set as
\[
 \GammaR(\sigma)\prod_{p\in S}(1-p^{-\sigma})^{-1}
 =\GammaR(\sigma)
   \sum_{\substack{n\geq1\\p\mid n\Rightarrow p\in S}}n^{-\sigma}.
\]
The last identity follows by multiplying finitely many absolutely
convergent geometric series and using unique factorization.
An increasing exhaustion of the primes exhausts the support.
Every positive integer eventually occurs. Monotone convergence therefore
identifies the full absolute integral with
\(\GammaR(\sigma)\sum_{n\geq1}n^{-\sigma}\), allowing the value infinity.
The series converges exactly for \(\sigma>1\), by comparison with
\(\int_1^\infty u^{-\sigma}\,du\).
For \(\sigma\leq0\), already the real integral on the finite-unit subset
diverges at zero.

In the convergent half-plane, dominated convergence justifies the same
exhaustion for the complex integrand. Finite factorization then proves
\eqref{eq:adeles-euler}. On a compact subset with
\(\operatorname{Re}s\geq1+\delta\), the Dirichlet series has majorant
\(\sum_n n^{-1-\delta}\). The finite Euler products converge uniformly
by the same majorant and exhaustion. Equivalently, their logarithms have
the absolutely convergent expansion
\[
 \sum_p\sum_{k\geq1}\frac{p^{-ks}}{k},\qquad
 \sum_p\sum_{k\geq1}\frac{|p^{-ks}|}{k}
 \leq\frac{1}{1-2^{-1-\delta}}\sum_{n\geq2}n^{-1-\delta}.
\]
This also verifies the usual absolute convergence of the Euler product.
\end{proof}

\begin{proposition}[An explicit idele class domain]
\label{prop:adeles-idele-domain}
Every idele has a unique expression
\begin{equation}
 x=q\,(t,u),\qquad q\in\mathbb Q^\times,\quad t\in\mathbb R_{>0},
 \quad u\in U.
 \label{eq:adeles-idele-domain}
\end{equation}
Here \((t,u)\) denotes an idele, and \(t=|x|_{\Adele}\).
In these coordinates the measure is counting measure on
\(\mathbb Q^\times\), multiplied by \(dt/t\) and probability measure
\(du\) on \(U\). Thus
\[
 \mathbb Q^\times\backslash\Adele^\times
 \simeq\mathbb R_{>0}\times U
\]
as topological groups, with quotient measure \((dt/t)du\).
\end{proposition}

\begin{proof}
Let \(r_p=v_p(x_p)\), with only finitely many nonzero values, and set
\[
 q=\operatorname{sgn}(x_\infty)\prod_p p^{r_p}.
\]
Then \(u_p=x_p/q\) is a unit at every prime, while
\(t=x_\infty/q>0\). The product formula gives
\(t=|x|_{\Adele}\). Conversely, the finite valuations and real sign recover
\(q\) from any expression of the stated form, so it is unique.
The subset \(D=\mathbb R_{>0}\times U\) is an open subgroup of the
idele group, and the sets \(qD\) are its disjoint open cosets.
Restriction of \(d^{\times} x\) to \(D\) is exactly \((dt/t)du\).
Multiplicative invariance proves the measure statement on each coset.
The quotient map restricted to \(D\) is a continuous open group
isomorphism. This proves the topological assertion as well.
\end{proof}

In particular, the subgroup of idele classes of norm one has volume one.
There is no extra factor from the two real signs: the rational factor
\(q\) in \eqref{eq:adeles-idele-domain} already records that sign.

\subsection{Continuation from the Gaussian theta identity}
\label{sec:adeles-continuation}

For \(\operatorname{Re}s>1\), the preceding domain and absolute convergence
allow unfolding over \(\mathbb Q^\times\):
\begin{align}
 Z(\phi,s)
 &=\int_U\int_0^\infty
      \sum_{q\in\mathbb Q^\times}\phi(q(t,u))t^s\,\frac{dt}{t}\,du
 \notag\\
 &=\int_0^\infty(\Theta(t)-1)t^s\,\frac{dt}{t}.
 \label{eq:adeles-mellin-theta}
\end{align}
Indeed, since every \(u_p\) is a unit, the finite factors require
\(q\in\mathbb Z\setminus\{0\}\). Both positive and negative integers
occur. The sum is independent of \(u\), whose total measure is one.
This proves the Mellin identity directly from the idele integral.

For \(t\geq1\), the elementary inequality \(n^2\geq n\) gives
\begin{equation}
 0<\Theta(t)-1
 =2\sum_{n\geq1}e^{-\pi n^2t^2}
 \leq\frac{2e^{-\pi t^2}}{1-e^{-\pi t^2}}
 \leq\frac{2}{1-e^{-\pi}}e^{-\pi t^2}.
 \label{eq:adeles-theta-bound}
\end{equation}
It follows that
\[
 E(s)=\int_1^\infty(\Theta(t)-1)t^s\,\frac{dt}{t}
\]
is entire. More explicitly, on every compact set of \(s\)-values and for
every integer \(j\geq0\), the differentiated integrand is bounded by
a constant times \(e^{-\pi t^2}t^{M-1}(\log t)^j\), for some fixed
\(M\). This bound is integrable on \([1,\infty)\). Differentiation under
the integral therefore proves holomorphy and all derivative formulas.

\begin{theorem}[The completed rational zeta integral]
\label{thm:adeles-tate}
\label{thm:tate}
With the character, measures, and test function defined above, the
absolutely convergent integral on \(\operatorname{Re}s>1\) equals
\[
 Z(\phi,s)=\GammaR(s)\zeta(s)=:\Lambda(s),\qquad
 \zeta(s)=\sum_{n\geq1}n^{-s}.
\]
Its unique meromorphic continuation to \(\mathbb C\) is
\begin{equation}
 \Lambda(s)=\frac1{s-1}-\frac1s+E(s)+E(1-s).
 \label{eq:adeles-continuation}
\end{equation}
It has precisely two poles, both simple, with
\[
 \operatorname*{Res}_{s=1}\Lambda(s)=1,\qquad
 \operatorname*{Res}_{s=0}\Lambda(s)=-1,
 \qquad \Lambda(s)=\Lambda(1-s).
\]
Consequently,
\[
 \xi_{\mathbb R}(s)=\frac12s(s-1)\Lambda(s)
\]
is entire and satisfies \(\xi_{\mathbb R}(s)=\xi_{\mathbb R}(1-s)\).
Its values at both \(0\) and \(1\) are \(1/2\).
\end{theorem}

\begin{proof}
Proposition~\ref{prop:adeles-euler-product} proves the identity in the
initial half-plane. Split \eqref{eq:adeles-mellin-theta} at \(t=1\).
For \(0<t<1\), the proved theta identity gives
\[
 \Theta(t)-1=t^{-1}(\Theta(t^{-1})-1)+t^{-1}-1.
\]
When \(\operatorname{Re}s>1\), each of the following integrals converges
absolutely. Substitution \(v=t^{-1}\) gives
\begin{align*}
 \int_0^1(\Theta(t)-1)t^s\,\frac{dt}{t}
 &=\int_1^\infty(\Theta(v)-1)v^{1-s}\,\frac{dv}{v}
   +\int_0^1(t^{s-1}-t^s)\,\frac{dt}{t}\\
 &=E(1-s)+\frac1{s-1}-\frac1s.
\end{align*}
This proves \eqref{eq:adeles-continuation} on the original domain.
Since \(E\) is entire, its right side defines a meromorphic function on
\(\mathbb C\). The identity theorem gives uniqueness of the continuation.
The rational part has residues \(1\) and \(-1\), and the remaining terms
are entire. Thus neither pole cancels, and no other pole occurs.

Replacing \(s\) by \(1-s\) exchanges the two entire terms and preserves
the rational part. This proves the functional equation. Finally,
\[
 \xi_{\mathbb R}(s)
 =\frac12+\frac12s(s-1)\bigl(E(s)+E(1-s)\bigr)
\]
proves the asserted entireness, symmetry, and endpoint values.
\end{proof}

The transformation in \eqref{eq:adeles-continuation} is the theta-integral
argument in Riemann's paper \cite{Riemann1859}*{p.~3 of the transcription}.
Tate's global theorem extends this mechanism to number fields and Hecke
characters \cite{Tate1950}*{Main Theorem~4.4.1, pp.~339--341}.
Here every summation and domain formula
needed for the rational Gaussian integral has been proved explicitly.
The continuation is a meromorphic scalar function. An operator family or
cohomological determinant requires additional constructions on its own
carrier.

\section{Equilibrium states and prime occupation spaces}
\label{sec:equilibrium}

The local integral gives a geometric series because the valuation takes
nonnegative integer values on the unit ball. A Hilbert space with one
basis vector for each such value gives the same series as a heat trace.
This second construction also has an observable algebra and a time
evolution. We first determine its equilibrium state, then add the phase
observables that distinguish the full arithmetic system.

All Hilbert spaces in this section are complex, with inner products linear
in the first variable. Their operator algebras use composition as product,
the Hilbert adjoint as involution, and the operator-norm topology.
They have cohomological degree zero; their frequency decomposition under
time evolution is an additional grading, not a cochain differential.

\subsection{The equilibrium condition}

Let \(A\) be a unital complex \(C^*\)-algebra and
\(\tau:\mathbb R\to\operatorname{Aut}(A)\) a strongly continuous group.
A state is a positive complex-linear functional \(\varphi:A\to\mathbb C\)
with \(\varphi(1)=1\). An eigenvector of frequency \(\omega\in\mathbb R\)
satisfies \(\tau_t(a)=e^{it\omega}a\). An element is entire analytic if
\(t\mapsto\tau_t(a)\) extends to an entire norm-holomorphic function.

\begin{definition}[KMS condition]\label{bc:def:kms}
For \(\beta>0\), a state \(\varphi\) is KMS\(_\beta\) if, for all
\(a,b\in A\), there is a bounded continuous function on the strip
\(0\leq\operatorname{Im}z\leq\beta\), holomorphic in its interior, such that
\[
 F_{a,b}(t)=\varphi(a\tau_t(b)),\qquad
 F_{a,b}(t+i\beta)=\varphi(\tau_t(b)a).
\]
For entire analytic elements the corresponding algebraic identity is
\begin{equation}\label{bc:eq:kms}
 \varphi(ab)=\varphi(b\tau_{i\beta}(a)).
\end{equation}
\end{definition}

\begin{lemma}[The analytic criterion]\label{bc:lem:analytic}
Suppose \(\mathcal B_\tau\) is a norm-dense star subalgebra of \(A\)
spanned by eigenvectors of \(\tau\). A state is KMS\(_\beta\) exactly when
\eqref{bc:eq:kms} holds for \(a,b\in\mathcal B_\tau\).
\end{lemma}

\begin{proof}
For \(a,b\in\mathcal B_\tau\), the function
\(\varphi(a\tau_z(b))\) is a finite sum of exponentials and is bounded
on the closed strip. The difference between this function and a KMS
strip function has zero lower boundary. Schwarz reflection across that
boundary, followed by the identity theorem, makes the difference zero.
Its upper boundary then gives \eqref{bc:eq:kms} after interchanging
\(a\) and \(b\).

Conversely, the identity gives both required boundaries for the finite
exponential sum. Their absolute values are at most \(\|a\|\|b\|\).
The bounded-strip maximum principle gives the same bound in the strip.
Approximate arbitrary \(a,b\) in norm by elements of \(\mathcal B_\tau\).
The same estimate applied to differences makes the strip functions
uniformly Cauchy. Their limit is continuous on the closed strip,
holomorphic in its interior, and has the required boundaries.
\end{proof}

\subsection{A prime oscillator}

Fix a prime \(p\).  Let \(\mathscr H_p=\ell^2(\mathbb N_0)\) with
orthonormal basis \((e_{p,r})_{r\geq0}\), and set
\[
 S_pe_{p,r}=e_{p,r+1},\qquad \mathcal T_p=C^*(S_p)
 \subset\mathcal B(\mathscr H_p).
\]
Thus \(S_p^*S_p=1\) but \(S_pS_p^*\ne1\).  Define the positive operator
\[
 N_pe_{p,r}=r\log p\,e_{p,r}
\]
on its maximal domain
\[
 \mathcal D(N_p)=
 \left\{\sum_{r\geq0}a_re_{p,r}:
 \sum_{r\geq0}r^2(\log p)^2|a_r|^2<\infty\right\}.
\]
Testing its adjoint against each basis vector gives exactly the same
maximal domain, so the diagonal operator is self-adjoint. Truncation of
the coordinate sequence converges in its graph norm; hence finite-support
vectors form a core.

\begin{proposition}[Prime heat trace and Toeplitz KMS state]
\label{prop:prime-kms}
The unitaries \(U_{p,t}=e^{itN_p}\) define a strongly continuous dynamics
\[
 \alpha_{p,t}(a)=U_{p,t}aU_{p,t}^*,\qquad
 \alpha_{p,t}(S_p)=p^{it}S_p
\]
on \(\mathcal T_p\).  For \(s\in\mathbb C\),
\(e^{-sN_p}\) is trace class exactly when \(\operatorname{Re}s>0\), and
then
\[
 \operatorname{Tr}_{\mathscr H_p}(e^{-sN_p})=(1-p^{-s})^{-1}.
\]
For every real \(\beta>0\), \((\mathcal T_p,\alpha_p)\) has the unique
KMS\(_\beta\) state
\[
 \varphi_{\beta,p}(a)
 =(1-p^{-\beta})
 \operatorname{Tr}_{\mathscr H_p}(e^{-\beta N_p}a).
\]
On the dense analytic star algebra
\(\mathcal P_p=\operatorname{span}\{S_p^mS_p^{*n}:m,n\geq0\}\),
\[
 \varphi_{\beta,p}(S_p^mS_p^{*n})
 =\delta_{m,n}p^{-\beta m}.
\]
\end{proposition}

\begin{proof}
The formula for the dynamics follows on the basis vectors and extends by
norm density from \(\mathcal P_p\).  The singular values of
\(e^{-sN_p}\) are \(p^{-r\operatorname{Re}s}\), so their sum is finite
exactly when \(\operatorname{Re}s>0\); summing the diagonal entries gives
the trace formula.

For \(q=p^{-\beta}\), the operator
\((1-q)e^{-\beta N_p}\) is positive and has trace one.  Every monomial in
\(\mathcal P_p\) is entire analytic, and
\[
 e^{-\beta N_p}a=\alpha_{p,i\beta}(a)e^{-\beta N_p}.
\]
Cyclicity of the trace gives
\(\varphi_{\beta,p}(ab)=
\varphi_{\beta,p}(b\alpha_{p,i\beta}(a))\), the KMS identity in this
convention.  Direct evaluation gives the monomial formula.

Conversely, applying \eqref{bc:eq:kms} with \(b=1\) makes every
nonzero-frequency monomial have expectation zero.  For the remaining monomials,
the KMS identity gives
\[
 \varphi(S_p^mS_p^{*m})
 =p^{-\beta m}\varphi(S_p^{*m}S_p^m)=p^{-\beta m}.
\]
Lemma~\ref{bc:lem:analytic} proves the strip condition, and density of
\(\mathcal P_p\) proves uniqueness.  These formulas agree with
the prime system of Bost and Connes
\cite[\S2, Propositions~7--8]{BostConnes1995}.
\end{proof}

\begin{remark}[Heat flow is not observable dynamics]
For \(t>0\), the assignment \(S_p\mapsto p^{-t}S_p\) cannot be a unital
star-homomorphism because
\[
 (p^{-t}S_p)^*(p^{-t}S_p)=p^{-2t}1\ne1.
\]
The observable evolution is the automorphism group
\(S_p\mapsto p^{it}S_p\).  The unnormalized partition trace is not a
state, since its value at the identity is \((1-p^{-\beta})^{-1}\).
\end{remark}

\subsection{The tensor product over the primes}

Write \(\mathbb N^\times=\{1,2,\ldots\}\). For a finite set \(F\) of
primes, put
\[
 \mathcal T_F=\bigotimes_{p\in F}^{\min}\mathcal T_p,
 \qquad
 \mathcal T_{\mathrm{fin}}=\varinjlim_F\mathcal T_F,
\]
where the transition maps are \(a\mapsto a\otimes1\), and use the product
dynamics.  Relative to the vacuum vectors \(e_{p,0}\), unique
factorization gives
\[
 \bigotimes_p^{e_{p,0}}\mathscr H_p\simeq\ell^2(\mathbb N^\times).
\]
The compatible tensor-product representations define
\[
 \pi_{\mathrm{fin}}:\mathcal T_{\mathrm{fin}}
 \longrightarrow\mathcal B(\ell^2(\mathbb N^\times)),
 \qquad \pi_{\mathrm{fin}}(S_p)e_n=e_{pn}.
\]
On the latter space define
\[
 He_n=(\log n)e_n,\qquad
 \mathcal D(H)=
 \left\{\sum_{n\geq1}a_ne_n:
 \sum_{n\geq1}(\log n)^2|a_n|^2<\infty\right\}.
\]

\begin{proposition}[The noninteracting prime subsystem]
\label{prop:prime-product}
For every \(\beta>0\), \((\mathcal T_{\mathrm{fin}},\alpha)\) has one
KMS\(_\beta\) state \(\Phi_\beta\), characterized by
\[
 \Phi_\beta|_{\mathcal T_F}
 =\bigotimes_{p\in F}\varphi_{\beta,p}.
\]
The operator \(H\) is positive and self-adjoint on the displayed domain,
with finite-support vectors as a core, and implements the dynamics in the precise sense
\[
 \pi_{\mathrm{fin}}(\alpha_t(a))
 =e^{itH}\pi_{\mathrm{fin}}(a)e^{-itH}
 \qquad(a\in\mathcal T_{\mathrm{fin}}).
\]
It satisfies
\[
 e^{-sH}\text{ is trace class}
 \quad\Longleftrightarrow\quad \operatorname{Re}s>1.
\]
In that half-plane,
\begin{equation}
 \operatorname{Tr}(e^{-sH})=\sum_{n\geq1}n^{-s}=\zeta(s).
\label{bc:eq:trace}
\end{equation}
For \(\beta>1\), but not for \(0<\beta\leq1\), the product state is the
normal Gibbs state
\[
 \Phi_\beta(a)=\zeta(\beta)^{-1}
 \operatorname{Tr}(\pi_{\mathrm{fin}}(a)e^{-\beta H}).
\]
\end{proposition}

\begin{proof}
On a finite tensor product, \eqref{bc:eq:kms} with \(b=1\) kills every
monomial except those with zero total frequency.  Unique factorization of integers makes zero
total frequency equivalent to zero frequency at every prime.  Repeated use
of the KMS identity then forces the product of the local values in
Proposition~\ref{prop:prime-kms}. Conversely, the tensor product of the
local states satisfies the KMS identity on products of local analytic
monomials, by applying the local identity in each factor.
Lemma~\ref{bc:lem:analytic} proves existence at each finite stage.
The states are compatible under \(a\mapsto a\otimes1\); they therefore
define a state on the norm-completed inductive limit, where the same
analytic criterion applies. This proves existence and uniqueness there.

Under the displayed unitary, the sum of the local energies is \(\log n\).
Testing the adjoint on the basis identifies its domain with
\(\mathcal D(H)\), and truncating a sequence converges in the graph norm.
The singular values of \(e^{-sH}\) are \(n^{-\operatorname{Re}s}\).  Their
sum is finite exactly for \(\operatorname{Re}s>1\).  Normalizing by this
sum gives the Gibbs formula in that half-plane.  Below it, the abstract
product KMS state still exists, but no trace-class density
\(e^{-\beta H}\) exists in this representation.
\end{proof}

\subsection{The universal phase algebra}

The prime shifts remember divisibility but do not distinguish the phases
of roots of unity. To retain these phases, adjoin the continuous characters
of the profinite integers. The resulting algebra has the same occupation
Hamiltonian, but a larger space of equilibrium states.

Let \(\bcN=\{1,2,\ldots\}\), and let \(\bcZhat=\varprojlim_N\mathbb Z/N\mathbb Z\)
with its compact inverse-limit topology. The Chinese remainder theorem
identifies this ring with the product \(\prod_p\mathbb Z_p\) used in the
adelic construction: compatible residues modulo every positive integer
are equivalent to compatible residues modulo every prime power.
Its unit group is
\(\bcUnits=\varprojlim_N(\mathbb Z/N\mathbb Z)^\times\).
For \(r=a/N\in\mathbb Q/\mathbb Z\) and \(x\in\bcZhat\), set
\[
 \bcphase_r(x)=\exp(2\pi i a x_N/N),
\]
where \(x_N\) is the residue of \(x\) modulo \(N\).
This definition is independent of the chosen presentation of \(r\).
The functions \(\bcphase_r\) separate points. Their complex linear span is
therefore uniformly dense in \(C(\bcZhat)\).

The universal unital complex \(C^*\)-algebra \(\bcA\) has isometries
\(\mu_n\), \(n\in\bcN\), and unitaries \(e(r)\),
\(r\in\mathbb Q/\mathbb Z\), with relations
\begin{equation}\label{bc:eq:relations}
\begin{gathered}
 \mu_n^*\mu_n=1,\qquad \mu_n\mu_m=\mu_{nm},\qquad
 \mu_n\mu_m^*=\mu_m^*\mu_n\quad(\gcd(n,m)=1),\\
 e(0)=1,\qquad e(r)^*=e(-r),\qquad e(r)e(s)=e(r+s),\\
 e(r)\mu_n=\mu_ne(nr),\qquad
 \mu_ne(r)\mu_n^*=\frac1n\sum_{n\delta=r}e(\delta).
\end{gathered}
\end{equation}
The sum in the last relation has exactly \(n\) terms. Concretely, take
the free unital star algebra on these symbols, impose the displayed
relations, and use the supremum of the operator norms over all star
representations by isometries and unitaries. Each polynomial has finite
supremum norm, bounded by the sum of the absolute values of its
coefficients. Quotient by the zero seminorm ideal and complete.
The representations below prove that this completion is nonzero.
The dynamics is
\begin{equation}\label{bc:eq:dynamics}
 \sigma_t(\mu_n)=n^{it}\mu_n,
 \qquad \sigma_t(e(r))=e(r),\qquad t\in\mathbb R.
\end{equation}
These are the presentation and time convention of
Bost--Connes \cite{BostConnes1995}*{Proposition~18 and \S5}.
The relations also imply \(\mu_1=1\): multiply
\(\mu_1^2=\mu_1\) on the left by \(\mu_1^*\).

\begin{proposition}[The phase representation]\label{bc:prop:representation}
For every \(u\in\bcUnits\), the formulas
\begin{equation}\label{bc:eq:representation}
 \pi_u(\mu_m)\varepsilon_k=\varepsilon_{mk},\qquad
 \pi_u(e(r))\varepsilon_k=\bcphase_r(ku)\varepsilon_k
\end{equation}
define a unital star representation of \(\bcA\) on
\(\mathscr H=\ell^2(\bcN)\). Its inner product is linear in the first variable.
The phase subalgebra
\(D=C^*(e(r):r\in\mathbb Q/\mathbb Z)\) is canonically \(C(\bcZhat)\),
with \(e(r)=\bcphase_r\).

For the operator \(H\) and its closed diagonal domain in
Proposition~\ref{prop:prime-product},
\[
 \pi_u(\sigma_t(a))=e^{itH}\pi_u(a)e^{-itH}.
\]
Its occupation trace is \eqref{bc:eq:trace}, with the exact domain
\(\operatorname{Re}s>1\).
\end{proposition}

\begin{proof}
The adjoint of the shift by \(m\) sends \(\varepsilon_k\) to
\(\varepsilon_{k/m}\) if \(m\mid k\), and to zero otherwise.
This proves the isometry and semigroup relations. If \(\gcd(m,n)=1\),
then \(m\mid nk\) if and only if \(m\mid k\). This proves the mixed
commutation relation. The phase multiplication and its covariance with
the shifts follow by evaluation on \(\varepsilon_k\).

For the last relation in \eqref{bc:eq:relations}, choose \(\delta_0\) with
\(n\delta_0=r\). Character orthogonality gives
\[
 \frac1n\sum_{n\delta=r}\bcphase_\delta(ku)
 =\bcphase_{\delta_0}(ku)\frac1n
      \sum_{j=0}^{n-1}\exp(2\pi i jku/n)
 =\begin{cases}
   \bcphase_r((k/n)u),&n\mid k,\\
   0,&n\nmid k.
  \end{cases}
\]
Here the finite sum uses the residue of \(u\) modulo \(n\).
Its invertibility makes its vanishing condition precisely \(n\mid k\).
This verifies all relations and also proves that the universal algebra
is nonzero. Its universal norm exists because every generator has norm
one in every representation.

The dual of the discrete group \(\mathbb Q/\mathbb Z\) is \(\bcZhat\).
Indeed, a character is specified by compatible values on \(1/N\),
which are precisely the residues \(x_N\). Thus \(D\) is a quotient
of \(C(\bcZhat)\). In \(\pi_1\), an element \(f\in C(\bcZhat)\) has
diagonal values \(f(k)\). The positive integers are dense in \(\bcZhat\),
so this representation is faithful on \(C(\bcZhat)\). The quotient has
zero kernel.

The implementation of \(\sigma\) holds on generators by evaluation on
basis vectors and then on the algebra by norm continuity. The diagonal
domain, core, and trace statement are those already proved in
Proposition~\ref{prop:prime-product}.
\end{proof}

At a finite set \(F\) of primes, the vectors indexed by
\(n=\prod_{p\in F}p^{k_p}\) have energies \(\sum_{p\in F}k_p\log p\).
Hence
\begin{equation}\label{bc:eq:finite-trace}
 \sum_{n\in\bcN_F}n^{-\beta}
   =\prod_{p\in F}(1-p^{-\beta})^{-1},\qquad \beta>0,
\end{equation}
where \(\bcN_F\) consists of positive integers with prime factors in \(F\).
For \(F=\{2,3\}\) and \(\beta=1\), this trace is \(3\).
The full trace diverges at the same inverse temperature. Classifying
equilibrium states therefore requires a construction beyond that trace.

\subsection{From the universal algebra to scaling measures}

For \(f\in C(\bcZhat)\) and \(n\in\bcN\), define
\begin{equation}\label{bc:eq:alpha}
 L_nf(x)=f(nx),\qquad
 \alpha_n(f)(x)=
 \begin{cases}f(x/n),&x\in n\bcZhat,\\0,&x\notin n\bcZhat.
 \end{cases}
\end{equation}
Multiplication by \(n\) is a homeomorphism from \(\bcZhat\) onto the
clopen subset \(n\bcZhat\). In particular, \(\alpha_n(f)\) is continuous.
The relations imply
\begin{equation}\label{bc:eq:covariance}
 f\mu_n=\mu_nL_nf,\qquad
 \mu_n^*f=(L_nf)\mu_n^*,\qquad
 \mu_nf\mu_n^*=\alpha_n(f).
\end{equation}
The last identity follows first for characters from the finite sum in
\eqref{bc:eq:relations}, and then for every continuous function by density.

\begin{lemma}[Normal forms and expectation]\label{bc:lem:expectation}
The linear span of \(\mu_mf\mu_n^*\), with \(m,n\in\bcN\) and
\(f\in D\), is a dense star subalgebra of \(\bcA\). It admits a
positive unital conditional expectation \(E:\bcA\longrightarrow D\)
with
\begin{equation}\label{bc:eq:expectation}
 E(\mu_mf\mu_n^*)=\begin{cases}\alpha_n(f),&m=n,\\0,&m\ne n.
 \end{cases}
\end{equation}
Every monomial can be written with coprime \(m,n\), after replacing
its coefficient in \(D\).
\end{lemma}

\begin{proof}
If \(d=\gcd(n,k)\), the isometry and coprime relations give
\(\mu_n^*\mu_k=\mu_{k/d}\mu_{n/d}^*\). Thus
\[
 (\mu_mf\mu_n^*)(\mu_kg\mu_l^*)
 =\mu_{mk/d}\big((L_{k/d}f)(L_{n/d}g)\big)\mu_{ln/d}^*.
\]
This proves closure under products. Adjoints and all generators also
belong to the span. If \(d=\gcd(m,n)\), then
\[
 \mu_mf\mu_n^*=\mu_{m/d}\alpha_d(f)\mu_{n/d}^*,
\]
which proves reduction to coprime indices.

The compact group \(K=\prod_{p\ \mathrm{prime}}\mathbb T\) acts by
\[
 \kappa_z(\mu_n)=z(n)\mu_n,\qquad
 z(n)=\prod_pz_p^{v_p(n)},\qquad \kappa_z(f)=f.
\]
All relations are preserved. The inverse is \(\kappa_{z^{-1}}\).
Norm continuity follows first on monomials and then by density.
Integrating this action against probability Haar measure on \(K\)
defines a positive unital map \(E\). Orthogonality of its characters
gives \eqref{bc:eq:expectation}. Its range is exactly \(D\), and the
identity \(E(fag)=fE(a)g\), for \(f,g\in D\), gives the conditional
expectation property.
\end{proof}

\begin{proposition}[Faithfulness and the prime subalgebra]
\label{bc:prop:faithful}
Each representation \(\pi_u\), \(u\in\bcUnits\), is faithful.
Its restriction identifies \(C^*(\mu_p:p\text{ prime})\) with
\(\mathcal T_{\mathrm{fin}}\), sending \(\mu_p\) to the prime shift.
\end{proposition}

\begin{proof}
The expectation \(E\) is faithful. Indeed, if \(a\geq0\) and \(E(a)=0\),
then for every state \(\omega\) the nonnegative continuous function
\(z\mapsto\omega(\kappa_z(a))\) has Haar integral zero. Haar measure has
full support, so its value at the identity is zero. States separate
positive elements, and hence \(a=0\).

On \(\ell^2(\bcN)\) let \(V_z\varepsilon_k=z(k)\varepsilon_k\).
The formula on generators gives
\(\pi_u(\kappa_z(a))=V_z\pi_u(a)V_z^*\).
If \(\pi_u(a)=0\), integration implies \(\pi_u(E(a^*a))=0\).
The set \(\{ku:k\geq1\}\) is dense in \(\bcZhat\), because the
positive integers are dense and multiplication by \(u\) is a
homeomorphism. Thus \(\pi_u\) is faithful on \(D\), so
\(E(a^*a)=0\), and faithfulness of \(E\) gives \(a=0\).

The shifts \(\pi_u(\mu_p)\) do not depend on \(u\).
For a finite set \(F\), unique factorization decomposes
\(\ell^2(\bcN)\) into the tensor product of the occupation spaces at
\(p\in F\) and the space indexed by integers with no prime factor in
\(F\). The shifts act on the first factor and as the identity on the
second. This is precisely the faithful spatial representation of
\(\bigotimes_{p\in F}^{\min}\mathcal T_p\).
Taking the norm closure of the increasing finite prime subalgebras
proves the claimed identification.
\end{proof}

The same normal forms with coefficients in \(\operatorname{span}\{e(r)\}\)
give a dense entire analytic star algebra \(\mathcal B\). Its monomials
have frequency \(\log(m/n)\). Preservation of the defining relations
also proves that \eqref{bc:eq:dynamics} extends to automorphisms, with
inverse \(\sigma_{-t}\). The dynamics is strongly continuous by density.

Lemma~\ref{bc:lem:analytic} applies to this algebra with \(\tau=\sigma\).

\begin{theorem}[The measure correspondence]\label{bc:thm:measures}
Fix \(\beta>0\). Restriction to \(D=C(\bcZhat)\) gives an affine bijection
between the KMS\(_\beta\) states on \(\bcA\) and the probability
Borel measures \(\nu\) on \(\bcZhat\) satisfying
\begin{equation}\label{bc:eq:scaling}
 \nu(nB)=n^{-\beta}\nu(B)
 \quad(n\in\bcN,\ B\subset\bcZhat\text{ Borel}).
\end{equation}
The inverse sends \(\nu\) to
\(\varphi_\nu(a)=\int_{\bcZhat}E(a)\,d\nu\).
\end{theorem}

\begin{proof}
For a monomial \(a\) of frequency \(\omega\), the KMS identity with
\(b=1\) gives \(\varphi(a)=e^{-\beta\omega}\varphi(a)\).
Thus \(\varphi\) vanishes on every nonzero frequency and
\(\varphi=\varphi|_D\circ E\). Riesz representation identifies
\(\varphi|_D\) with a probability measure \(\nu\).
The KMS identity applied to \(a=\mu_n\) and \(b=f\mu_n^*\) gives
\[
 \int\alpha_n(f)\,d\nu=n^{-\beta}\int f\,d\nu.
\]
This holds for continuous \(f\) by norm density. Uniqueness of regular
Borel measures and the homeomorphism \(x\mapsto nx\) identify it
with \eqref{bc:eq:scaling}.

Conversely, \(\varphi_\nu\) is a state because \(E\) is positive and
unital. It remains to verify \eqref{bc:eq:kms} on monomials.
Products of two monomials have zero expectation unless the sum of
their frequencies is zero. In that case, reduction to
coprime indices writes the two monomials as
\(a=\mu_mf\mu_n^*\) and \(b=\mu_ng\mu_m^*\). Consequently,
\[
 \varphi_\nu(ab)=m^{-\beta}\int fg\,d\nu,
 \qquad
 \varphi_\nu(ba)=n^{-\beta}\int gf\,d\nu.
\]
Since \(\sigma_{i\beta}(a)=(m/n)^{-\beta}a\), these equations are
exactly \eqref{bc:eq:kms}. In the other cases both sides vanish.
Lemma~\ref{bc:lem:analytic} completes the proof.
\end{proof}

\subsection{Finite-prime partitions}

For a finite set \(F\) of primes, put
\[
 W_F=\bcZhat\setminus\bigcup_{p\in F}p\bcZhat,
 \qquad c_F(\beta)=\prod_{p\in F}(1-p^{-\beta}).
\]

\begin{lemma}[Valuations and reconstruction]\label{bc:lem:finite}
Let \(\nu\) satisfy \eqref{bc:eq:scaling}. Then
\begin{equation}\label{bc:eq:valuations}
 \nu\{v_p(x)=k_p\text{ for all }p\in F\}
  =\prod_{p\in F}(1-p^{-\beta})p^{-k_p\beta}
   \quad(k_p\in\mathbb N_0).
\end{equation}
In particular, \(\nu(W_F)=c_F(\beta)>0\).
With \(\theta_F=c_F(\beta)^{-1}\nu|_{W_F}\), one has
\begin{equation}\label{bc:eq:finite-reconstruction}
 \int f\,d\nu=c_F(\beta)
   \sum_{n\in\bcN_F}n^{-\beta}\int_{W_F}f(nx)\,d\theta_F(x)
   \quad(f\in C(\bcZhat)).
\end{equation}
The series is absolutely convergent.
\end{lemma}

\begin{proof}
The intersection of \(p^{k_p}\bcZhat\), for \(p\in F\), is
\((\prod_{p\in F}p^{k_p})\bcZhat\). Its measure is
\(\prod_{p\in F}p^{-k_p\beta}\) by scaling.
Finite inclusion-exclusion proves \eqref{bc:eq:valuations}.
For each prime \(p\), the set with zero \(p\)-component has measure
\(\lim_kp^{-k\beta}=0\). Outside this null set for \(p\in F\),
every \(x\) lies in exactly one set \(nW_F\), \(n\in\bcN_F\).
Scaling on each set gives \eqref{bc:eq:finite-reconstruction}.
The total weights sum to one by \eqref{bc:eq:finite-trace}.
\end{proof}

\begin{proposition}[Existence and the invariant measure]\label{bc:prop:radial}
For every \(\beta>0\), there is a unique \(\bcUnits\)-invariant
probability measure \(\nu_\beta\) satisfying \eqref{bc:eq:scaling}.
Under \(\bcZhat=\prod_p\mathbb Z_p\), it is the product of the local measures
\begin{equation}\label{bc:eq:radial}
 \nu_{\beta,p}=(1-p^{-\beta})
       \sum_{k\geq0}p^{-k\beta}(u\mapsto p^ku)_*h_p,
\end{equation}
where \(h_p\) is probability Haar measure on \(\mathbb Z_p^\times\).
\end{proposition}

\begin{proof}
Each local series has total mass one, is invariant under multiplication
by a local unit, and satisfies
\(\nu_{\beta,p}(pB)=p^{-\beta}\nu_{\beta,p}(B)\).
It gives zero mass to zero. Consistent finite products define a positive
norm-one functional on the algebra of continuous cylinder functions.
Uniform extension and Riesz representation produce the product measure.
For an integer \(n\), multiplication by \(n\) is a local unit action
times multiplication by \(p^{v_p(n)}\) at each prime. Only finitely many
of these valuations are nonzero. This proves \eqref{bc:eq:scaling} first
on cylinders and then on all Borel sets by uniqueness of measures.

For uniqueness, fix \(N=\prod_{p\mid N}p^{a_p}\). The action of
\((\mathbb Z/N\mathbb Z)^\times\) on \(\mathbb Z/N\mathbb Z\) has
orbits determined by the truncated valuations
\(\min(v_p(x),a_p)\). The Chinese remainder theorem proves this
description and the surjectivity of \(\bcUnits\) onto the finite unit group.
Invariance makes the measure uniform within each orbit.
Equation~\eqref{bc:eq:valuations} and its tail sums determine every orbit
mass. Therefore every residue class modulo \(N\) has a fixed measure.
These cylinders determine a Borel probability measure on \(\bcZhat\).
\end{proof}

For reference, if a residue \(a\bmod p^j\) has valuation \(k<j\), then
\begin{equation}\label{bc:eq:cylinder}
 \nu_{\beta,p}(a+p^j\mathbb Z_p)
 =\frac{(1-p^{-\beta})p^{-k\beta}}{(p-1)p^{j-k-1}},
 \qquad
 \nu_{\beta,p}(p^j\mathbb Z_p)=p^{-j\beta}.
\end{equation}
These are local marginal probabilities. Products of them give the
probabilities for arbitrary moduli.

\subsection{Uniqueness at and above the critical temperature}

The arithmetic input is a divergence statement for characters of a
finite unit group. The argument only uses Euler products where their
series converge absolutely.

\begin{lemma}[Character divergence]\label{bc:lem:divergence}
Let \(\chi\) be a nontrivial complex character of
\((\mathbb Z/N\mathbb Z)^\times\). Then
\begin{equation}\label{bc:eq:divergence}
 \sum_{p\nmid N}\frac{1-\operatorname{Re}\chi(p)}p=\infty.
\end{equation}
Consequently the same sum with \(p^{-\beta}\) in place of \(p^{-1}\)
diverges whenever \(0<\beta\leq1\).
\end{lemma}

\begin{proof}
Extend \(\chi\) to the integers by zero on nonunits modulo \(N\).
Its sum over a period is zero: multiplication by a unit with nontrivial
character value permutes that sum and multiplies it by that value.
Hence \(A(t)=\sum_{n\leq t}\chi(n)\) satisfies \(|A(t)|\leq N\).
For real \(s>1\), summation by parts gives
\[
 L(s,\chi)=\sum_{n\geq1}\chi(n)n^{-s}
   =s\int_1^\infty A(t)t^{-s-1}\,dt,
 \qquad |L(s,\chi)|\leq N.
\]
The principal character has series
\[
 L_0(s)=\zeta(s)\prod_{p\mid N}(1-p^{-s})\longrightarrow\infty
 \quad(s\downarrow1).
\]
Here \(\zeta(s)\geq\int_1^\infty t^{-s}\,dt=(s-1)^{-1}\),
and the finite product is bounded below by a positive constant.
Thus \(|L(s,\chi)/L_0(s)|\to0\).

Unique factorization and absolute convergence give Euler products for
both series when \(s>1\). Their logarithms satisfy
\begin{equation}\label{bc:eq:log-ratio}
 \log\left|\frac{L(s,\chi)}{L_0(s)}\right|
 =-\sum_{p\nmid N}\frac{1-\operatorname{Re}\chi(p)}{p^s}
  +\sum_{p\nmid N}\sum_{j\geq2}
     \frac{\operatorname{Re}(\chi(p)^j)-1}{jp^{js}}.
\end{equation}
The second sum is bounded in absolute value uniformly for \(s>1\),
because it is at most
\(2\sum_p\sum_{j\geq2}p^{-j}\leq2\sum_{n\geq2}[n(n-1)]^{-1}\).
If \eqref{bc:eq:divergence} failed, the first sum would also be uniformly
bounded. The ratio could then never tend to zero. This contradiction
proves \eqref{bc:eq:divergence}. Comparison of nonnegative terms proves
the assertion for smaller positive \(\beta\).
\end{proof}

\begin{lemma}[Finite-prime contraction]\label{bc:lem:contraction}
Let \(\nu\) satisfy \eqref{bc:eq:scaling}. Suppose that \(f\) factors
through \(\mathbb Z/N\mathbb Z\) and obeys
\(f(vx)=\chi(v)f(x)\) for every unit \(v\) modulo \(N\).
For every finite set \(F\) of primes not dividing \(N\),
\begin{equation}\label{bc:eq:contraction}
 \left|\int f\,d\nu\right|
 \leq\|f\|_\infty
       \prod_{p\in F}\left|\frac{1-p^{-\beta}}
                              {1-\chi(p)p^{-\beta}}\right|
 \leq\|f\|_\infty
       \exp\left(-\sum_{p\in F}(1-\operatorname{Re}\chi(p))p^{-\beta}\right).
\end{equation}
\end{lemma}

\begin{proof}
All \(n\in\bcN_F\) are units modulo \(N\), so
\(f(nx)=\chi(n)f(x)\). Substitute this equality in
\eqref{bc:eq:finite-reconstruction} and sum the absolutely convergent
geometric series. The normalized measure \(\theta_F\) has mass one,
which proves the first inequality. For \(|z|=1\) and \(0<t<1\),
\[
 \log\left|\frac{1-t}{1-zt}\right|
   =\sum_{j\geq1}\frac{\operatorname{Re}(z^j)-1}{j}t^j
   \leq-(1-\operatorname{Re}z)t.
\]
Every summand is nonpositive. Apply this inequality at
\(z=\chi(p)\), \(t=p^{-\beta}\), and exponentiate.
\end{proof}

\begin{theorem}[Uniqueness for \(0<\beta\leq1\)]\label{bc:thm:high}
For every real \(0<\beta\leq1\), the system \((\bcA,\sigma)\) has
exactly one KMS\(_\beta\) state. It is
\(\varphi_\beta=\nu_\beta\circ E\), with \(\nu_\beta\) given by
\eqref{bc:eq:radial}.
\end{theorem}

\begin{proof}
Let \(\nu\) be any scaling measure and fix a modulus \(N\).
The space of complex functions on \(\mathbb Z/N\mathbb Z\) decomposes
into isotypic subspaces for its finite unit group \(G_N\).
Explicitly, its character projections are
\[
 (P_\chi f)(x)=\frac1{|G_N|}\sum_{v\in G_N}
                         \overline{\chi(v)}f(vx),
 \qquad f=\sum_\chi P_\chi f.
\]
Character orthogonality proves the second identity, and
\((P_\chi f)(vx)=\chi(v)(P_\chi f)(x)\).
For nontrivial \(\chi\), Lemmas~\ref{bc:lem:divergence}
and~\ref{bc:lem:contraction} imply \(\nu(P_\chi f)=0\), by exhausting
the primes outside \(N\). Only the trivial character survives.
Thus \(\nu\) is invariant under multiplication by units on every
finite quotient. Locally constant functions are uniformly dense in
\(C(\bcZhat)\), so \(\nu\) is \(\bcUnits\)-invariant.
Proposition~\ref{bc:prop:radial} gives \(\nu=\nu_\beta\).
Theorem~\ref{bc:thm:measures} supplies existence and uniqueness of the state.
\end{proof}

At \(\beta=1\), formula~\eqref{bc:eq:cylinder} assigns mass \(p^{-j}\)
to every residue modulo \(p^j\). Hence \(\nu_1\) is additive Haar
measure on \(\bcZhat\). In particular,
\begin{equation}\label{bc:eq:critical-phase}
 \varphi_1(e(r))=0\quad(r\ne0),\qquad \varphi_1(1)=1.
\end{equation}
The contraction proof establishes the endpoint directly. It requires
neither nonvanishing of \(L(1,\chi)\) nor a factor classification.

\subsection{Every low-temperature state is a Gibbs mixture}

\begin{theorem}[Classification for \(\beta>1\)]\label{bc:thm:low}
Fix \(\beta>1\). For \(\lambda\in\bcProb(\bcUnits)\), define
\begin{equation}\label{bc:eq:gibbs-measure}
 \nu_{\beta,\lambda}
  =\frac1{\zeta(\beta)}\sum_{n\geq1}n^{-\beta}(u\mapsto nu)_*\lambda,
 \qquad \varphi_{\beta,\lambda}=\nu_{\beta,\lambda}\circ E.
\end{equation}
This is an affine homeomorphism from \(\bcProb(\bcUnits)\), with its weak
topology, onto all KMS\(_\beta\) states, with their weak-star topology.
Consequently these states form a simplex. Its extreme points are
\(\varphi_{\beta,u}=\varphi_{\beta,\delta_u}\), \(u\in\bcUnits\), and
\begin{equation}\label{bc:eq:gibbs-state}
 \varphi_{\beta,u}(a)
   =\zeta(\beta)^{-1}\bcTr\big(\pi_u(a)e^{-\beta H}\big).
\end{equation}
\end{theorem}

\begin{proof}
Let \(\nu\) be any scaling measure. Since
\(\bcUnits=\bigcap_p(\bcZhat\setminus p\bcZhat)\), continuity from above and
Lemma~\ref{bc:lem:finite} give
\[
 \nu(\bcUnits)=\prod_p(1-p^{-\beta})=\zeta(\beta)^{-1}.
\]
The Euler product equality follows by expanding each finite product and
then exhausting the positive integers. Absolute convergence holds for
\(\beta>1\). The sets \(n\bcUnits\), \(n\geq1\), are pairwise disjoint:
their valuation vectors determine \(n\). By scaling,
\[
 \sum_{n\geq1}\nu(n\bcUnits)
 =\nu(\bcUnits)\sum_{n\geq1}n^{-\beta}=1.
\]
Their union therefore has full measure. Define
\(\lambda=\zeta(\beta)\nu|_{\bcUnits}\), a probability measure on the
compact set \(\bcUnits\). Applying scaling to every Borel subset of
\(\bcUnits\) proves \eqref{bc:eq:gibbs-measure}. This also proves uniqueness
of the parameter measure.

Conversely, the series in \eqref{bc:eq:gibbs-measure} defines a probability
measure. The point \(nu\) lies in \(m\bcZhat\) exactly when \(m\mid n\),
because \(u\) is a unit at every prime. Thus its mass on \(mB\) is
\[
 \frac1{\zeta(\beta)}\sum_{k\geq1}(mk)^{-\beta}
          \lambda\{u:ku\in B\}
 =m^{-\beta}\nu_{\beta,\lambda}(B).
\]
Theorem~\ref{bc:thm:measures} now gives all the claimed states.

For each \(f\in C(\bcZhat)\), the series
\(\zeta(\beta)^{-1}\sum_n n^{-\beta}f(nu)\) converges uniformly in
\(u\). It is continuous, which proves continuity of the affine map
on parameter measures. A continuous bijection from a compact space to
a Hausdorff space is a homeomorphism. The extreme points of
\(\bcProb(\bcUnits)\) are precisely its point masses: a measure that is not
a point mass has a Borel set of mass strictly between zero and one,
which gives a nontrivial convex decomposition. Every measure has the
unique decomposition \(\lambda=\int\delta_u\,d\lambda(u)\).
This proves the simplex assertion without invoking a general KMS
decomposition theorem.

Finally, on \(f\in D\) the normalized trace in
\eqref{bc:eq:gibbs-state} has value
\(\zeta(\beta)^{-1}\sum_nn^{-\beta}f(nu)\).
On \(\mu_mf\mu_n^*\) with \(m\ne n\), every diagonal matrix
coefficient is zero: a nonzero coefficient sends index \(k\) to
\(mk/n\), which differs from \(k\). For \(m=n\), the element is
\(\alpha_n(f)\in D\). The trace state and \(\varphi_{\beta,u}\)
therefore agree on a dense algebra and hence everywhere.
\end{proof}

\subsection{The cyclotomic carrier and the symmetry action}

Let \(\mathbb Q^{\mathrm{cycl}}\subset\mathbb C\) be the union of all
fields \(\mathbb Q(\zeta_N)\), with \(\zeta_N=\exp(2\pi i/N)\).
The group \(G=\bcGal(\mathbb Q^{\mathrm{cycl}}/\mathbb Q)\) carries
the inverse-limit topology of its finite cyclotomic Galois groups.

\begin{lemma}[Cyclotomic identification]\label{bc:lem:galois}
The map \(u\mapsto\gamma_u\), specified by
\begin{equation}\label{bc:eq:galois}
 \gamma_u(\zeta_N)=\zeta_N^{u_N},
\end{equation}
is a topological group isomorphism \(\bcUnits\longrightarrow G\).
\end{lemma}

\begin{proof}
The cyclotomic polynomial \(\Phi_N\) is monic with integer coefficients,
as follows inductively from \(X^N-1=\prod_{d\mid N}\Phi_d(X)\).
We recall the irreducibility argument to check the full parameter space.
Let \(f\in\mathbb Z[X]\) be the monic irreducible polynomial of a
primitive root \(\zeta_N\). It divides \(\Phi_N\), by Gauss's lemma.
If a prime \(p\nmid N\) sent this root to a root of a different
irreducible factor \(g\) of \(\Phi_N\), then
\(f(X)\mid g(X^p)\). Reduction modulo \(p\) gives
\(\overline f\mid\overline g^p\), so \(\overline f\) and
\(\overline g\) have a common irreducible factor.
It would occur twice in \(\overline\Phi_N\) and thus in \(X^N-1\)
over \(\mathbb F_p\). This is impossible because the derivative
\(NX^{N-1}\) is coprime to \(X^N-1\).

The same argument applies to each root of \(f\). Consequently every
map \(\zeta_N\mapsto\zeta_N^p\), \(p\nmid N\), preserves its roots.
Every positive integer coprime to \(N\)
is a product of such primes. Thus all primitive roots are roots of
\(f\), and \(\Phi_N=f\). It follows that each unit modulo \(N\)
defines a field automorphism by \eqref{bc:eq:galois}.
Compatible residues define compatible automorphisms, hence an
automorphism of their union. Conversely, every automorphism preserves
root orders and determines these compatible units. This proves the
bijection and the group law. The finite-quotient definitions of the
two topologies prove continuity in both directions.
\end{proof}

Under this identification, \eqref{bc:eq:representation} becomes
\[
 \pi_\gamma(\mu_m)\varepsilon_n=\varepsilon_{mn},\qquad
 \pi_\gamma(e(r))\varepsilon_n
      =\gamma(e^{2\pi i nr})\varepsilon_n.
\]
Theorem~\ref{bc:thm:low} therefore has exactly the cyclotomic Galois
carrier in the Bost--Connes classification
\cite{BostConnes1995}*{Propositions~23--24 and Theorem~25}.

\begin{theorem}[Bost--Connes equilibrium classification]
\label{thm:bc}
The full system \((\bcA,\sigma)\) has a unique KMS\(_\beta\) state for
\(0<\beta\leq1\). For \(\beta>1\), its KMS simplex is affinely
homeomorphic to \(\bcProb(\bcUnits)\), and its extreme states are the
normalized Gibbs states \eqref{bc:eq:gibbs-state}.
Via \eqref{bc:eq:galois}, their parameters are exactly
\(\operatorname{Gal}(\mathbb Q^{\mathrm{cycl}}/\mathbb Q)\).
\end{theorem}

\begin{proof}
Theorems~\ref{bc:thm:measures}, \ref{bc:thm:high}, and~\ref{bc:thm:low}
construct and classify the states. Lemma~\ref{bc:lem:galois} identifies
the parameter group. In particular, no normalization of a divergent
occupation trace is used at or below \(\beta=1\).
\end{proof}

\begin{proposition}[Symmetry and the Toeplitz restriction]\label{bc:prop:symmetry}
For \(v\in\bcUnits\), the formulas
\[
 \theta_v(\mu_n)=\mu_n,\qquad \theta_v(e(r))=e(vr)
\]
define an action of \(\bcUnits\) on \(\bcA\) by automorphisms that commute
with \(\sigma\). Here \(vr\) uses the residue of \(v\) modulo the
order of \(r\). Its action on the low-temperature extreme states is
\begin{equation}\label{bc:eq:symmetry}
 \varphi_{\beta,u}\circ\theta_v=\varphi_{\beta,vu},
 \qquad \beta>1.
\end{equation}
It is free and transitive on those states. For \(0<\beta\leq1\),
the unique KMS state is fixed by every \(\theta_v\).

For every \(\beta>0\), all KMS\(_\beta\) states have the same
restriction to \(\mathcal T=C^*(\mu_n:n\in\bcN)\). On its dense
monomials this restriction is
\begin{equation}\label{bc:eq:toeplitz}
 \varphi(\mu_m\mu_n^*)=\delta_{m,n}n^{-\beta}.
\end{equation}
On any finite set of primes it is the product of the geometric
Toeplitz states with weights \((1-p^{-\beta})p^{-k\beta}\).
\end{proposition}

\begin{proof}
Multiplication by a profinite unit permutes \(\mathbb Q/\mathbb Z\)
and each fiber of multiplication by \(n\). Hence \(\theta_v\)
preserves all the defining relations. Its inverse is
\(\theta_{v^{-1}}\), and it commutes with the dynamics on generators.
The identity \(\pi_u\circ\theta_v=\pi_{vu}\) gives
\eqref{bc:eq:symmetry}. The parameterization in
Theorem~\ref{bc:thm:low} proves freeness and transitivity.
At \(0<\beta\leq1\), invariance also follows from
Theorem~\ref{bc:thm:high} or from the invariant measure.

The KMS identity gives zero at unequal indices and \(n^{-\beta}\)
at equal indices. The monomials \(\mu_m\mu_n^*\) span a dense star
algebra in \(\mathcal T\), by the product calculation in
Lemma~\ref{bc:lem:expectation}. Thus \eqref{bc:eq:toeplitz} determines the
restriction. For finitely many primes, unique factorization identifies
the diagonal powers and gives the product of their local values.
Those local values are the geometric-state expectations, since
\(\sum_{k\geq j}(1-p^{-\beta})p^{-k\beta}=p^{-j\beta}\).
\end{proof}

There is exactly one symmetry-invariant KMS state for each \(\beta>0\).
For \(\beta>1\), it is the mixture in \eqref{bc:eq:gibbs-measure} with
probability Haar measure on \(\bcUnits\). It agrees with
\(\nu_\beta\circ E\) by Proposition~\ref{bc:prop:radial}.
For a nonzero \(r=a/b\) in lowest terms, its phase value is
\begin{equation}\label{bc:eq:phase-value}
 \varphi_\beta(e(r))
   =\prod_{p^j\Vert b}
         p^{-j\beta}\frac{1-p^{\beta-1}}{1-p^{-1}}.
\end{equation}
Indeed, the average of a primitive additive character of conductor
\(p^j\) on \(p^k\mathbb Z_p^\times\) is zero for \(k<j-1\),
is \(-1/(p-1)\) for \(k=j-1\), and is one for \(k\geq j\).
The first assertion follows by subtracting the complete additive sum
on multiples of \(p\) from the complete sum modulo \(p^j\).
Inserting these three values in \eqref{bc:eq:radial} gives the local factor.
For example, \(\varphi_\beta(e(1/2))=2^{1-\beta}-1\).

\subsection{Critical and zero-temperature limits}

\begin{proposition}[Limits of the equilibrium states]\label{bc:prop:limits}
For every \(a\in\bcA\),
\begin{equation}\label{bc:eq:critical-limit}
 \sup_{\lambda\in\bcProb(\bcUnits)}
       |\varphi_{\beta,\lambda}(a)-\varphi_1(a)|\longrightarrow0
       \quad(\beta\downarrow1).
\end{equation}
At zero temperature, define
\[
 \varphi_{\infty,u}(a)
  =\langle\pi_u(a)\varepsilon_1,\varepsilon_1\rangle,
 \qquad
 \varphi_{\infty,\lambda}=\int_{\bcUnits}\varphi_{\infty,u}\,d\lambda(u).
\]
These are states, and
\begin{equation}\label{bc:eq:zero-limit}
 \|\varphi_{\beta,\lambda}-\varphi_{\infty,\lambda}\|
       \leq2\bigl(1-\zeta(\beta)^{-1}\bigr)\longrightarrow0
       \quad(\beta\to\infty),
\end{equation}
uniformly in \(\lambda\). The weak-star limits of KMS states as
\(\beta\to\infty\) are exactly these mixtures. Their extreme points
are indexed by \(\bcUnits\), or equivalently by \(G\), and
\(\varphi_{\infty,u}(e(r))=\bcphase_r(u)\).
\end{proposition}

\begin{proof}
Suppose \eqref{bc:eq:critical-limit} fails. There are \(\beta_j\downarrow1\)
and parameter measures \(\lambda_j\) whose state values stay a fixed
positive distance from \(\varphi_1(a)\). The state space is weak-star
compact, so a subnet converges to a state \(\psi\). For
\(b,c\in\mathcal B\), the KMS identity passes to the limit because
\(\sigma_{i\beta_j}(b)\to\sigma_i(b)\) in norm. Therefore \(\psi\)
is KMS\(_1\) by Lemma~\ref{bc:lem:analytic}, and
Theorem~\ref{bc:thm:high} gives \(\psi=\varphi_1\). This contradicts
the choice of the subnet.

For the other limit, the Gibbs formula writes each state as a convex
sum of vector states at \(\varepsilon_n\). The weight at
\(\varepsilon_1\) is \(\zeta(\beta)^{-1}\). The remaining weights
have total mass \(1-\zeta(\beta)^{-1}\). Two states are at norm distance
at most two, which proves \eqref{bc:eq:zero-limit} and its integrated form.
The zeta series tends to one by dominated convergence, using any fixed
exponent greater than one as a summable majorant.

The map \(u\mapsto\varphi_{\infty,u}(a)\) is continuous first on
monomials and then on all \(a\) by uniform approximation. Thus mixtures
depend continuously on \(\lambda\). Compactness of \(\bcProb(\bcUnits)\)
and \eqref{bc:eq:zero-limit} identify every weak-star limit with such a
mixture. Conversely, keeping \(\lambda\) fixed produces every mixture.
The phase algebra separates measures on \(\bcUnits\), since its restricted
characters are uniformly dense in \(C(\bcUnits)\). Hence the mixture map
is injective and its extreme points are precisely the point measures.
\end{proof}

The zero-temperature assertion concerns limits of finite-temperature
KMS states; it does not classify all ground states of the dynamics.
The normalization in \eqref{bc:eq:gibbs-state} always uses the
trace-class half-plane \(\operatorname{Re}s>1\). The state at the
critical point is instead constructed by the scaling measure, and
\eqref{bc:eq:critical-limit} specifies its relation to the Gibbs family.


\subsection{Vacua and augmentations}

The occupation vacuum is the vector \(\varepsilon_1\), of energy zero.
Its vector state \(\varphi_{\infty,u}\) is a zero-temperature limit,
whereas a KMS\(_\beta\) state with finite \(\beta>0\) obeys
\(\varphi(\mu_n\mu_n^*)=n^{-\beta}\).
An algebra augmentation means a unital multiplicative map to the
coefficient field. Here the formulas
\[
 \epsilon(e(r))=1,\qquad \epsilon(\mu_n)=1
\]
preserve every defining relation and give a star augmentation
\(\epsilon:\bcA\to\mathbb C\). It has
\(\epsilon(\mu_n\mu_n^*)=1\), so it is not KMS\(_\beta\) for
\(\beta>0\). The vacuum state is not multiplicative either: for
\(n>1\), its value on \(\mu_n\) is zero, while its value on
\(\mu_n^*\mu_n\) is one. Thus the equal scalar Euler factor in the
local integral and occupation trace specifies none of these three
additional choices. A chiral differential or arithmetic cohomological
comparison must be constructed on further data.

\section{The divisor monoid and its integral bar}
\label{sec:divisor-bar}

The shifts $e_n\mapsto e_{pn}$ compose according to multiplication of
integers.  Fix a finite ordered set $F$ of primes and let
\[
 M_F=\left\{\prod_{p\in F}p^{r_p}:r_p\in\mathbb N_0\right\}.
\]
Unique factorization identifies this monoid with the effective divisors
$\sum_{p\in F}r_p[p]$.  Over any commutative ring $R$, its monoid algebra is
\[
 A_F=R[M_F]=R[u_p:p\in F],\qquad
 \epsilon(u_p)=0,\qquad I_F=\ker\epsilon.
\]
Here $u_p$ represents the divisor $[p]$ and multiplication adds divisors.
The augmentation takes the coefficient of the empty divisor.
The unit is the empty divisor. The coefficient modules have the discrete
topology, and all tensors in this section are algebraic.  All algebra
elements have homological degree zero.

For one prime, multiplication by $u_p$ gives an exact resolution
\[
 0\longrightarrow A_{\{p\}}\eta_p
 \xrightarrow{\ d\ }A_{\{p\}}
 \xrightarrow{\epsilon}R\longrightarrow0,
 \qquad d\eta_p=u_p.
\]
It is exact over arbitrary $R$: multiplication by an indeterminate is
injective, and the quotient is $R$.  The classical Koszul resolution for these polynomial generators is their
finite tensor product:
\[
 K_F=A_F\otimes_R\Lambda_R(\eta_p:p\in F),\qquad |\eta_p|=1,
 \qquad d\eta_p=u_p,
\]
with the graded derivation differential.  The exterior algebra imposes
$\eta_p^2=0$ also in characteristic two.

The normalized bar computing $R\otimes_{A_F}^{\mathbf L}R$ is the
algebraic direct sum
\[
 B(A_F)=\bigoplus_{n\geq0}I_F^{\otimes_R n},
 \qquad |[a_1|\cdots|a_n]|=n,
\]
with
\[
 d[a_1|\cdots|a_n]
 =\sum_{j=1}^{n-1}(-1)^j
       [a_1|\cdots|a_ja_{j+1}|\cdots|a_n].
\]
The empty word has degree zero.  Deconcatenation is its coproduct.

\begin{proposition}[Integral prime-set bar comparison]
\label{prop:divisor-bar}
Give $\Lambda_R(\eta_p:p\in F)$ zero differential and primitive generators
for its graded coalgebra structure.  The map
\[
 \eta_{p_1}\wedge\cdots\wedge\eta_{p_n}
 \longmapsto
 \sum_{\sigma\in S_n}\operatorname{sgn}(\sigma)
       [u_{p_{\sigma(1)}}|\cdots|u_{p_{\sigma(n)}}]
 \quad(p_1<\cdots<p_n)
\]
is a quasi-isomorphism of coaugmented differential graded coalgebras
\[
 \Lambda_R(\eta_p:p\in F)\longrightarrow B(A_F).
\]
It is compatible with inclusions of finite prime sets and with scalar
extension.  In particular,
\[
 \operatorname{Tor}^{A_F}_n(R,R)=\Lambda_R^n(R^F).
\]
\end{proposition}

\begin{proof}
The tensor product defining $K_F$ resolves $R$.  Indeed, each one-variable
resolution is a complex of free $R$-modules with an $R$-linear contraction
onto its constant term.  Their finite tensor product has the same property.
Its terms are free $A_F$-modules.

Lift the displayed map to $K_F\to A_F\otimes_R B(A_F)$ by retaining the
left coefficient.  The target is the normalized left bar resolution of
$R$ over $A_F$.  Its differential has the additional first face
$a[a_1|\cdots|a_n]\mapsto aa_1[a_2|\cdots|a_n]$.
The internal faces in the alternating sum cancel in pairs, because
$u_pu_q=u_qu_p$.  The first faces are exactly the terms of the Koszul
differential.  Thus this lift is a chain map inducing the identity on $R$.
Both complexes are free resolutions: the bar resolution has the usual
contraction inserting the augmentation-ideal part of the left coefficient.
A chain map lifting the identity between projective resolutions is a
chain homotopy equivalence.  To construct its inverse and the two
homotopies, lift each cycle through the next surjective differential,
inductively in homological degree.  Projectivity permits every lift.
Tensoring over $A_F$ with $R$ preserves this homotopy equivalence and
annuls the Koszul differential.

For the coalgebra assertion, cut each alternating word after its $j$th
entry.  Grouping terms by the subset of generators before the cut gives
exactly the signed subset coproduct of the exterior coalgebra.  This
also verifies the empty cuts and the counit.  The formulas preserve
prime inclusions and coefficients.  No factorial is inverted.
\end{proof}

For two primes the degree-two representative is
$[u_p|u_q]-[u_q|u_p]$.  Its differential is
$-u_pu_q+u_qu_p=0$, and its reduced coproduct is
$[u_p]\otimes[u_q]-[u_q]\otimes[u_p]$.
For three distinct primes the exterior degree-three class is nonzero
whenever $R\ne0$.  The untwisted divisor construction therefore has no
uniform concentration in degrees zero through two.

Filtered union over finite prime sets gives the same comparison for
$A=R[u_p:p\text{ prime}]$ and the direct-sum exterior coalgebra on all
prime generators.  Every bar tensor involves only finitely many monomials
and therefore finitely many primes.  Thus both coalgebras and the comparison
are the filtered colimits of the finite-set constructions.  Filtered
colimits of $R$-modules are exact, so their comparison remains a
quasi-isomorphism.  This uses no product completion in word length.

Over $\mathbb C$, the orthonormal monomial completion of $A_F$ is
$\ell^2(M_F)$.  Multiplication by $u_p$ extends to $S_p$.
The self-adjoint diagonal operator
$N_Fu^{\mathbf r}=(\sum_pr_p\log p)u^{\mathbf r}$, with
$\mathbf r=(r_p)_{p\in F}$, has domain
\[
 \mathcal D(N_F)=\left\{(c_{\mathbf r})\in\ell^2(M_F):
 \sum_{\mathbf r}\left(\sum_{p\in F}r_p\log p\right)^2
                 |c_{\mathbf r}|^2<\infty\right\}.
\]
Its heat trace is
$\prod_{p\in F}(1-p^{-s})^{-1}$ for $\operatorname{Re}s>0$.
This constructs the algebra before its scalar trace.  The polynomial
augmentation has $\epsilon(u_p)=0$, whereas the Gibbs state satisfies
$\varphi_{\beta,p}(S_p^*S_p)=1$ and
$\varphi_{\beta,p}(S_p)=0$.  It is a positive functional, not a
multiplicative augmentation.  Adjoining the Hilbert adjoints therefore
changes the available boundary data for a bar.

\section{Arithmetic comparison maps}
\label{sec:arithmetic-comparison}

Proposition~\ref{prop:divisor-bar} uses the multiplication of effective
divisors and its augmentation module. To compare it with arithmetic cochains,
one must specify the coefficient category and the boundary modules there.  Let \((\mathcal C,\widehat\otimes,\one)\) be a stable, closed,
\(\mathbb C\)-linear symmetric monoidal category with a dg model.  Assume
that it admits geometric realizations and filtered colimits and that
tensoring preserves them separately in each variable. Here $E_1$ means
associative multiplication with its coherent higher homotopies. For an
\(E_1\)-algebra \(A\), a right \(A\)-module \(M\), and a left
\(A\)-module \(N\), the simplicial two-sided bar has
\[
 \operatorname{Bar}_A(M,N)_n
 =M\widehat\otimes A^{\widehat\otimes n}\widehat\otimes N.
\]
Its faces use the multiplication and module actions, while its degeneracies
use the unit.  Its realization computes \(M\widehat\otimes_A N\) under the
stated hypotheses
\cite[Construction~4.4.2.7 and Theorem~4.4.2.8]{LurieHA}.  An augmentation
\(\epsilon:A\to\one\) makes \(\one\) into both boundary modules and hence
defines the reduced bar \(B(A)=\operatorname{Bar}_A(\one,\one)\).

An associative chiral duality must compare this multiplication
and its bar with another associative object in a specified category of
factorization data. A Lie--commutative chiral duality has different input
and output operations and supplies no such comparison without an
additional functor and compatibility maps.

\begin{proposition}[One heat trace, inequivalent bars]
\label{prop:two-bars}
Fix \(p\) and let
\(V_p=\bigoplus_{r\geq0}\mathbb C e_r\), with
\(N_pe_r=r\log p\,e_r\).  Give $e_r$ weight $r$ and put every $e_r$ in homological degree zero.
The same pointed weight-graded vector space carries the
two augmented associative algebra structures
\[
 A_p^{\mathrm{pol}}=\mathbb C[u],\quad e_r=u^r,
 \qquad
 A_p^{\mathrm{sq}}=\mathbb C\oplus I,\quad
 I=\bigoplus_{r\geq1}\mathbb C e_r,\quad I^2=0,
\]
with unit \(e_0\) and augmentation \(e_0\mapsto1\),
\(e_r\mapsto0\) for \(r>0\).  In both cases \(N_p\) is a derivation, and
on the common Hilbert completion both algebras give the heat trace
\[
 \operatorname{Tr}(e^{-sN_p})=(1-p^{-s})^{-1},
 \qquad \operatorname{Re}s>0.
\]
Nevertheless, their reduced bars are not quasi-isomorphic.  For the
augmentation module \(\mathbb C\),
\[
 \operatorname{Tor}^{A_p^{\mathrm{pol}}}_n(\mathbb C,\mathbb C)
 =\begin{cases}
   \mathbb C,&n=0,\\
   \mathbb C\eta_p,&n=1,\\
   0,&n\geq2,
  \end{cases}
 \qquad
 \operatorname{Tor}^{A_p^{\mathrm{sq}}}_n(\mathbb C,\mathbb C)
 =I^{\otimes n}\quad(n\geq0).
\]
Thus the energy spectrum and its trace do not select an \(E_1\)-algebra or
a bar construction.
\end{proposition}

\begin{proof}
The two-term Koszul resolution of \(\mathbb C\) over \(\mathbb C[u]\)
gives the first calculation.  In the normalized bar of
\(\mathbb C\oplus I\), every multiplication face on the augmentation ideal
vanishes because \(I^2=0\), so the reduced bar differential is zero and its
degree-\(n\) term is \(I^{\otimes n}\).  The derivation identity is immediate
for the polynomial product and is still valid when two positive-weight
elements have zero product.  The heat trace depends only on the common
diagonal operator.
\end{proof}

\begin{remark}
The Toeplitz algebra does not remove this ambiguity.  It has no augmentation
sending \(S_p\) to zero, since \(S_p^*S_p=1\).  Every character sends
\(S_p\) to a scalar of modulus one, and no such character is invariant under
\(S_p\mapsto p^{it}S_p\).  A Toeplitz bar therefore requires explicitly
chosen boundary modules; it cannot silently inherit the polynomial
augmentation.
\end{remark}

\begin{openproblem}[Arithmetic-chiral realization]
\label{problem:bar}
Extend the divisor construction of Section~\ref{sec:divisor-bar} to
arithmetic local coefficients by constructing the following data and
compatibilities.  In particular, identify the maps from the divisor
algebras or their bars to the proposed arithmetic objects.

\begin{enumerate}
\item Fix one category
\((\mathcal C,\widehat\otimes,\one)\) as above, including the completion
functor or tensor topology.  For every place \(v\), give an augmented
\(E_1\)-algebra \((A_v,\eta_v,\epsilon_v)\), or give explicit right and
left boundary modules when no augmentation is intended.  At infinity,
also give the realization functor from the chosen chiral or factorization
algebra to \(\mathcal C\).

\item Construct each local bar through simplicial length three, with its
grading and signs.  For a finite set \(S\) of places containing infinity,
form the declared completed tensor product \(B_S\) of the local bars.

\item Use cohomological indexing for this twisting construction, so the
homological bar above is reindexed as \(B^{-n}=B_n\), with the same
differential.  Give a coaugmented conilpotent dg coalgebra
\[
 P_S\in\operatorname{CoAlg}_{\mathrm{conil}}(\mathcal C),
 \qquad \eta_{P_S}:\one\longrightarrow P_S,
\]
set
\(E_S=\underline{\operatorname{End}}_{\mathcal C}(B_S)\), and construct a
continuous degree-one map
\[
 \tau_S:P_S\longrightarrow E_S.
\]
Require the normalization \(\tau_S\circ\eta_{P_S}=0\).  Equip these
objects with complete separated descending filtrations for which the
differentials, coproduct, composition, and action on \(B_S\) are continuous.
Require the resulting twisting summand \(T_{\tau_S}\) either to be locally
nilpotent before completion or to satisfy
\(T_{\tau_S}(F^r)\subseteq F^{r+1}\), so that every iterated expression in
the completed tensor product converges.
For homogeneous maps use the declared convolution product and Hom
differential
\[
 (f\star g)(c)=
 \sum(-1)^{|g||c_{(1)}|}f(c_{(1)})g(c_{(2)}),
 \qquad
 d_{\mathrm{Hom}}f=d_Ef-(-1)^{|f|}f d_P.
\]
Prove the Maurer--Cartan equation
\[
 d_{\mathrm{Hom}}\tau_S+\tau_S\star\tau_S=0,
\]
write the induced twisted differential on
\(P_S\widehat\otimes B_S\) with all Koszul signs, and verify its square
directly.

\item If \(\tau_S\) is to come from tame symbols, construct its
chain-level factorization through a specified residue-unit group algebra
and an action of that algebra on \(B_S\).  A multiplicative reciprocity
identity alone is not a Maurer--Cartan proof.

\item For every inclusion \(S\subset T\), construct a chain map
\[
 j_{S,T}:(P_S\widehat\otimes B_S,D_{\tau_S})
 \longrightarrow(P_T\widehat\otimes B_T,D_{\tau_T})
\]
with strict or specified homotopy-coherent composition.  Prove that it
intertwines the twisting operators.  Only then form the filtered colimit
and its cohomology.  Any claimed concentration in degrees \(0,1,2\) must
be proved; three untwisted polynomial prime factors already have the
nonzero class \(\eta_p\wedge\eta_q\wedge\eta_r\) in degree three.

\item Relate a finite-place algebra to the Toeplitz carrier by an actual
representation on the dense space \(\bigoplus_{r\geq0}\mathbb C e_{p,r}\)
and an Euler derivation \(E_p\) satisfying
\[
 [N_p,\pi_p(a)]=\pi_p(E_pa)
\]
on a common invariant domain.  For the divisor algebra this part is
already supplied by \(\pi_p(u_p)=S_p\) and
\(E_p(u_p^r)=r\log p\,u_p^r\): direct evaluation on the finite basis span
gives \([N_p,S_p^r]=r\log p\,S_p^r\).
The additional arithmetic coefficient algebra must carry its own such
representation and comparison with this divisor model.
\end{enumerate}

No system with the stated arithmetic, archimedean-chiral, tame-symbol, and
transition properties is constructed here.  Until one is constructed, its
colimit and cohomology are specifications rather than arithmetic-chiral
invariants.
\end{openproblem}

\section{The Euler product as an analytic determinant}
\label{sec:euler-det}

There are two Hilbert spaces in an Euler product: one has a vector for
each prime, and the other has a vector for each occupation configuration.
For the two primes \(2,3\), put
\[
 K_{\{2,3\}}(s)=\operatorname{diag}(2^{-s},3^{-s})
 \quad\text{on }\mathbb C^2.
\]
The occupation basis is \(e_{a,b}\), with \(a,b\in\mathbb N_0\).
Its energy is \(a\log2+b\log3\), so for \(\operatorname{Re}s>0\),
\[
 \sum_{a,b\geq0}2^{-as}3^{-bs}
 =\det(I-K_{\{2,3\}}(s))^{-1}.
\]
At \(s=1\), the trace is \(3\) and the determinant is \(1/3\).
Passing to all primes requires convergence in an operator ideal.

\subsection{The prime operator and its Fock space}

Let \(\mathscr K=\ell^2(\mathcal P)\), where \(\mathcal P\) is the set
of rational primes.  Define
\[
 Ae_p=(\log p)e_p,\qquad
 \mathcal D(A)=\left\{x\in\mathscr K:
 \sum_p(\log p)^2|x_p|^2<\infty\right\}.
\]
The inner product is linear in its first argument.  This diagonal
operator is closed: convergence of a sequence and of its images gives
coordinatewise convergence, and the limiting image has coordinates
\((\log p)x_p\), placing the limit in the displayed domain.
Its adjoint has the same domain, since testing against the finitely
supported vectors forces those coordinates and their square summability.
Thus \(A\) is positive and self-adjoint.  The resolvent
\((A+I)^{-1}\) is the norm limit of finite diagonals because its
entries \(1/(1+\log p)\) tend to zero; hence it is compact.
For \(\operatorname{Re}s>0\), let
\[
 K(s)=e^{-sA},\qquad K(s)e_p=p^{-s}e_p.
\]
For an integer \(m\geq1\), the Schatten ideal \(\mathcal S_m\) consists
of compact operators with \(m\)-summable singular values.  The singular
values \(\mu_j(T)\) are the nonnegative square roots of the eigenvalues
of \(T^*T\), counted with multiplicity.  The ideal norm is
\(\|T\|_m=(\sum_j\mu_j(T)^m)^{1/m}\); \(\mathcal S_1\) is the
trace class.  For a trace-class diagonal \(Te_j=t_je_j\), the
Fredholm determinant used here is the convergent product
\(\det(I-T)=\prod_j(1-t_j)\).  Absolute summability of the entries
makes this the limit of the finite diagonal determinants.

\begin{theorem}[Analytic Euler determinant]
\label{thm:euler-fredholm}
For \(m\geq1\), the map \(s\mapsto K(s)\) is holomorphic with values
in \(\mathcal S_m\) precisely on the half-plane
\(\operatorname{Re}s>1/m\).  There,
\[
 \|K(s)\|_m^m=\sum_p p^{-m\operatorname{Re}s}.
\]
Let
\(\mathscr F_{\mathrm{sym}}(\mathscr K)=\bigoplus_{q\geq0}\operatorname{Sym}^q
\mathscr K\) be the Hilbert bosonic Fock space, with normalized occupation
basis \(|(r_p)_p\rangle\) indexed by finitely supported nonnegative
integers.  The summand \(q\) has particle number \(q=\sum_p r_p\);
the vacuum is the unit vector in the \(q=0\) summand.  The map
\[
 U:\mathscr F_{\mathrm{sym}}(\mathscr K)\longrightarrow\ell^2(\mathbb N^\times),
 \qquad U|(r_p)_p\rangle=e_{\prod_p p^{r_p}}
\]
is unitary.  Second quantization satisfies
\[
 U\left(\bigoplus_{q\geq0}K(s)^{\otimes_s q}\right)U^{-1}
 =e^{-sH},\qquad \operatorname{Re}s>0.
\]
These Fock operators are trace class exactly when
\(\operatorname{Re}s>1\), and on that half-plane
\[
 \operatorname{Tr}_{\mathscr F_{\mathrm{sym}}(\mathscr K)}
 \left(\bigoplus_{q\geq0}K(s)^{\otimes_s q}\right)
 =\det(I-K(s))^{-1}=\zeta(s).
\]
The determinant is the ordinary Fredholm determinant on \(\mathscr K\),
normalized to tend to \(1\) as real \(s\to+\infty\).
\end{theorem}

\begin{proof}
The singular values of \(K(s)\) are \(p^{-\operatorname{Re}s}\).
\label{proof:prime-harmonic}
The prime series \(\sum_p p^{-x}\) converges for \(x>1\) by comparison
with \(\sum_{n\geq2}n^{-x}\).  It diverges for \(x\leq1\).
Indeed, convergence of \(\sum_p1/p\) would bound all products
\(\prod_{p\leq R}(1-1/p)^{-1}\): the series
\(-\log(1-1/p)\) is at most \(2/p\).  Their geometric expansions contain
\(\sum_{n\leq R}1/n\), which is unbounded.  Comparison proves the
remaining divergence claims.

For a compact subset of \(\operatorname{Re}s>1/m\), choose
\(\sigma_0>1/m\) below its real parts.  The finite diagonal truncations
of \(K(s)\) converge uniformly there in \(\mathcal S_m\), because their
\(m\)-norm tails are bounded by \(\sum_{p>R}p^{-m\sigma_0}\).
They are holomorphic, so their limit is holomorphic in that Banach space.
At \(\operatorname{Re}s=0\), the diagonal operator is unitary and
noncompact.  For \(\operatorname{Re}s<0\), it is unbounded on its
maximal diagonal domain.

Unique factorization makes \(U\) a bijection of orthonormal bases.
On an occupation vector, second quantization has eigenvalue
\(\prod_p p^{-sr_p}\), which proves the operator equality.
Its trace-class domain follows from Proposition~\ref{prop:prime-product}.
For a finite set of primes \(F\), summing all occupations gives
\(\prod_{p\in F}(1-p^{-s})^{-1}\).  The projections onto configurations
supported in \(F\) tend to the identity.  For \(\operatorname{Re}s>1\),
the omitted absolute diagonal sum tends to zero, so the convergence is
in trace norm.  On \(\mathscr K\), the finite compressions converge in
\(\mathcal S_1\).  Their determinants converge to
\(\prod_p(1-p^{-s})\), the Fredholm determinant of the diagonal operator.
Both products converge locally uniformly.  This proves
the trace--determinant identity, including its analytic meaning and normalization.
\end{proof}

\subsection{Regularization and its exact boundary}

When \(0<\operatorname{Re}s\leq1\), the prime operator remains compact.
The divergent first terms of its logarithmic determinant can be removed.
For \(m\geq1\), define the entire function
\[
 E_m(z)=(1-z)\exp\left(\sum_{k=1}^{m-1}\frac{z^k}{k}\right),
\]
where an empty sum is zero.  Since \(E_m(z)-1=O(z^m)\) at zero,
\(E_m(K)-I\) is trace class for the diagonal \(K\in\mathcal S_m\)
used here.  Define
\[
 \det_m(I-K)=\det\bigl(E_m(K)\bigr).
\]
This fixes the regularization, including every exponential factor.

\begin{theorem}[Regularized prime determinants]
\label{thm:prime-regularization}
For every integer \(m\geq1\),
\[
 \Delta_m(s):=\det_m(I-K(s))
 =\prod_p E_m(p^{-s})
\]
is nonzero and holomorphic on \(\operatorname{Re}s>1/m\).  It has the
unique holomorphic logarithm tending to zero along the positive real axis,
\[
 L_m(s):=\log \Delta_m(s)
 =-\sum_p\sum_{k\geq m}\frac{p^{-ks}}{k}.
\]
If \(P(z)=\sum_p p^{-z}\) for \(\operatorname{Re}z>1\), then
\[
 \Delta_m(s)=\zeta(s)^{-1}
 \exp\left(\sum_{k=1}^{m-1}\frac{P(ks)}{k}\right),
 \qquad \operatorname{Re}s>1.
\]
For \(\sigma=\operatorname{Re}s>1/m\) and integer \(R\geq2\),
\[
 \left|L_m(s)-\sum_{p\leq R}\log E_m(p^{-s})\right|
 \leq\frac{R^{1-m\sigma}}
 {m(m\sigma-1)(1-2^{-\sigma})}.
\]
On the same initial half-plane, the full logarithmic derivative is
\[
 L_m'(s)=\sum_p(\log p)\frac{p^{-ms}}{1-p^{-s}}.
\]
The logarithms in the tail estimate use the power series
\(\log E_m(z)=-\sum_{k\geq m}z^k/k\) for \(|z|<1\).
\end{theorem}

\begin{proof}
For \(|z|\leq2^{-\sigma}<1\), the defining series gives
\[
 |\log E_m(z)|\leq\frac{|z|^m}{m(1-2^{-\sigma})}.
\]
Summation over primes proves absolute and locally uniform convergence
of the logarithmic series.  Exponentiation gives a nonzero holomorphic
product.  The trace-class diagonal \(E_m(K(s))-I\) has the same Fredholm
product, so this is the stated operator determinant.  Its limit at
\(+\infty\) is \(1\), and the logarithm has limit zero.
On \(\operatorname{Re}s>1\), add back the omitted terms and apply
Theorem~\ref{thm:euler-fredholm}.  Finally,
\[
 \sum_{p>R}p^{-m\sigma}
 \leq\sum_{n>R}n^{-m\sigma}
 \leq\int_R^\infty x^{-m\sigma}\,dx
\]
proves the stated tail bound.  On a smaller closed half-plane within
this domain, differentiation is justified by
\(\sum_{n\geq2}(\log n)n^{-m\sigma}<\infty\).
Differentiating each summand and summing its geometric series gives the
stated full logarithmic derivative.
\end{proof}

The two-prime example already distinguishes the normalizations:
\[
 \det_2(I-K_{\{2,3\}}(s))
 =\det(I-K_{\{2,3\}}(s))\exp(2^{-s}+3^{-s}).
\]
At \(s=1\), its value is \(e^{5/6}/3\).  A constant cannot remove this
parameter-dependent correction.

\begin{theorem}[Fractional zero at the convergence boundary]
\label{thm:fractional-zero}
For \(m\geq2\), there is a disk about \(a_m=1/m\) and a nonvanishing
holomorphic function \(G_m\) on that disk such that
\[
 \Delta_m(s)^m=(ms-1)G_m(s),\qquad \operatorname{Re}s>a_m.
\]
Moreover, \(G_m(a_m)>0\).  Thus \(\Delta_m\) has no single-valued meromorphic
continuation to a neighborhood of \(a_m\).  On the real axis,
\[
 \Delta_m(s)\sim C_m(ms-1)^{1/m}\quad(s\downarrow a_m),\qquad C_m>0.
\]
\end{theorem}

\begin{proof}
Put
\[
 R_0(z)=\sum_{k\geq2}\frac{P(kz)}k,\qquad
 B_m(s)=\sum_{k\geq m+1}\frac{P(ks)}k.
\]
The estimate in Theorem~\ref{thm:prime-regularization} shows that these
functions are holomorphic on \(\operatorname{Re}z>1/2\) and
\(\operatorname{Re}s>1/(m+1)\), respectively.
For \(\operatorname{Re}z>1\),
\[
 \log\zeta(z)=P(z)+R_0(z).
\]
The residue of \(\zeta\) at \(1\) is one.  One can see its normalization
directly from
\[
 \zeta(z)-\frac1{z-1}
 =\sum_{n\geq1}\left(n^{-z}-\int_n^{n+1}x^{-z}\,dx\right).
\]
The right side converges locally uniformly for \(\operatorname{Re}z>0\),
because its summands are \(O(n^{-\operatorname{Re}z-1})\) on compact
sets.  Consequently \(Q(z)=(z-1)\zeta(z)\) is holomorphic near \(1\),
with \(Q(1)=1\), and is nonzero on a sufficiently small disk.

The logarithmic series of Theorem~\ref{thm:prime-regularization} gives, for
\(\operatorname{Re}s>a_m\),
\[
 \Delta_m(s)^m
 =\zeta(ms)^{-1}\exp\bigl(R_0(ms)-mB_m(s)\bigr)
 =(ms-1)\frac{\exp\bigl(R_0(ms)-mB_m(s)\bigr)}{Q(ms)}.
\]
The last quotient is the required \(G_m\).  It is holomorphic and nonzero
near \(a_m\), with
\[
 G_m(a_m)=\exp\bigl(R_0(1)-mB_m(a_m)\bigr)>0.
\]
If \(\Delta_m\) were meromorphic there, the order of its \(m\)-th power
would be divisible by \(m\).  The displayed simple zero contradicts this.
Since every factor is positive for real \(s>a_m\), the asymptotic follows
with \(C_m=\exp(R_0(1)/m-B_m(a_m))\).
\end{proof}

The fractional zero suggests the coordinate \(t^m=ms-1\), in which
the determinant has an ordinary zero.  The next section constructs its
continuation surface and the comparison maps for different regularization
orders from the convergent remainder functions.

\section{Branched continuation of the prime determinant}
\label{sec:prime-continuation}

A fractional zero can become an ordinary zero after changing the
coordinate on the domain.  The square root of \(s-a\) returns with the
opposite sign after one counterclockwise circuit around \(a\).
On the coordinate surface \(s=a+t^2\), it is the single-valued function
\(t\).  For the prime determinant, the analogous construction will require
infinitely many different ramification orders.

Fix an integer \(m\geq2\) and retain precisely \(E_m\), \(L_m\), and
\(\Delta_m\) from Theorem~\ref{thm:prime-regularization}.
Thus the initial domain is \(\operatorname{Re}s>1/m\),
\(p^{-s}=\exp(-s\log p)\) uses the real logarithm of \(p\), and
\(L_m(s)\to0\) as real \(s\to+\infty\).
Our base point is \(s_*=2\), with the positive value \(\Delta_m(2)\).
The continuation concerns this fixed holomorphic germ; the summability
domain of the operator in Theorem~\ref{thm:euler-fredholm} stays fixed.

We work over \(\mathbb C\) and use its usual topology.  A meromorphic
function is locally a quotient of holomorphic functions.  At a point \(a\),
its order is the integer \(v\) in \(f(s)=(s-a)^v u(s)\), where \(u\)
is holomorphic and nonzero.  Positive orders are zeros and negative
orders are poles.  The divisor records these points with their integer
orders; its support is the set of points with nonzero order.
A covering map is locally a disjoint union of copies
of an open set.  A branched covering also permits local coordinates
\(s=a+t^q\), with \(q\) a positive integer.  We will construct these
coordinates and their gluing maps explicitly.

We use convergent complex power series and integration along piecewise
smooth paths.  The local logarithms and residue formula needed for the
continuation maps are developed next.

\subsection{Logarithms, residues, and zeta factors}

\subsubsection{Local logarithms and integrals}

A locally uniform limit of holomorphic functions is holomorphic.
For completeness, on a circle contained in the domain, uniform convergence
allows passage to the limit in the Cauchy integral formula.
Expanding \((z-s)^{-1}\) as a geometric series on a smaller concentric
disk gives a convergent power series for the limit.
The same argument gives locally uniform convergence of derivatives there.
The integral formula itself follows by integrating the power-series
primitive on small disks and subdividing the region between the circle
and a smaller circle about \(s\).  Internal edges cancel, while the
smaller circle contributes \(2\pi i f(s)\).

If \(u(a)\ne0\), shrink a disk about \(a\) until
\(|u(s)/u(a)-1|<1\).  The power series for \(\log(1+z)\), together
with one logarithm of \(u(a)\), gives a holomorphic logarithm of \(u\).
Two logarithms on a connected overlap differ by \(2\pi i n\),
\(n\in\mathbb Z\), because their exponentials agree and the difference
is continuous with discrete values.  Roots of a nonvanishing function
are obtained by dividing such a logarithm by the required integer.

For a closed path \(\gamma\) avoiding \(a\), define
\[
 \operatorname{wind}(\gamma,a)
 =\frac1{2\pi i}\int_\gamma\frac{ds}{s-a}.
\]
To see that this is an integer, continue a local logarithm of \(s-a\)
along finitely many small disks covering the path.  At the endpoint its
change is \(2\pi i n\).  Integrating its derivative gives the formula.
Subdividing a homotopy into small rectangles shows that the integral
of a holomorphic one-form is unchanged under path homotopy with fixed
endpoints: each rectangle lies in a disk with a power-series primitive.

We need a related residue formula.  Suppose a one-form on a half-plane
is holomorphic except at a locally finite set of simple poles \(a\),
with residues \(r_a\).  A closed path lies inside a compact convex polygon
in that half-plane whose boundary avoids the poles.
There are only finitely many poles in the polygon.
Subtract \(\sum_a r_a\,ds/(s-a)\).  The remaining form is holomorphic
on a slightly smaller polygon containing the path and has zero integral
around closed paths by subdivision.  It follows that
\[
 \int_\gamma\omega
 =2\pi i\sum_a r_a\operatorname{wind}(\gamma,a).
\]
Only finitely many terms contribute.  These elementary integral identities
will determine the entire continuation space.

\subsubsection{Continuation of the ordinary zeta function}

Write \(\Hp_c=\{s\in\mathbb C:\Re s>c\}\), and \(\Hp=\Hp_0\).
For \(\Re z>1\), let \(\zeta(z)=\sum_{n\geq1}n^{-z}\).
Finite geometric products and unique factorization give
\[
 \zeta(z)=\prod_p(1-p^{-z})^{-1},\qquad
 Z(z):=\log\zeta(z)=\sum_p\sum_{j\geq1}\frac{p^{-jz}}j.
\]
Both limits are locally uniform and absolute on \(\Hp_1\).
For example, all omitted integers in a finite prime product are bounded
in absolute sum by the tail of \(\sum n^{-\Re z}\).
The absolute logarithmic sum is bounded by
\((1-2^{-\Re z})^{-1}\sum_{n\geq2}n^{-\Re z}\).
Its exponential is the product, so \(\zeta\) is nonzero on \(\Hp_1\).
The branch \(Z\) tends to zero on the positive real axis.

The subtraction formula in the proof of
Theorem~\ref{thm:fractional-zero} continues \(\zeta\) meromorphically
to \(\Hp\), with its only pole at \(1\), simple and of residue one.
Its locally uniform derivative bound proves holomorphy of the subtracted
series throughout \(\Hp\).  This construction is sufficient for every
zeta factor used below.

Zeros are isolated.  A holomorphic function that is not identically zero
has a first nonzero Taylor coefficient at each zero, which factors it as
\((z-a)^v u(z)\) with \(u(a)\ne0\).  An accumulation of zeros would
make every coefficient vanish and then, by overlapping disks, make the
function identically zero.  The zeta function is nonzero on \(\Hp_1\),
so this cannot occur.

We will also need the absence of zeros on the positive real axis.
For \(x>0\), the paired series
\[
 \eta(x)=\sum_{r\geq1}\bigl((2r-1)^{-x}-(2r)^{-x}\bigr)
\]
converges and is strictly positive.  Its complex analogue converges
locally uniformly on \(\Hp\) by the same derivative estimate.
Splitting even and odd terms for \(\Re z>1\), and then continuing,
gives \(\eta(z)=(1-2^{1-z})\zeta(z)\).
Thus \(\zeta(x)<0\) for \(0<x<1\), while \(\zeta(x)>0\) for \(x>1\).

\subsection{The first continuation surface}

The proof of Theorem~\ref{thm:fractional-zero} isolates the first
divergent part without continuing the prime series itself.
Retain its functions \(P\), \(R_0\), and \(B_m\), holomorphic on
\(\Hp_1\), \(\Hp_{1/2}\), and \(\Hp_{1/(m+1)}\), respectively.
The same calculation gives, on \(\Hp_{1/m}\),
\[
 mL_m(s)=-P(ms)-mB_m(s)
 =-Z(ms)+R_0(ms)-mB_m(s).
\]
Consequently the single-valued meromorphic function
\[
 F_m(s)=\frac{\exp(R_0(ms)-mB_m(s))}{\zeta(ms)}
 \quad\text{on }U_m:=\Hp_{1/(m+1)}
\]
extends \(\Delta_m^m\).  The stated domain follows because
\(m/(m+1)>1/2\).  The numerator is holomorphic and nonzero throughout it.
The function \(F_m\) has a simple zero at \(a_m=1/m\).
At a point \(\rho/m\) where \(\zeta(\rho)=0\) has multiplicity \(v\)
and \(\Re\rho>m/(m+1)\), it has a pole of order \(v\).
There are no other zeros or poles on \(U_m\).

\begin{theorem}[Minimal first surface]
\label{thm:first-surface}
Let \(U_m^\circ\) be \(U_m\) with the divisor support of \(F_m\)
removed.  The equation
\[
 Y_m=\{(s,y):s\in U_m^\circ,\ y\in\mathbb C^*,\ y^m=F_m(s)\}
\]
defines a connected \(m\)-sheeted holomorphic cover.
It carries the continuation of \(\Delta_m\) as the function \(y\),
with the initial sheet specified by \(y(2)=\Delta_m(2)>0\).
It admits a meromorphic branched filling over every divisor point.
If \(v=\ord_a F_m\ne0\) and \(d=\gcd(m,|v|)\), the filling has
\(d\) points over \(a\).  Each has ramification index \(m/d\), and
\(y\) has order \(v/d\) there.
At \(a_m\), this is the minimal local \(m\)-fold cover, and \(y\)
has a simple zero.
\end{theorem}

\begin{proof}
On a small disk avoiding the divisor, choose \(\log F_m\).
The \(m\) functions
\(y_j=\exp((\log F_m+2\pi i j)/m)\), \(0\leq j<m\), are disjoint
local sheets.  A circuit about the simple zero at \(a_m\) permutes them
by \(j\mapsto j+1\).  Thus all sheets are connected by continuation.

Near a divisor point, write \(F_m(s)=(s-a)^v u(s)\) and choose
\(h^m=u\) with \(h\) holomorphic and nonzero.  The charts are
\[
 s=a+t^{m/d},\qquad
 y=\xi\,t^{v/d}h(a+t^{m/d}),\qquad \xi^m=1.
\]
Negative powers of \(t\) mean a meromorphic pole.  Changing \(t\) by
an \((m/d)\)-th root of unity identifies exactly \(m/d\) choices of
\(\xi\), because \(v/d\) and \(m/d\) are coprime.  Hence there are
\(d\) charts.  They parametrize all the punctured sheets and fill their
ends with smooth disk coordinates.  On overlaps their changes of
coordinate are the roots of unity just described, so they glue.

At \(a_m\), write \(F_m(s)=(ms-1)G_m(s)\), where
\[
 G_m(s)=\frac{\exp(R_0(ms)-mB_m(s))}{(ms-1)\zeta(ms)}.
\]
The denominator is nonzero near \(a_m\) and equals one there.
Thus \(G_m(a_m)>0\).  The coordinate \(t^m=ms-1\) gives
\(y=t\,G_m((1+t^m)/m)^{1/m}\), using the root positive at \(t=0\).
One circuit multiplies \(y\) by \(e^{2\pi i/m}\).
On any local cover \(s-a_m=t^e\) where \(y\) is meromorphic,
its order must be \(e/m\), an integer.  Therefore \(m\) divides \(e\).
Our chart attains this bound.
\end{proof}

For \(m=2\), the two sheets meet smoothly above \(1/2\), and one
circuit changes the sign.  At a zero of \(\zeta(2s)\) of order two,
the same construction gives two unramified simple poles.
A simple zero of \(\zeta(2s)\) would instead give one ramified simple
pole upstairs.  These are conditional descriptions of any such zeros
in \(U_2\), not assertions that they exist there.

\subsection{Finite extraction and continuation through the positive half-plane}

The preceding surface uses the full domain of the two displayed remainder
series.  To go farther left, remove finitely many more logarithmic factors.
The coefficients determine every branch.

In the extraction and the cover construction that follows, allow any
integer \(m\geq1\); order one will recover the reciprocal zeta function.
Define the integer \(\mu(n)\) as zero when a prime square divides \(n\),
and as \((-1)^r\) when \(n\) is a product of \(r\) distinct primes.
Set \(\mu(1)=1\).  Choosing subsets of the prime divisors proves
\(\sum_{d\mid n}\mu(d)=0\) for \(n>1\), and the sum is one for \(n=1\).
For \(n\geq1\), put
\[
 c_m(n)=\frac1n\sum_{\substack{k\mid n\\ k\geq m}}\mu(n/k).
\]
Then \(c_m(n)=0\) for \(n<m\), \(c_m(m)=1/m\), and
\(|c_m(n)|\leq1\).  Summing over divisors in the opposite order gives
\[
 \sum_{d\mid n}d\,c_m(d)=\mathbf1_{n\geq m}.
\]
Indeed, the contribution for a fixed \(k\mid n\), \(k\geq m\), is
\(\sum_{e\mid n/k}\mu(e)\), so only \(k=n\) survives.

\begin{theorem}[Finite logarithmic extraction]
\label{thm:extraction}\label{mono:prop:extraction}
On \(\Hp_{1/m}\),
\[
 L_m(s)=-\sum_{n\geq m}c_m(n)Z(ns).
\]
For \(N\geq m\), the remainder
\[
 T_{m,N}(s)=-\sum_{n>N}c_m(n)Z(ns)
\]
is single-valued and holomorphic on \(U_N=\Hp_{1/(N+1)}\).
On every open disk in \(U_N\) avoiding zeros and poles of
\(\zeta(ns)\), \(m\leq n\leq N\), the formula
\[
 f_{m,N}(s)=\exp\left(T_{m,N}(s)
       -\sum_{n=m}^N c_m(n)\ell_n(s)\right),
 \qquad e^{\ell_n(s)}=\zeta(ns),
\]
gives a local branch of the continued determinant, when its logarithms
are continued from the initial half-plane.
The formulas for successive \(N\) agree with these continued choices.
\end{theorem}

\begin{proof}
For \(x>1\) and \(\Re z=x\),
\[
 |Z(z)|\leq\frac1{1-2^{-x}}\sum_{r\geq2}r^{-x}
 \leq\frac{2^{-x}+2^{1-x}/(x-1)}{1-2^{-x}}.
\]
For \(x\geq2\), this is at most \(4\,2^{-x}\).
On a compact subset of \(U_N\), all \(ns\), \(n>N\), lie in
\(\Hp_1\).  The large \(n\) terms therefore have a geometric bound,
which proves local uniform convergence of \(T_{m,N}\) and its derivatives.

On \(\Hp_{1/m}\), expand each \(Z(ns)\) into its prime series.
Absolute convergence follows from the same bound, so sums can be reordered.
The coefficient of \(p^{-js}\) becomes
\[
 -\sum_{n\mid j}\frac{c_m(n)}{j/n}
 =-\frac{\mathbf1_{j\geq m}}j.
\]
This is the defining coefficient of \(L_m\).
The local formula follows by keeping its first \(N\) terms separate.
On \(U_N\),
\[
 T_{m,N}=T_{m,N+1}-c_m(N+1)Z((N+1)s).
\]
The last logarithm has its original single-valued branch there.
For \(m=1\), divisor cancellation gives
\(c_1(n)=\mathbf1_{n=1}\), so the same identity reduces to
\(L_1=-Z\).  The compatibility therefore holds for every positive
integer order.
\end{proof}

For example, \(c_2(n)=-\mu(n)/n\) for \(n>1\).
The first coefficients are
\[
 c_2(2)=\frac12,\quad c_2(3)=\frac13,\quad c_2(4)=0,
 \quad c_2(5)=\frac15,\quad c_2(6)=-\frac16.
\]
The first continuation takes a square root of an inverse zeta factor.
The next takes a cube root as well.  The negative coefficient at six
reverses the contribution of the pole of \(\zeta(6s)\).

\subsubsection{The logarithmic derivative and coincident divisors}

Let \(S_m\subset\Hp\) consist of points where \(\zeta(ns)\) has a
zero or pole for some \(n\geq m\) with \(c_m(n)\ne0\).
This set is locally finite.  On a compact set with \(\Re s\geq\epsilon>0\),
only \(n\leq1/\epsilon\) can contribute, since \(\zeta\) is nonzero
and holomorphic on \(\Hp_1\).  Each of these finitely many factors
has finitely many zeros and poles on the compact set.
Put \(\Omega_m=\Hp\setminus S_m\).

The locally finite punctures do not disconnect the half-plane.
A polygonal path between two points meets only finitely many punctures
in a compact neighborhood, and small detours avoid them.
Hence \(\Omega_m\) is path connected.

On \(\Omega_m\), define the holomorphic one-form
\[
 \omega_m=dT_{m,N}
 -\sum_{n=m}^N c_m(n)\frac{n\zeta'(ns)}{\zeta(ns)}\dd
 \quad\text{where }s\in U_N.
\]
Theorem~\ref{thm:extraction} proves independence of \(N\).
For every \(a\in\Hp\), define the rational number
\[
 \alpha_m(a)=-\sum_{n\geq m}c_m(n)\ord_a\zeta(ns).
\]
Only finitely many terms occur.  This is the residue of the meromorphic
extension of \(\omega_m\), with value zero wherever that extension is
holomorphic.  The signs include a zeta pole with order
\(-1\).  Coincident zeros and poles contribute to the same sum.
Each logarithmic derivative has a simple pole, so if this sum vanishes,
\(\omega_m\) extends holomorphically at \(a\).
For a loop \(\gamma\) in \(\Omega_m\), set
\[
 \chi_m(\gamma)=\exp\left(\int_\gamma\omega_m\right)
 =\exp\left(2\pi i\sum_{a\in S_m}
       \alpha_m(a)\operatorname{wind}(\gamma,a)\right).
\]
This is a homomorphism from based loops modulo homotopy to
the unit complex numbers.  Every value is a root of unity.
It is the monodromy character of the determinant.

\subsection{The minimal cover and the determinant line}

A Riemann surface is a Hausdorff space with a countable atlas of open
subsets of \(\mathbb C\) and holomorphic coordinate changes.
A pointed germ here means a holomorphic function near \(2\), with
two such functions identified if they agree on a smaller neighborhood.
A torsor for a group is a set on which the group acts so that exactly
one group element carries any specified point to any other.
For loops based at \(2\), multiplication means concatenation, and
reversing a path gives the inverse.  The kernel of \(\chi_m\) consists
of loops with character one.

\begin{theorem}[Continuation on a minimal cover]
\label{thm:minimal-cover}\label{mono:prop:local}
There is a connected holomorphic covering
\(\pi:\mathcal E_m\to\Omega_m\) and a nonzero holomorphic function
\(\mathcal D_m\) on it, with the following properties.
The pointed germ at \(2\) is \(\Delta_m\), and
\[
 \frac{d\mathcal D_m}{\mathcal D_m}=\pi^*\omega_m.
\]
The fiber of \(\pi\) is a torsor for \(\operatorname{im}\chi_m\).
Every pointed connected cover on which that germ becomes single-valued
maps uniquely to \(\mathcal E_m\) over \(\Omega_m\), preserving the germ.
Thus \(\mathcal E_m\) has the fewest sheets among such covers.
It has a meromorphic branched filling over every point of \(S_m\).
If \(\alpha_m(a)=b/q\) in lowest terms with \(q>0\), each local end
has ramification index \(q\), and \(\mathcal D_m\) has order \(b\).
For \(\alpha_m(a)=0\), take \(b=0,q=1\).
Every meromorphic realization with local ramification index \(e\)
requires \(q\mid e\).
\end{theorem}

\begin{proof}
Consider paths \(\gamma\) in \(\Omega_m\) from \(2\).
Two paths are equivalent when they have the same endpoint and
\(\exp(\int_\gamma\omega_m-\int_\eta\omega_m)=1\).
Write \([\gamma]\) for the class, and set
\[
 \pi[\gamma]=\gamma(1),\qquad
 \mathcal D_m[\gamma]=\Delta_m(2)
       \exp\left(\int_\gamma\omega_m\right).
\]
The equivalence relation makes the function well-defined.
For a small disk \(V\subset\Omega_m\) about \(\gamma(1)\), append
paths inside \(V\).  A primitive of \(\omega_m\) on \(V\) shows
that the resulting classes form a copy of \(V\).
Two such copies are either disjoint or identical over a connected overlap,
with change of coordinates equal to the identity on the base.
These charts define the covering and its complex structure.
They separate distinct points in a fiber, so the space is Hausdorff.
There are at most countably many classes in a fiber, because finite sums
of rational residues give a countable monodromy image.
A countable disk basis on the base therefore gives a countable atlas.
Every class is joined to the base class by the initial portions of its
representing path.  Thus the cover is connected.

The displayed exponential is holomorphic in each chart, with the stated
logarithmic derivative.  On the initial half-plane it agrees with
\(\Delta_m\), since the logarithmic derivatives and the values at \(2\)
agree.  Appending loops gives the simply transitive action of
\(\operatorname{im}\chi_m\) on each fiber.

Let \(X\to\Omega_m\) be another pointed connected cover carrying the
continued germ.  A path in \(X\) from its base point projects to a path
\(\gamma\).  Its endpoint value is forced by the exponential formula.
Two paths ending at the same point therefore have trivial exponential
period.  Sending that point to \([\gamma]\) is well-defined, holomorphic,
and unique.  It is onto because every path on the base lifts to \(X\).
This proves minimality.

Near a puncture, the Laurent calculation gives
\[
 \omega_m=\frac bq\frac{ds}{s-a}+dh(s)
\]
for a holomorphic \(h\) on a disk.  Each branch is consequently
\(C(s-a)^{b/q}e^{h(s)}\), with \(C\ne0\).
The punctured sheets form cycles of length \(q\), since \(b,q\) are
coprime.  Fill each cycle with the coordinate \(s=a+t^q\).
The function becomes \(Ct^b e^{h(a+t^q)}\), which is meromorphic.
Different choices of the coordinate differ by a \(q\)-th root of unity.
This gives the compatible local filling and the stated order.
On a chart \(s-a=t^e\), a meromorphic realization would have integer
order \(eb/q\); coprimality of \(b,q\) forces \(q\mid e\).
\end{proof}

The formula resolves cancellations before assigning ramification.
For example, local exponents \(1/2\) and \(-1/2\) give a nonzero
holomorphic branch, not a double cover.  A total exponent \(2/3\)
gives a triple cover with a double zero.  An exponent \(-2/3\) gives
the same cover with a double pole.  The denominator of the reduced total
exponent, rather than any individual factor, is the minimal local degree.

\subsubsection{A finite description on each half-plane}

The construction has finite equations on every \(U_N\).
Let \(M_N\) be the least common multiple of the denominators of
\(c_m(m),\ldots,c_m(N)\), in lowest terms.  Define
\[
 A_{m,N}(s)=\exp(M_N T_{m,N}(s))
       \prod_{n=m}^N\zeta(ns)^{-M_Nc_m(n)}.
\]
This is meromorphic on \(U_N\), since all powers are integers.
Over \(U_N\setminus S_m\), the equation
\(y^{M_N}=A_{m,N}(s)\) is connected and has exactly \(M_N\) sheets.
To prove this, first note the identity
\[
 \alpha_m(1/n)=c_m(n),\qquad n\geq m.
\]
Only \(\zeta(ns)\) has a pole at \(1/n\); all other scaled zeta
values there are nonzero finite values on the positive real axis.
For \(m\leq n\leq N\), these points lie in \(U_N\).
Small loops about them generate the subgroup of roots of unity generated
by \(\exp(2\pi i c_m(n))\).  Its order is \(M_N\): a positive integer
annihilates every exponent modulo \(\mathbb Z\) exactly when it is
divisible by every reduced denominator.  Thus continuation reaches all
the roots of the polynomial equation.  Conversely every residue on
\(U_N\) has denominator dividing \(M_N\), so no further sheets occur.
The path construction identifies this equation, with base value
\(y(2)=\Delta_m(2)\), with the minimal continuation using paths in
that half-plane.

For order two, the first degrees are
\[
 \begin{array}{c|ccccc}
 N&2&3&4&5&6\\ \hline
 M_N&2&6&6&30&30
 \end{array}
\]
The point \(1/4\) introduces no sheet, since \(c_2(4)=0\).
At \(1/6\), the exponent \(-1/6\) gives a simple pole after the
sixfold local pullback.  Both facts occur within the degree-thirty surface
over \(U_6\).  For order three, the exponents \(1/3\) and \(1/4\)
already force degree twelve over \(U_4\).

Since \(M_N\mid M_{N+1}\), the remainder identity implies
\[
 A_{m,N+1}|_{U_N}=A_{m,N}^{\,M_{N+1}/M_N}.
\]
The transition between the pointed components is exactly
\((s,y)\mapsto(s,y)\).  Its image is the pointed component over
\(U_N\) inside the next surface.  The maps also extend over filled
divisor points, since both use the same reduced residue and local coordinate.

Every path in \(\Omega_m\) is compact and has positive minimum real
part.  It therefore lies in some \(U_N\).  The union of these compatible
pointed surfaces is \(\mathcal E_m\).  This proves continuation through
the whole positive half-plane with its described branched fillings.
It is the full domain supplied by this exhaustion of convergent
remainders.  No assertion of maximal continuation outside \(\Hp\) is
needed or proved.

\subsubsection{Branches as a section of a line bundle}

There is also a description without enlarging the base.
Choose small disks \(V_i\) covering \(\Omega_m\), and continue the
initial determinant to a branch \(f_i\) on each disk along a chosen path.
On a connected component of \(V_i\cap V_j\), put
\[
 g_{ij}=f_i/f_j.
\]
Both branches have logarithmic derivative \(\omega_m\), so \(g_{ij}\)
is constant.  It is a value of \(\chi_m\), and
\(g_{ij}g_{jk}=g_{ik}\) on triple overlaps.
Identify \((s,v_j)\in V_j\times\mathbb C\) with
\((s,v_i)\in V_i\times\mathbb C\) when \(v_i=g_{ij}v_j\).
The cocycle identity makes the quotient a holomorphic line bundle
\(\mathcal L_m\).  The values \(f_i\) define a nonzero holomorphic
section of it.  All bundle and section maps have now been specified.

Because the transition constants have zero derivative, ordinary
differentiation of local coordinates defines a flat connection \(\nabla\)
on \(\mathcal L_m\).  Flat means that continuing a locally constant
coordinate around a sufficiently small loop returns the same coordinate.
The determinant section satisfies \(\nabla f=\omega_m f\).
Its nonvanishing trivializes the underlying holomorphic line bundle.
In that trivialization, the connection is \(d+\omega_m\).
Its parallel transport around \(\gamma\) multiplies by
\(\chi_m(\gamma)^{-1}\).  Thus the continuation information lies in the
specified connection as well as the section; it is not a claim that the
underlying holomorphic line bundle is nontrivial.
The branched charts of Theorem~\ref{thm:minimal-cover} describe its
meromorphic extension after the required local pullback.

\subsection{Changing the regularization order}

The order of regularization changes the function.  The comparison must
therefore include a multiplier and the space on which it is defined.
For \(n>m\), define on \(\Hp_{1/m}\)
\[
 H_{m,n}(s)=\exp\left(\sum_{k=m}^{n-1}\frac{P(ks)}k\right).
\]
Subtracting the absolutely convergent defining series gives
\[
 \Delta_n=H_{m,n}\Delta_m,\qquad
 H_{m,n}H_{n,r}=H_{m,r}\quad(m<n<r).
\]
For instance, \(H_{2,3}=\exp(P(2s)/2)\).
At \(s=1/2\), its continued local exponent is \(-1/2\), canceling
the exponent \(1/2\) of \(\Delta_2\).  The resulting \(\Delta_3\)
is already holomorphic and nonzero there by its initial definition.

\begin{theorem}[Comparison on common covers]
\label{thm:order-comparison}
On \(\Omega=\Omega_m\cap\Omega_n\), the multiplier \(H_{m,n}\)
continues with logarithmic derivative \(\omega_n-\omega_m\) and
character \(\chi_n\chi_m^{-1}\).  On the pointed cover defined by
\(\ker\chi_m\cap\ker\chi_n\), all three functions are single-valued
and \(\mathcal D_n=H_{m,n}\mathcal D_m\).
At a divisor point, the multiplier has local exponent
\(\alpha_n-\alpha_m\).  The comparisons preserve their cocycle identity
on common covers and in the corresponding line bundles.

Multiplication descends to a map from the minimal \(m\)-cover to the
minimal \(n\)-cover precisely when
\(\ker\chi_m\subseteq\ker\chi_n\), with characters restricted to
\(\Omega\).  Here each minimal cover uses paths in \(\Omega\) from
the fixed base point.  Such descent is not automatic.
\end{theorem}

\begin{proof}
The initial identity gives the difference of logarithmic derivatives and
fixes the multiplier at \(2\).  Path integration therefore defines its
unique continuation.  The exponential period is the quotient of the two
characters.  It is trivial on their common kernel, so multiplication is
well-defined on the stated common cover.  The local exponent follows by
subtracting residues.  For three orders, the logarithmic derivatives
telescope and the base values multiply correctly.  Uniqueness of the
continued germ proves the cocycle identity.

For the line bundles, choose the same continuation paths on a common disk
cover.  The multiplier has transition constants
\(g_{ij}^{(n)}/g_{ij}^{(m)}\), so it is a nonzero section of
\(\mathcal L_n\otimes\mathcal L_m^{-1}\).
Multiplying local coordinates gives an isomorphism of holomorphic line
bundles carrying the determinant section to the other one.  The difference
of the two flat connections is the specified logarithmic derivative.

A point of the minimal \(m\)-cover identifies paths whose quotient loop
lies in \(\ker\chi_m\).  Sending it to the corresponding path class
in the \(n\)-cover is well-defined exactly when every such loop lies
in \(\ker\chi_n\).  This proves the descent criterion in both directions.
\end{proof}

For a concrete failure, compare orders \(2\) and \(3\) around \(s=1/4\).
The coefficient formulas give \(c_2(4)=0\) and \(c_3(4)=1/4\).
All other scaled zeta factors are finite and nonzero at that positive
real point.  Thus \(\chi_2=1\) and \(\chi_3=i\) on its small positive
loop.  The order-three branch cannot be a function on the minimal
order-two cover near this loop.  The common cover in the theorem supplies
the required comparison without discarding this monodromy.

\subsection{Why the number of sheets is unbounded}

The finite equations determine the whole sheet group, including its
prime-power orders.  For a prime \(p\), let \(e_p\) be the least
positive integer such that \(p^{e_p}\geq m\).  Equivalently,
\(p^{e_p-1}<m\leq p^{e_p}\).
For a positive integer \(d\), let \(\mu_d\subset\mathbb C^*\) be the
group of \(d\)-th roots of unity.

\begin{theorem}[Exact monodromy group]
\label{thm:monodromy-group}
For \(m\geq2\), the image of the continuation character is
\[
 \operatorname{im}\chi_m
 =\bigcup_{N\geq m}\mu_{M_N}
 =\{z\in\mathbb C^*:z\text{ has finite order }d,
       \ v_p(d)\leq e_p\text{ for every prime }p\}.
\]
Here \(v_p(d)\) is the exponent of \(p\) in the prime factorization of
\(d\).  As an abstract group, this is
\(\bigoplus_p\mathbb Z/p^{e_p}\mathbb Z\): only finitely many
coordinates are nonzero in each element.
\end{theorem}

\begin{proof}
Every loop lies in some \(U_N\), whose monodromy image is exactly
\(\mu_{M_N}\).  This proves the first equality.
For \(n>1\), divisor cancellation rewrites the coefficient as
\[
 c_m(n)=-\frac1n\sum_{\substack{k\mid n\\k<m}}\mu(n/k).
\]
If \(v_p(n)>e_p\), then \(k<m\) implies \(v_p(k)\leq e_p-1\).
Thus \(p^2\mid n/k\) in every term, so \(c_m(n)=0\).
Every nonzero coefficient consequently has denominator with
\(p\)-valuation at most \(e_p\).

Conversely, take \(n=p^{e_p}\).  Its only divisor at least \(m\) is
\(n\) itself, so \(c_m(n)=1/p^{e_p}\).
The loop about \(1/p^{e_p}\) gives a generator of \(\mu_{p^{e_p}}\).
These loops for finitely many distinct primes generate all roots with
order dividing their product.  Indeed, cyclic groups of coprime orders
intersect trivially and their product has the product order.
This proves the second equality and the direct-sum description.
\end{proof}

For \(m=2\), every \(e_p=1\), so the possible orders are precisely
the squarefree positive integers.  For \(m=3\), one has \(e_2=2\),
while \(e_p=1\) for every odd prime.  The allowed orders divide four
times an odd squarefree integer.  Increasing the regularization can change
this group, as well as the local determinant values.

\begin{theorem}[Prime ramification and the boundary at zero]
\label{thm:infinite-degree}
Fix \(m\geq2\).  For every prime \(q\geq m\),
\[
 \alpha_m(1/q)=1/q.
\]
Hence the minimal continuation cover over \(\Omega_m\) has infinite
degree.  No finite-degree branched cover of the whole positive half-plane
carries the continued determinant as a meromorphic function.
Nor can one finite-degree branched cover of a neighborhood of zero carry
all these continued germs at the positive points \(1/q\).
\end{theorem}

\begin{proof}
At \(s=1/q\), the factor \(\zeta(ns)\) has a pole only for \(n=q\).
Every other \(n/q>0\) is a nonzero finite zeta value by the real-axis
calculation above.  Since \(q\) is prime and \(m\geq2\), its only
divisor at least \(m\) is \(q\), so \(c_m(q)=1/q\).
The residue formula gives the claimed exponent.

A small loop about \(1/q\) therefore has character
\(e^{2\pi i/q}\), of order \(q\).  Infinitely many primes exceed \(m\):
if only finitely many did, all primes would form a finite list, while
one plus their product has a prime divisor outside that list.
Thus the image of \(\chi_m\) is infinite.

On a meromorphic realization over a local chart
\(s-1/q=t^e\), the order is \(e/q\), which must be an integer.
Therefore its ramification index is at least \(q\).
A finite-degree cover has all local indices bounded by its degree.
Arbitrarily large primes contradict such a bound, both on \(\Hp\)
and in any neighborhood of zero containing their reciprocals.
\end{proof}

Multiplying by a single-valued meromorphic factor cannot remove this
obstruction.  Its local order is an integer, so it changes
\(\alpha_m(a)\) by an integer and leaves the character unchanged.
The result concerns this determinant germ and its explicitly constructed
continuation.  It neither identifies a global arithmetic cohomology nor
asserts that the entire imaginary axis is a natural boundary.
The continuation through \(\Hp\), its minimal local ramification,
and its order-comparison maps are determined by the formulas above.


\section{Finite covers and domain restrictions}
\label{sec:finite-covers}

The continuation of a determinant depends on the permitted paths.
A germ may have finitely many branches under paths in one domain even
when paths that leave the domain produce infinitely many additional values.
We determine the least cover for the first choice and the components
created by the second.  The definitions of \(\Delta_m,c_m,\omega_m,\alpha_m\)
and their local charts are those of Section~\ref{sec:prime-continuation},
including \(m=1\).  Write \(\mathcal G_m=\operatorname{im}\chi_m\).
A finite branched cover of degree \(d\) is a proper surjective holomorphic
map of connected Riemann surfaces with discrete branch values and total
local degree \(d\) over every point.  Proper means that inverse images
of compact sets are compact; the total local degree is the sum of the
ramification indices in a fibre.

\subsection{An exact criterion for a finite cover}

Let \(D\subset\Hp\) be a connected open set, and put
\(D^\circ=D\setminus S_m\).
Choose \(b\in D^\circ\) and a path in \(\Omega_m\) from \(2\) to \(b\).
It fixes a nonzero continued germ \(f_b\).
Restrict the differential and define
\[
 \mathcal G_m(D)=\left\{\exp\int_\gamma\omega_m:
       \gamma\text{ is a loop in }D^\circ\text{ based at }b\right\}.
\]
Changing the initial path changes \(f_b\) by a constant in \(\mathcal G_m\),
but it leaves this subgroup unchanged.
Let \(\mathcal M(D)\) denote the field of single-valued meromorphic
functions on \(D\).
A germ is algebraic over \(\mathcal M(D)\) if it satisfies a nonzero
polynomial with coefficients in that field.
Its algebraic degree is the least degree of such a polynomial.
A polynomial is irreducible if it has no factorization into polynomials
of smaller positive degrees over the same field.

\begin{theorem}[Finite-cover criterion]
\label{mono:thm:criterion}
The following conditions are equivalent.
\begin{enumerate}
\item The group \(\mathcal G_m(D)\) is finite.
\item There is an integer \(k\geq1\) such that
\(k\int_\gamma\omega_m\in2\pi i\mathbb Z\) for every loop in \(D^\circ\).
\item Some positive integer power of \(f_b\) extends meromorphically to \(D\).
\item A connected finite branched cover of \(D\) carries a meromorphic
function agreeing with the continued germ near a point over \(b\).
\item The germ \(f_b\) is algebraic over \(\mathcal M(D)\).
\end{enumerate}
When these conditions hold, let \(k=|\mathcal G_m(D)|\).
Then \(k\) is the least possible power in condition (iii), the least
cover degree in condition (iv), and the algebraic degree in condition (v).
There is \(A_D\in\mathcal M(D)\) such that the minimal cover is the
branched filling of \(y^k=A_D(s)\), with \(y(b)=f_b(b)\).
The polynomial \(X^k-A_D\) is irreducible over \(\mathcal M(D)\).
\end{theorem}

\begin{proof}
A finite subgroup of \(\mathbb C^*\) of size \(k\) equals the group
\(\mu_k\) of \(k\)-th roots of unity.
Indeed, multiplying all its elements before and after multiplication
by one element \(g\) gives \(g^k=1\).
There are exactly \(k\) roots, so the inclusion is equality.
Conversely a subgroup annihilated by \(k\) has at most \(k\) elements.
This proves the equivalence of (i) and (ii).

Under (ii), define on \(D^\circ\)
\[
 A_D(s)=f_b(b)^k\exp\left(k\int_b^s\omega_m\right).
\]
The period condition makes it path independent and holomorphic.
At a puncture \(a\in D\cap S_m\), a small positive loop gives
\(k\alpha_m(a)\in\mathbb Z\).
Theorem~\ref{mono:prop:local} consequently makes \(A_D\) meromorphic there.
This proves (iii).  Conversely, continuing a meromorphic power around a
loop gives \(g^k=1\), so (iii) implies (ii).

For \(k=|\mathcal G_m(D)|\), the equation \(y^k=A_D\) over \(D^\circ\)
is a connected \(k\)-sheeted cover.
Local roots give its sheets, and continuation realizes every element
of \(\mu_k\), so every sheet is reached from the specified one.
For completeness, at a zero or pole write
\(A_D(s)=(s-a)^v u(s)\) and choose \(h^k=u\).
If \(d=\gcd(k,|v|)\), its local charts are
\[
 s=a+t^{k/d},\qquad
 y=\xi\,t^{v/d}h(a+t^{k/d}),\qquad \xi^k=1.
\]
There are \(d\) distinct charts: changing \(t\) by a \((k/d)\)-th
root identifies exactly \(k/d\) choices of \(\xi\).
They fill all punctured sheets and give a meromorphic \(y\).
When \(v=0\), use \(d=k\) and ordinary unramified charts.
These local changes glue to the given cover on \(D^\circ\).
Above a sufficiently small closed disk there are finitely many compact
coordinate disks.  A finite covering of a compact subset of \(D\)
by such disks proves properness.  This proves (iv).
The polynomial equation also proves (v).

Suppose (iv) holds on a cover \(\pi:X\to D\) of degree \(d\),
with meromorphic function \(F\) and possibly additional branch values
in \(D^\circ\).
Move the base point slightly to avoid them.
Every loop in \(D^\circ\) can be perturbed to avoid the discrete extra
branch values without changing its period.
Lifting such a loop sends the initial sheet to one of at most \(d\) sheets.
The continued value is multiplied by its monodromy factor.
Distinct factors give distinct nonzero values at that point.
Thus \(|\mathcal G_m(D)|\leq d\), proving (i) and the lower bound for cover degree.
The identity \(dF/F=\pi^*\omega_m\), initially valid near the chosen
point, extends as a meromorphic identity on the connected cover.
Thus \(F\) has no zero or pole above \(D^\circ\), and its agreement
with the continued germ persists along the lifted paths.

Finally suppose a polynomial of degree \(d\) vanishes on the germ.
Choose a point in \(D^\circ\) where its leading coefficient is nonzero
and all coefficients are finite.  Such points exist because the leading
coefficient is not identically zero and its exceptional sets are discrete.
Analytic continuation of the polynomial identity shows that all monodromy
values at that point are its roots.
Again these values are distinct, so \(|\mathcal G_m(D)|\leq d\).
This proves (v) implies (i), and every annihilating polynomial has degree
at least \(k\).
The polynomial \(X^k-A_D\) attains that degree and is irreducible:
otherwise one of its smaller-degree factors would vanish on the germ.
A power is single-valued exactly when it annihilates \(\mu_k\), which
requires divisibility by \(k\).  All three minimality claims follow.
\end{proof}

The criterion uses periods of loops in \(D\), not only residues at
points inside \(D\).
For example, choose a small disk about \(a=1/p\) containing no other
point of \(S_m\), where \(p\geq m\geq2\) is prime.
An annulus \(r<|s-a|<R\) inside this disk has no divisor points.
Nevertheless its central loop has residue integral \(2\pi i/p\),
so its minimal cover has degree \(p\).
Deleting the center has removed the local pole from the domain but
has retained the loop around it.

\subsection{Order one, higher orders, and finite restrictions}

For \(m=1\), divisor cancellation gives
\(c_1(1)=1\) and \(c_1(n)=0\) for \(n>1\).
Consequently
\[
 \Delta_1=\zeta^{-1},\qquad
 \omega_1=-\frac{\zeta'}{\zeta}\,ds,\qquad
 \alpha_1(a)=-\ord_a\zeta.
\]
The inverse zeta function is meromorphic on \(\Hp\), and the
scalar continuation of Theorem~\ref{thm:tate} extends it meromorphically
to \(\mathbb C\).
Every residue is an integer and \(\mathcal G_1(D)=\{1\}\) for every connected
\(D\subset\Hp\).  All three degrees in Theorem~\ref{mono:thm:criterion}
are one.  In particular, \(\alpha_1(1/p)=0\) for every prime \(p\).
At \(1\) it has a simple zero, and at a zeta zero of multiplicity \(v\)
it has a pole of order \(v\).

\begin{corollary}[Algebraic obstruction for higher orders]
\label{mono:thm:global}
For \(m\geq2\), the determinant germ is transcendental over
\(\mathcal M(\Hp)\).  No connected finite branched cover of \(\Hp\)
carries its coherent continuation as a meromorphic function, and no
nonzero single-valued meromorphic multiplier changes this conclusion.
\end{corollary}

\begin{proof}
Theorem~\ref{thm:infinite-degree} gives monodromy of order \(p\) at
\(1/p\) for every prime \(p\geq m\).  The image is therefore infinite,
and Theorem~\ref{mono:thm:criterion} excludes both algebraicity and a
finite branched realization.  A meromorphic multiplier adds integer
local orders and integral multiples of \(2\pi i\) to the periods,
so its continuation character is trivial.
\end{proof}

For a reduced rational number \(x\), let \(\operatorname{den}(x)\) be
its positive denominator, with \(\operatorname{den}(0)=1\).
For \(\epsilon>0\), define the finite integer
\[
 M_m(\epsilon)=\operatorname{lcm}
 \{\operatorname{den}(c_m(n)):n\geq1,\ n\epsilon<1\},
\]
where the least common multiple of an empty set is one.

\begin{theorem}[Exact degree on a smaller half-plane]
\label{mono:thm:halfplane}
For continuation along paths in \(\Hp_\epsilon\setminus S_m\),
\[
 \mathcal G_m(\Hp_\epsilon)=\mu_{M_m(\epsilon)}.
\]
Hence the least cover degree and algebraic degree on \(\Hp_\epsilon\)
are exactly \(M_m(\epsilon)\).
Every relatively compact connected domain \(D\subset\Hp\) admits a
finite branched realization of its chosen continued germ.
Its degree divides \(M_m(\epsilon)\) whenever \(D\subset\Hp_\epsilon\).
\end{theorem}

\begin{proof}
Choose \(N\geq m\) with \(1/(N+1)\leq\epsilon\).
For every \(n\leq N\) satisfying \(n\epsilon\geq1\),
\(\zeta(ns)=e^{Z(ns)}\) is holomorphic and nonzero on \(\Hp_\epsilon\).
Its logarithm can be absorbed into the holomorphic remainder.
The only logarithmic factors left have \(n\epsilon<1\).
Every residue is therefore an integer combination of the corresponding
\(c_m(n)\).
The residue-period formula on this convex half-plane puts the monodromy
inside \(\mu_{M_m(\epsilon)}\).
Conversely, the point \(1/n\) lies inside precisely when \(n\epsilon<1\).
Its loop gives \(e^{2\pi i c_m(n)}\).
The subgroup generated by these finitely many roots has order equal to
the least common multiple of their orders: an integer annihilates all
of them precisely when every reduced denominator divides it.
This proves equality.

A compact subset of \(\Hp\) has positive minimum real part.
Thus the closure of a relatively compact \(D\) lies in some
\(\Hp_\epsilon\).  Each loop in \(D\setminus S_m\) is also a loop
in that half-plane, so \(\mathcal G_m(D)\) is a subgroup of the finite cyclic group.
Its order divides \(M_m(\epsilon)\).
Apply Theorem~\ref{mono:thm:criterion}.
\end{proof}

For \(m=2\), the coefficients begin
\[
 c_2(2)=\tfrac12,\quad c_2(3)=\tfrac13,\quad c_2(4)=0,
 \quad c_2(5)=\tfrac15,\quad c_2(6)=-\tfrac16.
\]
Thus the exact half-plane degrees are
\[
 \begin{array}{c|ccccc}
 \epsilon&1/3&1/4&1/5&1/6&1/7\\\hline
 M_2(\epsilon)&2&6&6&30&30
 \end{array}
\]
The strict inequality excludes \(n=3\) at \(\epsilon=1/3\).
The zero coefficient at four contributes no ramification.
At \(1/6\), the exponent \(-1/6\) gives a simple pole after the
sixfold local substitution.
These examples exhibit boundary exclusion, cancellation, and a pole
within the same determinant.

\subsection{Restricting every globally continued branch}

Let \(D\subset\Hp\) be connected, with the path and base point chosen
as in Theorem~\ref{mono:thm:criterion}.
The preimage \(\pi^{-1}(D^\circ)\subset\mathcal E_m\) includes every
branch obtained by paths in the full \(\Omega_m\).
It can have more than one component.  The component containing the
chosen point uses only paths in \(D^\circ\).

\begin{theorem}[Components after restriction]
\label{mono:thm:restriction}
The connected components of \(\pi^{-1}(D^\circ)\) are indexed by the
cosets \(\mathcal G_m/\mathcal G_m(D)\).
Each component is a copy of the minimal pointed continuation cover for
paths in \(D^\circ\), with degree \(|\mathcal G_m(D)|\).
On a component indexed by a representative \(g\in \mathcal G_m\), the determinant
is the chosen component's function multiplied by \(g\).
These identifications extend over the local branched fillings in \(D\).

If \(m\geq2\) and \(D\) is relatively compact, there are infinitely many
components, each of finite degree.
A finite cover carrying a single meromorphic function cannot encode all
globally obtained branch values over \(D\).
Nevertheless the minimal finite cover of \(D\) supports, as separate
functions, all the multiples \(g f\), \(g\in \mathcal G_m\).
\end{theorem}

\begin{proof}
At \(b\), identify the fiber with \(\mathcal G_m\) using the chosen initial path.
Lifting loops in \(D^\circ\) moves a point exactly through its coset
of \(\mathcal G_m(D)\).
Points in the same coset are connected by these lifts.
A path in the preimage connecting two fiber points projects to a loop
in \(D^\circ\), so points in different cosets cannot be connected.
Every component meets that fiber by lifting a path to \(b\).
This proves the component classification.

Prepending a loop with character \(g\) defines a holomorphic map between
the chosen component and its translate.
It has the inverse obtained from \(g^{-1}\), commutes with the projection,
and multiplies the determinant by \(g\).
This proves the stated identification of functions and covers.
At a puncture their exponents are the same \(\alpha_m(a)\).
Multiplication by \(g\) changes only the nonzero constant in
Theorem~\ref{mono:prop:local}, so the maps extend across its local charts.
A local cycle stays within its component, because its loop belongs to
\(\mathcal G_m(D)\).  Filling cannot merge different components.

For relatively compact \(D\), Theorem~\ref{mono:thm:halfplane} makes
\(\mathcal G_m(D)\) finite, whereas Corollary~\ref{mono:thm:global} makes \(\mathcal G_m\) infinite.
The coset set is therefore infinite.
At any ordinary point there are infinitely many distinct global branch
values, all nonzero scalar multiples indexed by \(\mathcal G_m\).
One meromorphic function on a degree-\(d\) cover has at most \(d\)
values above an ordinary point.  It cannot contain that entire set.
On the other hand, each separately specified function \(g f\) is
meromorphic on the same finite cover as \(f\).
This proves both assertions without identifying the two meanings of
carrying all branches.
\end{proof}

For example, on a small disk about \(2\), the group \(\mathcal G_m(D)\) is
trivial and a single local function suffices for paths confined to the disk.
For \(m\geq2\), the global surface restricts to infinitely many disjoint
copies of that disk, with the functions \(g\Delta_m\), \(g\in \mathcal G_m\).
Going around a distant reciprocal-prime point and returning to the disk
changes which copy contains the endpoint.

Theorem~\ref{thm:infinite-degree} also excludes a finite-degree cover
of a neighborhood of zero carrying the specified germs at all sufficiently
small positive reciprocal-prime points.  That conclusion uses their
unbounded local indices.  It gives no natural-boundary statement at any
other point of the imaginary axis.

\subsection{Compatibility with the finite root equations}

The finite degree can also be read directly from the equations
\(y^{M_N}=A_{m,N}(s)\) in Section~\ref{sec:prime-continuation}.
The comparison depends on the initial branch as well as the domain.

\begin{proposition}[The pointed component on a finite half-plane]
\label{prop:finite-root-comparison}
Fix \(m\geq2\), \(N\geq m\), and \(D=U_N=\Hp_{1/(N+1)}\).
Choose \(b=2\) and the initial value \(f_b(2)=\Delta_m(2)\).
Then
\[
 M_m(1/(N+1))=M_N,\qquad A_D=A_{m,N}.
\]
The minimal cover of Theorem~\ref{mono:thm:criterion} identifies with the
finite root surface by \((s,y)\mapsto(s,y)\), including its branched filling.
Inside the global surface, this is the component containing the initial
point.  The components over \(U_N\) are indexed by
\(\mathcal G_m/\mu_{M_N}\).
Changing the initial germ by a factor \(g\in\mathcal G_m\) changes
\(A_D\) by \(g^{M_N}\).
\end{proposition}

\begin{proof}
The strict condition \(n/(N+1)<1\) is exactly \(n\leq N\).
Coefficients below \(m\) vanish, so the two least common multiples agree.
The finite product is the explicitly defined meromorphic function
\[
 A_{m,N}(s)=\exp(M_NT_{m,N}(s))
       \prod_{n=m}^N\zeta(ns)^{-M_Nc_m(n)}.
\]
Both \(A_D\) and \(A_{m,N}\) have logarithmic derivative \(M_N\omega_m\)
on \(U_N\setminus S_m\).
At \(2\), both equal \(\Delta_m(2)^{M_N}\), by the initial extraction
formula.  Their ratio has derivative zero on this connected domain and
value one at \(2\), proving equality there and meromorphically on \(U_N\).
The root equations and initial roots coincide, and their local charts use
the same reduced residue.  The identity map therefore extends over all
filled points.
Theorem~\ref{mono:thm:halfplane} identifies the subgroup with \(\mu_{M_N}\),
and Theorem~\ref{mono:thm:restriction} gives the component set.
Finally the definition of \(A_D\) contains the factor \(f_b(2)^{M_N}\).
Replacing that value by \(g f_b(2)\) multiplies \(A_D\) by \(g^{M_N}\).
Thus paths leaving \(U_N\) can select another component, while paths
confined to \(U_N\) retain the pointed one.
\end{proof}


\section{Simultaneous continuation of several orders}
\label{sec:simultaneous-covers}

Two individually finite covers need not have a common cover whose degree
is their least common multiple.  For the prime determinants of orders two
and three on \(\Hp_{1/5}\), the individual degrees are six and twelve.
The common continuation has degree twenty-four.  The additional sheets
record how the two branches change together, information that their
individual degrees do not retain.

\subsection{The vector of period characters}

Fix positive integers \(\mathbf m=(m_1,\ldots,m_r)\), a real number
\(\epsilon>0\), and \(D=\Hp_\epsilon\).  All functions and covers in
this section are complex analytic.  Their products are ordinary products
of scalar meromorphic functions, and their algebraic degrees refer to the
field \(\mathcal M(D)\).  No cohomological grading is assigned to them.
Put
\[
 J_\epsilon=\{n\in\mathbb Z_{\geq1}:n\epsilon<1\},\qquad
 v_n=(c_{m_1}(n),\ldots,c_{m_r}(n))\in\mathbb Q^r.
\]
Let \(S=\bigcup_i(S_{m_i}\cap D)\) and \(D^\circ=D\setminus S\).
Choose a real base point \(b>\max(1,\epsilon)\), with the positive
values \(\Delta_{m_i}(b)\) given by the initial products.  For a loop
\(\gamma\) in \(D^\circ\) based at \(b\), define its additive period
character by
\[
 \theta(\gamma)=
 \left(\frac1{2\pi i}\int_\gamma\omega_{m_i}\bmod\mathbb Z\right)_{i=1}^r
 \in(\mathbb Q/\mathbb Z)^r.
\]
Rationality follows from the residue formula.  Write
\(G_{\mathbf m}(\epsilon)=\operatorname{im}\theta\).
Its \(i\)-th coordinate, exponentiated by \(x\mapsto e^{2\pi ix}\),
is the character of the \(i\)-th determinant.

\begin{theorem}[Exact simultaneous character]
\label{thm:simultaneous-character}
One has
\[
 G_{\mathbf m}(\epsilon)
 =\left\langle v_n\bmod\mathbb Z^r:n\in J_\epsilon\right\rangle
 \subset(\mathbb Q/\mathbb Z)^r.
\]
This is a finite abelian group, which need not be cyclic.
The connected pointed cover corresponding to \(\ker\theta\) is the
minimal cover on which all the initial germs become single-valued.
It has degree \(|G_{\mathbf m}(\epsilon)|\), carries the functions
\(\mathcal D_{m_i}\), and admits a finite proper branched filling over
\(D\).  These are the least possible degree and the joint algebraic
degree of the germs over \(\mathcal M(D)\).

At \(a\in S\), form the total rational vector
\[
 \alpha(a)=(\alpha_{m_i}(a))_{i=1}^r
 =-\sum_{n\in J_\epsilon}\ord_a\zeta(ns)\,v_n.
\]
Let \(e(a)\) be the least positive integer for which
\(e(a)\alpha(a)\in\mathbb Z^r\).  The filling has
\(|G_{\mathbf m}(\epsilon)|/e(a)\) points above \(a\).
Each has ramification index \(e(a)\), and the \(i\)-th determinant
has order \(e(a)\alpha_{m_i}(a)\) there.
\end{theorem}

\begin{proof}
Choose \(N\geq\max_i m_i\) with \(1/(N+1)\leq\epsilon\).
In the extraction formula, every factor with \(n\epsilon\geq1\)
has its original holomorphic logarithm \(Z(ns)\) on \(D\).
Absorb these terms into the holomorphic functions
\[
 h_i(s)=T_{m_i,N}(s)
       -\sum_{\substack{1\leq n\leq N\\n\epsilon\geq1}}
          c_{m_i}(n)Z(ns).
\]
Terms with zero coefficient are omitted.  Then
\[
 \omega_{m_i}=dh_i-
       \sum_{n\in J_\epsilon}c_{m_i}(n)\,d\log\zeta(ns).
\]
The logarithmic differential is single-valued even where the logarithm
is not.  Summing its orders at each point gives the displayed vector
\(\alpha(a)\), including every coincident zero.  The residue-period
formula on the convex half-plane \(D\) shows that every period is an
integer combination of these vectors.  Thus the image is contained in
the group generated by the \(v_n\).

For the reverse inclusion, \(1/n\in D\) exactly when
\(n\epsilon<1\).  At \(1/n\), only \(\zeta(ns)\) has a pole.
Every other scaled argument is positive real and different from one,
so its zeta value is finite and nonzero.  A small positive loop about
\(1/n\), connected to \(b\) by any path in \(D^\circ\), therefore
has period vector \(v_n\bmod\mathbb Z^r\).  Zero vectors need no
loop.  This proves equality without any assumption on complex zeta zeros.
There are finitely many rational generators, so the group is finite.

Use paths from \(b\), identifying two with the same endpoint exactly
when the integrals of every \(\omega_{m_i}\) differ by
\(2\pi i\mathbb Z\).  The disk charts of
Theorem~\ref{thm:minimal-cover} give a connected cover, and on it put
\[
 \mathcal D_{m_i}[\gamma]
 =\Delta_{m_i}(b)\exp\int_\gamma\omega_{m_i}.
\]
The fibre is a torsor for the vector image.  A pointed cover carrying
all these functions maps to this one by projected paths; the functions
make the map well-defined, and disk charts prove holomorphy and uniqueness.
Distinct period vectors give distinct tuples of nonzero values above
\(b\).  Any finite realization of all germs must therefore have at
least \(|G_{\mathbf m}(\epsilon)|\) sheets.

Near \(a\), write
\(\omega_{m_i}=\alpha_{m_i}(a)\,ds/(s-a)+dq_i\), with \(q_i\)
holomorphic on a disk.  The local sheets form cycles of length \(e(a)\).
Fill each cycle by
\[
 s=a+t^{e(a)},\qquad
 \mathcal D_{m_i}=C_i t^{e(a)\alpha_{m_i}(a)}
                    e^{q_i(a+t^{e(a)})},\qquad C_i\ne0.
\]
The constants are fixed by the continued germs.  These charts give all
points and orders stated in the theorem.  Changing a root coordinate
multiplies \(t\) by a root of unity, so the charts glue to the
punctured cover.  They also prove properness: the inverse image of a small
closed disk is a finite union of compact disks, and finitely many such
base disks cover any compact subset of \(D\).
Any meromorphic local realization requires all its pulled-back orders
to be integers, hence a local index divisible by \(e(a)\).

The functions on this finite cover are algebraic over \(\mathcal M(D)\):
the product of \(X-F\) over its sheets has single-valued meromorphic
coefficients for each meromorphic \(F\).  Every finitely generated subextension of its meromorphic function
field is finite algebraic.  In characteristic zero it has a primitive
element: for two generators, all but finitely many complex linear
combinations separate their finitely many embeddings, and induction
handles any finite list.  The sheet polynomial bounds the degree of
that primitive element by the cover degree.  Thus the whole function
field has degree at most the cover degree: any larger linearly
independent finite list would lie in such a subextension.
For a direct verification of equality for the field generated by our
coordinates, choose complex constants \(\lambda_i\) so that the values
of \(F=\sum_i\lambda_i\mathcal D_{m_i}\) are distinct on the fibre
over \(b\).  This excludes finitely many proper linear hyperplanes,
since the tuples on that fibre are distinct.  The continued germ of
\(F\) has exactly \(|G_{\mathbf m}(\epsilon)|\) distinct values there.
Every polynomial annihilating it has at least that degree, as in the
proof of Theorem~\ref{mono:thm:criterion}.  Its sheet product attains
that degree.  The coordinate field contains \(F\) and lies in the
function field of the cover, proving the asserted joint degree.
\end{proof}

The strict cutoff in this proof is essential.  If \(n\epsilon=1\),
then \(\operatorname{Re}(ns)>1\) everywhere in \(D\); its pole is on
the excluded boundary, and its logarithm contributes no period in \(D\).
At a common complex zero, the sum defining \(\alpha(a)\) must be taken
before denominators are reduced.  Thus coincident divisors can reduce
local ramification, while the separate reciprocal-real poles already
force the full group in the theorem.

\subsection{Integer lattices and coordinates}

The theorem turns the degree into a finite integer calculation.
Choose a positive integer \(Q\) divisible by the denominators of every
coordinate of the finitely many \(v_n\).  Let \(A\) be the integer
matrix with columns \(Qv_n\), and define the full-rank lattice
\[
 \Lambda=Q\mathbb Z^r+A\mathbb Z^{J_\epsilon}\subset\mathbb Z^r.
\]
Here \(\mathbb Z^{J_\epsilon}\) is the free group with one generator
per index in \(J_\epsilon\).  Reduction modulo \(Q\mathbb Z^r\)
gives
\[
 G_{\mathbf m}(\epsilon)\cong\Lambda/Q\mathbb Z^r,
 \qquad
 |G_{\mathbf m}(\epsilon)|=\frac{Q^r}{[\mathbb Z^r:\Lambda]}.
\]
One exact procedure is integer column reduction of \([QI_r\ A]\)
to a basis matrix \(B\) for \(\Lambda\).  Then
\(C=B^{-1}QI_r\) is integral, and
\[
 G_{\mathbf m}(\epsilon)\cong\mathbb Z^r/C\mathbb Z^r.
\]
Apply integer row and column operations of determinant \(\pm1\) to
\(C\), obtaining its Smith diagonal
\(d_1,\ldots,d_r\), with \(d_i>0\) and \(d_i\mid d_{i+1}\).
The group is \(\bigoplus_i\mathbb Z/d_i\mathbb Z\), with factors
of order one omitted, and its degree is \(\prod_i d_i\).

These reductions use only the Euclidean algorithm.  To recall why they
terminate, move a nonzero entry to the first position and use division
with remainder to reduce its row and column until it divides them.
If it fails to divide an entry in the remaining block, adding that
entry's row to the first row and repeating produces a smaller positive
pivot.  Eventually the pivot divides every remaining entry.  Clear its
row and column and repeat in the smaller block.  The resulting diagonal
has the stated divisibility.  Every operation changes a basis in the
source or target free group, so it preserves the quotient.  Applied to
the rectangular generating matrix, the same operations produce a lattice
basis.  This proves the procedure and its degree formula, including
redundant or zero generators.

For explicit meromorphic coordinates put
\[
 k_i=M_{m_i}(\epsilon),\qquad
 A_i(s)=e^{k_i h_i(s)}
        \prod_{n\in J_\epsilon}\zeta(ns)^{-k_i c_{m_i}(n)}.
\]
The powers are integers.  Over \(D^\circ\), form the product of the
root equations
\[
 y_i^{k_i}=A_i(s)\quad(1\leq i\leq r),\qquad
 y_i(b)=\Delta_{m_i}(b).
\]
Coordinatewise multiplication by roots of unity gives an action of
\(G_0=\prod_i(\tfrac1{k_i}\mathbb Z/\mathbb Z)\) on the product,
simply transitive on every fibre. Continuation acts by translations
in the subgroup \(G_{\mathbf m}(\epsilon)\).
Consequently the product has
\(\prod_i k_i/|G_{\mathbf m}(\epsilon)|\) connected components.
The component containing the displayed base point is the minimal
simultaneous cover.  Closing the equations in
\(D\times(\mathbb P^1)^r\) and replacing each local branch by the
smooth coordinates of Theorem~\ref{thm:simultaneous-character} gives
its normalization and branched filling.

The selected component also has finite equations beyond the individual
root equations.  A character of \(G_0\) is represented by integers
\(a_1,\ldots,a_r\), with \(a_i\) taken modulo \(k_i\), and has
value \(\sum_i a_i g_i\bmod\mathbb Z\) at \(g\in G_0\).
It vanishes on the simultaneous subgroup exactly when
\[
 d_n(a):=\sum_i a_i c_{m_i}(n)\in\mathbb Z
 \quad\text{for every }n\in J_\epsilon.
\]
For each such character define
\[
 R_a(s)=\exp\left(\sum_i a_i h_i(s)\right)
          \prod_{n\in J_\epsilon}\zeta(ns)^{-d_n(a)}.
\]
It is meromorphic on \(D\).  The equations
\(\prod_i y_i^{a_i}=R_a(s)\) select the pointed component.
Indeed, on any disk the root tuples differ by \(G_0\), and these
equations retain exactly the common kernel of the characters vanishing
on \(G_{\mathbf m}(\epsilon)\).  That kernel is the subgroup itself:
Smith reduction of the quotient \(G_0/G_{\mathbf m}(\epsilon)\)
provides a character detecting every nonzero quotient element.
A finite generating set of its character group suffices.
Changing \(a_i\) by \(k_i\) multiplies both sides by \(A_i\), so
this description is independent of the chosen representatives.

Coordinate projection maps the minimal simultaneous cover onto the
minimal \(m_i\)-cover, with degree \(|G_{\mathbf m}(\epsilon)|/k_i\).
More generally, deleting coordinates gives the cover for their projected
period group.  In the local coordinates at a divisor, such a map is
\(t\mapsto t^{e/e'}\), up to the root of unity fixed by the base
point, where \(e'\) is the local period order after projection.
This also proves extension over every filled point.

The order-change multiplier is the meromorphic function \(y_j/y_i\).
Its period character is \(g\mapsto g_j-g_i\), its logarithmic differential is
\(\omega_{m_j}-\omega_{m_i}\), and its local order on the simultaneous
cover is \(e(a)(\alpha_{m_j}(a)-\alpha_{m_i}(a))\).
The ratios satisfy the cocycle identity by cancellation.
Its minimal cover is the quotient by the kernel of this difference
character.  A map from the individual \(i\)-cover to the individual
\(j\)-cover preserving continued paths exists exactly when
\(\ker(g\mapsto g_i)\subseteq\ker(g\mapsto g_j)\).
Thus the scalar descent condition is retained, while the common cover
is governed by the whole vector group.

If \(\epsilon'>\epsilon\), then \(J_{\epsilon'}\subseteq J_\epsilon\)
and the smaller half-plane has the subgroup generated by the smaller
list.  Restricting the larger continuation surface gives components
indexed by \(G_{\mathbf m}(\epsilon)/G_{\mathbf m}(\epsilon')\).
The positive initial branches select one component and its compatible
inclusion.  For a connected domain with holes, the same statement uses
its actual period subgroup; the finite list of divisor points lying
inside the domain alone does not determine it.

\subsection{Orders two and three on \texorpdfstring{\(\Hp_{1/5}\)}{Re(s) > 1/5}}
\label{subsec:simultaneous-two-three}

Here \(J_{1/5}=\{1,2,3,4\}\).  The nonzero period generators are
\[
 g_2=(\tfrac12,0),\qquad
 g_3=(\tfrac13,\tfrac13),\qquad
 g_4=(0,\tfrac14).
\]
They have orders two, three and four.  If
\(a g_2+b g_3+c g_4=0\), with \(0\leq a<2\), \(0\leq b<3\),
\(0\leq c<4\), the first coordinate gives
\(3a+2b\equiv0\pmod6\).  Reduction modulo two and three forces
\(a=b=0\); the second coordinate then forces \(c=0\).
Therefore
\[
 G_{(2,3)}(1/5)
 \cong\mathbb Z/2\mathbb Z\oplus\mathbb Z/3\mathbb Z
                    \oplus\mathbb Z/4\mathbb Z
 \cong\mathbb Z/2\mathbb Z\oplus\mathbb Z/12\mathbb Z.
\]
The first coordinate image has order six and the second has order
twelve.  Their least common multiple, twelve, is strictly smaller than
the simultaneous degree twenty-four.

The lattice calculation gives the same invariants by a different route.
With \(Q=12\), the generator columns are \((6,0),(4,4),(0,3)\).
Together with \(12\mathbb Z^2\), they span
\[
 \Lambda=\{(x,y)\in\mathbb Z^2:x\equiv0\pmod2,
                                      \ x-y\equiv0\pmod3\}.
\]
A basis and its relation matrix are
\[
 B=\begin{pmatrix}6&4\\0&1\end{pmatrix},\qquad
 B^{-1}(12I_2)=\begin{pmatrix}2&-8\\0&12\end{pmatrix}.
\]
The greatest common divisor of the latter entries is two and its
determinant is twenty-four.  Its Smith invariants are consequently two
and twelve.  The subgroup consists of the twenty-four pairs
\((a/6,b/12)\) with \(b-2a\equiv0\pmod3\).

To see the functions, abbreviate \(T_2=T_{2,4}\), \(T_3=T_{3,4}\),
and \(z_n(s)=\zeta(ns)\).  On the simultaneous cover there are roots
\[
 u^2=z_2^{-1},\qquad v^3=z_3^{-1},\qquad w^4=z_4^{-1},
\]
chosen positive at \(s=2\).  Their product cover is connected: the
loops about \(1/2,1/3,1/4\) change precisely one of these three
roots and generate all \(2\cdot3\cdot4\) sheets.  The determinant
coordinates are
\[
 y_2=e^{T_2}uv,\qquad y_3=e^{T_3}vw.
\]
No two root triples have the same pair.  For a ratio of two triples,
\(uv=vw=1\) and \(u^2=v^3=w^4=1\) force \(v=u=w=1\).
One can also recover the functions directly: writing
\(q_2=e^{-T_2}y_2\), \(q_3=e^{-T_3}y_3\), one has
\[
 v=\frac1{z_2z_3q_2^2},\qquad u=q_2/v,\qquad w=q_3/v.
\]
These identities are meromorphic and continue across the filled points.
Thus the root triples and the determinant pair describe the same
normalized connected surface, including its functions.

The individual equations are
\[
 y_2^6=A:=e^{6T_2}z_2^{-3}z_3^{-2},\qquad
 y_3^{12}=B_3:=e^{12T_3}z_3^{-4}z_4^{-3}.
\]
Their fibre product has seventy-two sheets before a component is
selected.  Put
\[
 C=e^{4(T_3-T_2)}z_2^2z_4^{-1}.
\]
Then \(C^3=B_3/A^2\), and the three components are distinguished by
\[
 y_3^4=\xi C y_2^4,\qquad \xi^3=1.
\]
The positive values at \(2\) select \(\xi=1\).  On that component
\(y_2\) and \(y_3\) have the specified germs; its projections have
degrees four and two.  The normalization is therefore essential to the
geometric comparison, beyond the separate degrees of the equations.

The local data at the reciprocal-real poles are
\[
\begin{array}{c|c|c|c|c}
 a&\alpha(a)&e(a)&\text{points above }a&
        (\ord y_2,\ord y_3)\\\hline
 1/2&(1/2,0)&2&12&(1,0)\\
 1/3&(1/3,1/3)&3&8&(1,1)\\
 1/4&(0,1/4)&4&6&(0,1)
\end{array}
\]
At any coincident complex zeros, the summed vector in
Theorem~\ref{thm:simultaneous-character} gives the corresponding
ramification and orders without a simplicity assumption.

Finally \(H_{2,3}=y_3/y_2\) has character values generated by
\(-1/2\) and \(1/4\), hence a four-sheeted minimal cover and a kernel
of order six on the simultaneous cover.  Its fourth power is \(C\).
The loop at \(1/4\) is trivial on \(y_2\) and nontrivial on \(y_3\),
whereas the loop at \(1/2\) is trivial on \(y_3\) and nontrivial on
\(y_2\).  Neither individual minimal cover maps to the other while
preserving continued paths.  The simultaneous cover supplies both
projections and their multiplier with the least possible degree.
This comparison concerns scalar determinant germs.  A map to an
arithmetic complex would additionally have to preserve its differential,
products and arithmetic classes; none is supplied by these root equations.

\section{The archimedean oscillator determinant}
\label{sec:oscillator}

\subsection{An equally spaced spectrum}

The Gaussian Mellin integral defines the real archimedean factor.
An additive operator model for its reciprocal must account for an infinite
sequence of equally spaced zeros. The diagonal oscillator is the simplest
such model. Its eigenvalue product diverges, so the normalization must come
from a convergent spectral sum and its analytic continuation.

Let \(H=\ell^2(\mathbb N_0;\mathbb C)\), where
\(\mathbb N_0=\{0,1,2,\ldots\}\). Its inner product is
\(\langle x,y\rangle=\sum_{r\geq0}x_r\overline{y_r}\), linear in
the first variable, and \(e_r\) is the coordinate basis vector.
For positive real \(b\), define
\[
 B_b e_r=br e_r,\qquad
 \mathcal D(B_b)=
 \left\{x\in H:\sum_{r\geq0}b^2r^2|x_r|^2<\infty\right\}.
\]
For positive real \(c\) and \(\operatorname{Re}s>0\), put
\[
 T_{b,c}(s):\mathcal D(B_b)\longrightarrow H,
 \qquad T_{b,c}(s)=\frac{sI+B_b}{c}.
\]
The domain is independent of \(s\). This is a complex operator pencil.
For real \(s>0\), it is positive and self-adjoint.

To check the domain claims, finite sequences give a dense domain.
Coordinate convergence shows that a limit of graph pairs \((x,B_bx)\)
is again a graph pair, so \(B_b\) is closed.
If \(y\) is in the adjoint domain, testing against \(e_r\) forces
\((B_b^*y)_r=br y_r\). This sequence belongs to \(H\) exactly when
\(y\in\mathcal D(B_b)\). The converse follows from the inner product
formula, proving self-adjointness.
For complex \(s\), the same coordinate computation proves normality of
\(T_{b,c}(s)\). Its inverse has diagonal entries \(c/(s+br)\),
which tend to zero. Truncation in the coordinate basis approximates this
inverse in operator norm by finite-rank operators, so it is compact.
Its range lies in \(\mathcal D(B_b)\), since \(br/(s+br)\) is bounded.
The eigenvalues of the pencil are exactly \((s+br)/c\), each of multiplicity
one. Indeed, away from this discrete set the diagonal inverse of the
shifted pencil is bounded and has range in the same domain.

Every eigenvalue lies in the right half-plane. We use its principal
logarithm, with argument in \((-\pi,\pi)\), to define
\begin{equation}
 Z(w,s)=\sum_{r=0}^{\infty}
 \exp\left(-w\operatorname{Log}\frac{s+br}{c}\right),
 \qquad \operatorname{Re}w>1.                     \label{eq:osc-spectral-sum}
\end{equation}
For real \(s>0\), the logarithms of the eigenvalues tend to \(+\infty\),
so their ordinary product diverges. If \(Z\) continues holomorphically
through \(w=0\), its zeta-regularized determinant is, by definition,
\begin{equation}
 D_{b,c}(s)=\det_\infty T_{b,c}(s)
 :=\exp\bigl(-\partial_w Z(0,s)\bigr).             \label{eq:osc-det-definition}
\end{equation}
For the real Hilbert space \(\ell^2(\mathbb N_0;\mathbb R)\), this
definition uses complexification. Each real eigenline becomes one complex
eigenline, so its multiplicity remains one.

\subsection{Continuation and the gamma normalization}

Write \(U=\{a\in\mathbb C:\operatorname{Re}a>0\}\).
On this half-plane, \(\operatorname{Log}\) denotes the principal logarithm,
which is real on the positive axis. All powers of \(a+x\), for \(x\geq0\),
use this logarithm. For \(\operatorname{Re}w>1\), define
\[
 \zeta_H(w,a)=\sum_{r=0}^{\infty}(a+r)^{-w},\qquad a\in U.
\]
On every compact subset of this domain, the summands are bounded by
\(C(1+r)^{-1-\varepsilon}\) for some \(C,\varepsilon>0\).
Thus the sum is jointly holomorphic in \((w,a)\).

\begin{lemma}[A summation formula with a controlled remainder]
\label{lem:osc-summation}
For real \(t\), let \(\{t\}=t-\lfloor t\rfloor\) and set
\(q(t)=\{t\}(1-\{t\})/2\). For every integer \(N\geq0\),
\begin{align}
 \zeta_H(w,a)
 &=\sum_{r=0}^{N-1}(a+r)^{-w}
   +\frac{(N+a)^{1-w}}{w-1}+\frac12(N+a)^{-w}
   +R_N(w,a),                                      \label{eq:osc-cutoff}\\
 R_N(w,a)
 &=w(w+1)\int_N^\infty q(t)(t+a)^{-w-2}\,dt.       \label{eq:osc-remainder}
\end{align}
The integral is jointly holomorphic for \(\operatorname{Re}w>-1\),
\(a\in U\). Formula~\eqref{eq:osc-cutoff} gives a meromorphic continuation
there, with only a simple pole at \(w=1\), of residue \(1\).
In particular,
\begin{equation}
 \zeta_H(0,a)=\frac12-a,\qquad a\in U.              \label{eq:osc-zeta-zero}
\end{equation}
\end{lemma}

\begin{proof}
On \([n,n+1]\), the polynomial \(q(t)=(t-n)(n+1-t)/2\)
vanishes at both endpoints. Its endpoint derivatives are \(1/2\) and
\(-1/2\), and its second derivative is \(-1\).
For any twice continuously differentiable complex function \(f\),
two integrations by parts give
\[
 \int_n^{n+1}q(t)f''(t)\,dt
 =\frac{f(n)+f(n+1)}2-\int_n^{n+1}f(t)\,dt.
\]
Sum this identity from \(n=N\) to \(n=M-1\). The endpoints give
\[
 \sum_{r=N}^{M-1}f(r)
 =\int_N^M f(t)\,dt+\frac{f(N)-f(M)}2
   +\int_N^M q(t)f''(t)\,dt.
\]
Take \(f(t)=(t+a)^{-w}\) and first assume \(\operatorname{Re}w>1\).
Then \(f''(t)=w(w+1)(t+a)^{-w-2}\).
The terms at \(M\) vanish as \(M\to\infty\), while
\(\int_N^\infty f(t)\,dt=(N+a)^{1-w}/(w-1)\).
This proves~\eqref{eq:osc-cutoff} in the initial domain.

For joint continuation, fix compact sets \(K\subset U\) and
\(W\subset\{\operatorname{Re}w>-1\}\).
Put \(\delta=\min_K\operatorname{Re}a>0\) and
\(\sigma=\min_W\operatorname{Re}w>-1\).
The bounds \(0\leq q\leq1/8\) and
\(t+\delta\leq|t+a|\leq t+\max_K|a|\) imply
\[
 |q(t)(t+a)^{-w-2}|
 \leq C_{K,W}(1+t)^{-\sigma-2}.
\]
Indeed, \(|\arg(t+a)|<\pi/2\), and \(\operatorname{Im}w\) is bounded
on \(W\). The same bound for any fixed parameter derivatives acquires
only powers of \(\log(2+t)\) and additional nonpositive powers of \(1+t\).
These bounds remain integrable. Differentiation under the integral is
therefore valid locally in both variables. They also give
\[
 |R_N(w,a)|\leq C_{K,W}|w(w+1)|(1+N)^{-\sigma-1}.
\]
All other terms in~\eqref{eq:osc-cutoff} are holomorphic except the displayed
pole at \(w=1\). Its residue is \((N+a)^0=1\).
The continued formulas for different \(N\) agree, since they agree for
\(\operatorname{Re}w>1\). At zero, the remainder vanishes, and the
remaining terms give \(N-(N+a)+1/2\).
\end{proof}

The derivative at zero involves logarithms of finite products.
The normalization is that of the Gamma integral already used in the
local real factor:
\begin{equation}
 \Gamma(a)=\int_0^\infty e^{-x}x^{a-1}\,dx,
 \qquad a\in U,                                   \label{eq:osc-gamma-integral}
\end{equation}
where \(x^{a-1}=\exp((a-1)\log x)\) uses the real logarithm.

\begin{lemma}[The gamma logarithm and its reciprocal]
\label{lem:osc-gamma}
The integral~\eqref{eq:osc-gamma-integral} is holomorphic and nonzero on
\(U\). There is a unique holomorphic logarithm \(L\) of \(\Gamma\)
which is real on the positive axis. It satisfies, locally uniformly on \(U\),
\begin{equation}
 L(a)=\lim_{N\to\infty}
 \left\{\log((N-1)!)+a\log N
              -\sum_{r=0}^{N-1}\operatorname{Log}(a+r)\right\}.
 \label{eq:osc-gamma-log-limit}
\end{equation}
The reciprocal of \(\Gamma\) extends to the entire function
\begin{equation}
 E(a)=a e^{\gamma a}\prod_{n=1}^{\infty}(1+a/n)e^{-a/n},
 \qquad
 \gamma=\lim_{N\to\infty}\left(\sum_{n=1}^N\frac1n-\log N\right).
 \label{eq:osc-gamma-product}
\end{equation}
Its zeros are simple and occur exactly at \(0,-1,-2,\ldots\).
Moreover, \(\Gamma(1)=1\) and \(\Gamma(a+1)=a\Gamma(a)\) on \(U\).
\end{lemma}

\begin{proof}
On a compact set with \(0<\delta\leq\operatorname{Re}a\leq A\),
the integral and each fixed derivative have an integrable majorant
\[
 e^{-x}(x^{\delta-1}+x^{A-1})(1+|\log x|^k).
\]
This proves holomorphy. Direct integration gives \(\Gamma(1)=1\).
Integration by parts, with vanishing endpoints, gives the recurrence.

For \(N\geq1\), repeated integration by parts on \([0,1]\) gives
\[
 \int_0^1 t^{a-1}(1-t)^N\,dt
 =\frac{N!}{a(a+1)\cdots(a+N)}.
\]
At each step, integrate \(t^{a+j-1}\) and differentiate the remaining
integer power of \(1-t\). The last integral is \(1/(a+N)\).
The substitution \(x=Nt\) yields
\[
 \frac{N!N^a}{a(a+1)\cdots(a+N)}
 =\int_0^N x^{a-1}(1-x/N)^N\,dx
 \longrightarrow\Gamma(a).
\]
For \(0\leq x<N\), one has \((1-x/N)^N\leq e^{-x}\).
Extend the integrand by zero for \(x\geq N\).
The preceding majorants prove convergence locally uniformly in \(a\),
including any fixed number of derivatives.

Let \(H_N=\sum_{n=1}^N1/n\).
The inequality \(\log(1+1/N)>1/(N+1)\) shows that \(H_N-\log N\)
decreases. Integral comparison gives \(H_N\geq\log(N+1)\), so the
sequence has a finite limit \(\gamma\).
The reciprocals of the finite expressions above are
\[
 aN^{-a}\prod_{n=1}^N(1+a/n)
 =a\exp\bigl(a(H_N-\log N)\bigr)
       \prod_{n=1}^N(1+a/n)e^{-a/n}.
\]
For \(|u|\leq1/2\), the logarithm series gives
\( |\operatorname{Log}(1+u)-u|\leq |u|^2\).
Thus the logarithms of the product tails converge uniformly on each
compact subset of \(\mathbb C\). Each tail is the exponential of a
holomorphic sum and is nonzero. The finite initial factors prove that
\(E\) is entire and identify every zero and its order.
Multiplying the two finite expressions and passing to their locally
uniform limits gives \(\Gamma E=1\) on \(U\).

On \(U\), define explicitly
\[
 L(a)=-\operatorname{Log}a-\gamma a
       -\sum_{n=1}^{\infty}
        \bigl(\operatorname{Log}(1+a/n)-a/n\bigr).
\]
The locally uniform convergence proves holomorphy, and
\(\exp L=1/E=\Gamma\). Every term is real for \(a>0\).
Two holomorphic logarithms differ by a constant in \(2\pi i\mathbb Z\),
so the stated real condition proves uniqueness.
Finally, the expression in braces in~\eqref{eq:osc-gamma-log-limit} equals
\[
 -\operatorname{Log}a+a(\log N-H_{N-1})
 -\sum_{n=1}^{N-1}
        \bigl(\operatorname{Log}(1+a/n)-a/n\bigr).
\]
Here \(H_{N-1}-\log N\to\gamma\), which proves the limit with its
specified branch.
\end{proof}

\begin{lemma}[The factorial constant]
\label{lem:osc-stirling}
As positive integers \(N\) tend to infinity,
\begin{equation}
 \log((N-1)!)
 =\left(N-\frac12\right)\log N-N
      +\frac12\log(2\pi)+o(1).                    \label{eq:osc-stirling}
\end{equation}
Also \(\Gamma(1/2)=\sqrt\pi\).
\end{lemma}

\begin{proof}
Let \(J=\int_{\mathbb R}e^{-u^2/2}\,du\), which is finite and positive.
The product integral in polar coordinates gives
\[
 J^2=\int_{\mathbb R^2}e^{-(u^2+v^2)/2}\,du\,dv
 =2\pi\int_0^\infty e^{-r^2/2}r\,dr=2\pi.
\]
Positivity gives \(J=\sqrt{2\pi}\).
Substitution \(x=u^2\) in~\eqref{eq:osc-gamma-integral}, followed by
rescaling this Gaussian integral, gives \(\Gamma(1/2)=\sqrt\pi\).

The recurrence for \(\Gamma\) gives \(n!=\int_0^\infty x^n e^{-x}\,dx\).
Set \(x=n+\sqrt n\,u\). Then
\begin{equation}
 \frac{n!}{n^n e^{-n}\sqrt n}
 =\int_{-\sqrt n}^{\infty}
   \exp\left(n\left[\log(1+u/\sqrt n)-u/\sqrt n\right]\right)du.
 \label{eq:osc-laplace}
\end{equation}
For each fixed \(u\), the exponent tends to \(-u^2/2\).
We give bounds that justify integrating this limit.
For \(-1<x\leq0\), differentiation shows
\(\log(1+x)-x\leq-x^2/2\).
For \(x\geq0\), integration of \(t/(1+t)\) gives
\[
 x-\log(1+x)=\int_0^x\frac{t}{1+t}\,dt
 \geq\frac{x^2}{2(1+x)}.
\]
Consequently the integrand in~\eqref{eq:osc-laplace} is bounded by
\(e^{-u^2/2}\) on \((-\sqrt n,0]\), and by \(e^{-u^2/4}\) on
\([0,\sqrt n]\). Extend these portions by zero to apply dominated
convergence on the two half-lines.

For the remaining tail, \(\log(1+x)/x\) decreases for \(x>0\).
Its derivative has the sign of
\(x/(1+x)-\log(1+x)<0\), as integration of \(1/(1+t)\) shows.
Thus, for \(x\geq1\),
\(\log(1+x)-x\leq-\kappa x\), where \(\kappa=1-\log2>0\).
The tail in~\eqref{eq:osc-laplace} is at most
\[
 \int_{\sqrt n}^{\infty}e^{-\kappa\sqrt n\,u}\,du
 =\frac{e^{-\kappa n}}{\kappa\sqrt n}\longrightarrow0.
\]
It follows that \(n!/(n^n e^{-n}\sqrt n)\to\sqrt{2\pi}\).
Take real logarithms, then use \((N-1)!=N!/N\), to obtain
\eqref{eq:osc-stirling}.
\end{proof}

\begin{proposition}[The two Hurwitz values]
\label{prop:osc-hurwitz-values}
The continuation in Lemma~\ref{lem:osc-summation} is holomorphic near
\(\{0\}\times U\), and
\begin{equation}
 \zeta_H(0,a)=\frac12-a,\qquad
 \partial_w\zeta_H(0,a)=L(a)-\frac12\log(2\pi).
 \label{eq:osc-hurwitz-values}
\end{equation}
Here \(L=\log\Gamma\) is exactly the holomorphic branch constructed in
Lemma~\ref{lem:osc-gamma}.
\end{proposition}

\begin{proof}
The first identity and joint holomorphy were proved in
Lemma~\ref{lem:osc-summation}. Differentiating its exact formula at zero
gives, with \(A=N+a\),
\begin{align*}
 \partial_w\zeta_H(0,a)
 ={}&-\sum_{r=0}^{N-1}\operatorname{Log}(a+r)
       +(A-\tfrac12)\operatorname{Log}A-A\\
    &+\int_N^\infty\frac{q(t)}{(t+a)^2}\,dt.
\end{align*}
The last integral has absolute value at most \(1/(8(N+\delta))\)
when \(\operatorname{Re}a\geq\delta>0\).
It therefore tends to zero locally uniformly on \(U\).
Use~\eqref{eq:osc-gamma-log-limit} to rewrite the remaining expression
as the convergent approximation to \(L(a)\), plus
\[
 -\log((N-1)!)-a\log N
 +(N+a-\tfrac12)\operatorname{Log}(N+a)-(N+a).
\]
By~\eqref{eq:osc-stirling}, this equals
\[
 -\frac12\log(2\pi)
 +(N+a-\tfrac12)\operatorname{Log}(1+a/N)-a+o(1).
\]
For \(a\) in a fixed compact set,
\(\operatorname{Log}(1+a/N)=a/N+O(N^{-2})\) uniformly.
The last two terms therefore tend to zero locally uniformly.
This proves the derivative identity for complex \(a\in U\) directly.
\end{proof}


\subsection{The determinant and the zero mode}

\begin{theorem}[The normalized oscillator determinant]
\label{thm:osc-det}\label{prop:oscillator-det}
For positive real \(b,c\), the spectral sum~\eqref{eq:osc-spectral-sum}
has a jointly meromorphic continuation to
\(\{\operatorname{Re}w>-1\}\times\{\operatorname{Re}s>0\}\).
Its only pole in \(w\) is at \(1\), with residue \(c/b\).
It is jointly holomorphic near \(w=0\), and
\begin{equation}
 D_{b,c}(s)
 =\left(\frac bc\right)^{1/2-s/b}
       \frac{\sqrt{2\pi}}{\Gamma(s/b)}.
 \label{eq:osc-det-formula}
\end{equation}
The power uses the real logarithm of \(b/c\).
This function extends to an entire function of \(s\), with simple zeros
exactly at \(s=-br\), \(r\in\mathbb N_0\).
\end{theorem}

\begin{proof}
Set \(a=s/b\) and \(k=b/c>0\).
Since positive scaling preserves the principal argument,
\(\operatorname{Log}(k(a+r))=\log k+\operatorname{Log}(a+r)\).
Absolute local uniform convergence gives
\[
 Z(w,s)=k^{-w}\zeta_H(w,a).
\]
Lemma~\ref{lem:osc-summation} proves the joint continuation and its
residue. Proposition~\ref{prop:osc-hurwitz-values} gives
\[
 -\partial_w Z(0,s)
 =\left(\frac12-a\right)\log k-L(a)+\frac12\log(2\pi).
\]
Exponentiation proves~\eqref{eq:osc-det-formula} on \(\operatorname{Re}s>0\).
For all complex \(s\), define its continuation by
\[
 \sqrt{2\pi}\exp\bigl((\tfrac12-s/b)\log k\bigr)E(s/b).
\]
The exponential is entire and nonzero. Lemma~\ref{lem:osc-gamma}
identifies the simple zeros of \(E(s/b)\).
\end{proof}

The continuation in this theorem is a scalar entire continuation.
At \(s=-br\), the operator has a zero eigenvalue.
Across a negative real eigenvalue, the principal logarithm has a cut,
so the theorem does not assert a holomorphic principal-cut spectral family
on the whole \(s\)-plane. The specified right half-plane fixes the entire
scalar function uniquely.

\begin{corollary}[Scale and removal of the lowest mode]
For \(d>0\), scaling the operator gives
\begin{equation}
 \det_\infty(dT_{b,c}(s))
 =d^{1/2-s/b}D_{b,c}(s).                           \label{eq:osc-scale}
\end{equation}
The determinant on the closed span of \(e_1,e_2,\ldots\) is
\(D_{b,c}(s+b)\), and
\begin{equation}
 D_{b,c}(s)=\frac{s}{c}D_{b,c}(s+b),\qquad
 D'_{b,c}(0)=\sqrt{\frac{2\pi}{bc}}.               \label{eq:osc-mode}
\end{equation}
The continued determinant with the lowest mode removed has value
\(D_{b,c}(b)=\sqrt{2\pi c/b}\) at \(s=0\).
\end{corollary}

\begin{proof}
Positive scaling replaces \(Z\) by \(d^{-w}Z\).
Differentiating at zero gives~\eqref{eq:osc-scale}, since
\(Z(0,s)=1/2-s/b\).
Reindexing the tail identifies its spectral zeta function with the one
at \(s+b\). The omitted eigenvalue contributes \((s/c)^{-w}\), whose
derivative at zero is \(-\operatorname{Log}(s/c)\).
This proves the first identity in~\eqref{eq:osc-mode} initially in the
right half-plane and then everywhere by entire continuation.
Since \(E(a)/a\to1\) at zero by~\eqref{eq:osc-gamma-product},
\eqref{eq:osc-det-formula} gives the stated slope.
Setting \(s=b\) in that formula gives the final value.
\end{proof}

More generally, recurrence and \(E(a)/a\to1\) give
\[
 D'_{b,c}(-br)=\frac{\sqrt{2\pi}}{b}
       \left(\frac bc\right)^{r+1/2}(-1)^r r!.
\]
Indeed, \(E(a)=a(a+1)\cdots(a+r)E(a+r+1)\), and \(E(1)=1\).
This also verifies the multiplicity at every zero directly.

\subsection{Three scales and the archimedean comparison}

\begin{example}[Unit spacing and unit scale]
For \(b=c=1\), the eigenvalues are \(s,s+1,s+2,\ldots\), and
\[
 D_{1,1}(s)=\frac{\sqrt{2\pi}}{\Gamma(s)},\qquad
 Z(0,s)=\frac12-s.
\]
The integral values \(\Gamma(1)=1\) and \(\Gamma(1/2)=\sqrt\pi\) give
\[
 D_{1,1}(1)=\sqrt{2\pi},\qquad D_{1,1}(1/2)=\sqrt2.
\]
At \(s=1/2\), the exponent in~\eqref{eq:osc-scale} vanishes, so positive
operator rescaling leaves the determinant unchanged. At \(s=1\),
rescaling by \(d\) multiplies it by \(d^{-1/2}\).
The mode identity reads \(D_{1,1}(s)=sD_{1,1}(s+1)\).
It gives a simple zero at zero with slope \(\sqrt{2\pi}\), while the
tail determinant there equals \(\sqrt{2\pi}\).
\end{example}

\begin{example}[The real archimedean scale]
Use Lebesgue measure \(dx\) on \(\mathbb R\), multiplicative measure
\(d^\times x=dx/|x|\) on \(\mathbb R^\times\), and the test function
\(e^{-\pi x^2}\). For \(\operatorname{Re}s>0\), substitution
\(u=\pi x^2\) gives the absolutely convergent integral
\begin{align}
 \Gamma_{\mathbb R}(s)
 &=\int_{\mathbb R^\times}e^{-\pi x^2}|x|^s\,d^\times x
   =2\int_0^\infty e^{-\pi x^2}x^{s-1}\,dx\notag\\
 &=\pi^{-s/2}\int_0^\infty e^{-u}u^{s/2-1}\,du
   =\pi^{-s/2}\Gamma(s/2).                         \label{eq:osc-real-factor}
\end{align}
The factor \(2\) comes from the two real components, and cancels the
Jacobian factor \(1/2\).
For spacing \(b=2\) and operator scale \(c=2\pi\),
\begin{equation}
 \Delta_\infty(s):=D_{2,2\pi}(s)
 =\pi^{s/2-1/2}\frac{\sqrt{2\pi}}{\Gamma(s/2)}
 =\frac{\sqrt2}{\Gamma_{\mathbb R}(s)}.             \label{eq:osc-archimedean}
\end{equation}
Thus \(\Delta_\infty(1)=\sqrt2\) and
\(\Delta_\infty(2)=\sqrt2\,\pi\).
The entire function has zeros at \(0,-2,-4,\ldots\), with
\(\Delta_\infty'(0)=1/\sqrt2\).
Removing the zero mode instead gives value \(\sqrt2\,\pi\) at zero.
These constants use both the measure and the operator scale just stated.
\end{example}

\begin{corollary}[Completed Euler determinant]
\label{cor:completed-euler}
With the measures of Proposition~\ref{prop:local-tate} and scale $2\pi$,
\[
 \frac1{\Delta_\infty(s)\Delta_1(s)}
 =2^{-1/2}\Lambda(s),\qquad \operatorname{Re}s>1.
\]
Both sides continue as single-valued meromorphic scalar functions.
\end{corollary}
\begin{proof}
Theorem~\ref{thm:euler-fredholm} gives $\Delta_1=1/\zeta$ on the
initial half-plane. Equation~\eqref{eq:osc-archimedean} gives
$\Delta_\infty=\sqrt2/\GammaR$. Multiplying proves the identity;
the scalar continuations prove its meromorphic extension.
\end{proof}
The meromorphic identity changes neither the Schatten domain of the
prime diagonal nor the principal spectral cut of the oscillator.
It fixes the normalization a cohomological realization must recover.

\begin{example}[Prime logarithmic spacing]
Let \(p\) be a prime, \(b=\log p\), and \(c=2\pi\).
The heat operator is the diagonal bounded operator with entries
\(e^{-tbr}\). For \(\operatorname{Re}t>0\), these entries have summable
absolute values. Finite truncation therefore gives its trace as
\[
 \operatorname{Tr}(e^{-tB_b})=\sum_{r\geq0}p^{-tr}
 =\frac1{1-p^{-t}}.
\]
The additive determinant has a different spectral parameter:
\[
 D_{\log p,2\pi}(s)
 =\left(\frac{\log p}{2\pi}\right)^{1/2-s/\log p}
       \frac{\sqrt{2\pi}}{\Gamma(s/\log p)}.
\]
At \(t=1\), the trace is \(p/(p-1)\).
At \(s=\log p\), the determinant is \(2\pi/\sqrt{\log p}\).
Its slope at zero is \(1/\sqrt{\log p}\).
After identifying the two parameter planes by \(t=s\), the reciprocal
heat trace has simple zeros at \(2\pi i\mathbb Z/\log p\).
The additive determinant has simple zeros at \(-\mathbb N_0\log p\).
For example, \(-\log p\) is a zero only of the additive determinant.
Hence a nonzero constant cannot identify the two entire functions.
\end{example}


\subsection{Limits of ordinary spectral traces}

\begin{proposition}[Obstructions to ordinary spectral constructions]
\label{prop:spectral-obstructions}
Let \(H\) be the operator of Proposition~\ref{prop:prime-product}.
For every real \(a>0\), its shifted spectral sum
\[
 \sum_{n\geq1}(a+\log n)^{-w}
\]
has no right half-plane of convergence.  In particular, it does not define
a determinant by the spectral-zeta prescription of
\eqref{eq:osc-det-definition}.
There is also no self-adjoint operator \(C\) on a Hilbert space satisfying
\[
 e^{-tC}\in\mathcal S_1,\qquad
 \operatorname{Tr}(e^{-tC})=\GammaR(t)\quad(t>0).
\]
\end{proposition}

\begin{proof}
For real \(w>0\), the terms with \(e^j\leq n<e^{j+1}\) have sum at
least
\[
 \bigl(\lfloor e^{j+1}\rfloor-\lceil e^j\rceil\bigr)
 (a+j+1)^{-w},
\]
which tends to infinity.  For real \(w\leq0\), the summands do not tend
to zero.  Every right half-plane contains such divergent real parameters.
The spectral zeta function has no initial half-plane to continue.

For the second assertion, boundedness of \(e^{-2C}\) forces the
self-adjoint \(C\) to be bounded below, say by \(E_0\).  Functional
calculus and positivity then give
\[
 \operatorname{Tr}(e^{-(2n+2)C})
 \leq e^{-2E_0}\operatorname{Tr}(e^{-2nC}).
\]
But the gamma recurrence gives
\[
 \frac{\GammaR(2n+2)}{\GammaR(2n)}=\frac n\pi,
\]
which is unbounded.  This contradiction proves the claim.
\end{proof}

The first obstruction persists although \(H\) has compact resolvent.
The second concerns ordinary Hilbert traces of a fixed self-adjoint
operator.  The oscillator determinant in
Proposition~\ref{prop:oscillator-det} supplies a regularized construction
with the required archimedean factor.


\section{Formal coefficients and analytic traces}
\label{chap:formal-analytic-determinants}

A formal logarithm and a holomorphic logarithm satisfy the same
coefficient identities, but they answer different questions. A formal
logarithm determines each coefficient by finitely many operations.
A holomorphic logarithm also has a domain of convergence and a choice
of branch. Continuing that branch can require a new surface even when
the original coefficients are completely determined.

The prime determinant makes this distinction explicit. Its coefficients
belong to a natural Banach algebra precisely in its initial
regularization half-plane. The continued germ then lives on the surfaces
constructed in Sections~\ref{sec:prime-continuation} and~\ref{sec:finite-covers}. The same
diagonal matrices act on a formal product and on a Hilbert space, but
their traces have different continuity properties in those two settings.
We construct these comparisons, beginning with one Gaussian variable.

\subsection{Specializing coefficients without evaluating a formal variable}
\label{sec:da-specialization}

Let $R$ be a commutative algebra over $\mathbb Q$ and let $t$ be a
commuting formal variable. The ring $R[[t]]$ consists of sequences
$\sum_{n\geq0}a_nt^n$ with $a_n\in R$. The coefficient of $t^n$ in
a product is $\sum_{j=0}^na_jb_{n-j}$. Thus multiplication uses a finite
sum at every degree. The $t$-adic topology declares two series close
when their coefficients agree through a prescribed degree.

For $u\in tR[[t]]$, set
\[
 \exp(u)=\sum_{r\geq0}\frac{u^r}{r!},\qquad
 \log(1+u)=\sum_{r\geq1}\frac{(-1)^{r+1}}r u^r.
\]
At degree $n$, only $r\leq n$ can contribute. These are therefore
formal operations with no condition on the size of $a_n$.
Their usual inverse identities hold coefficientwise. For completeness,
write $E(z)=\sum z^r/r!$ and $L(z)=\sum_{r\geq1}(-1)^{r+1}z^r/r$.
Formal differentiation gives $E'=E$ and $L'=(1+z)^{-1}$.
The product and chain rules give
\[
 (L(E(z)-1))'=1,\qquad
 \bigl(E(L(z))(1+z)^{-1}\bigr)'=0.
\]
Their constant terms are zero and one. Characteristic zero therefore
gives $L(E(z)-1)=z$ and $E(L(z))=1+z$.
Substitution of $u\in tR[[t]]$ preserves these finite coefficient
identities.

\begin{proposition}[Coefficient specialization and its limit]
\label{prop:da-specialization}
A unital algebra homomorphism $\phi:R\to\mathbb C$ extends coefficientwise
to an algebra homomorphism
\[
 \phi[[t]]:R[[t]]\longrightarrow\mathbb C[[t]].
\]
It commutes with the displayed formal exponential and logarithm.
For any $t_0\in\mathbb C^\times$, there is no unital
$\mathbb C$-algebra homomorphism $\mathbb C[[t]]\to\mathbb C$
that sends $t$ to $t_0$.
\end{proposition}
\begin{proof}
Each product, exponential, or logarithm coefficient is a finite
polynomial expression over $\mathbb Q$ in the input coefficients.
Applying $\phi$ to that expression proves the first assertions.
For the last assertion, the series
\[
 (1-t/t_0)^{-1}=\sum_{n\geq0}(t/t_0)^n
\]
is an inverse in $\mathbb C[[t]]$. The proposed homomorphism would
send $1-t/t_0$ to zero and its product with this inverse to one.
This is impossible.
\end{proof}

For the commutative scalar sector of formal Gaussian elimination, take
$A=\mathbb C[x_1,\ldots,x_d]$. Its coefficient algebra is
$A[\hbar,\hbar^{-1}][[t]]$. Each $t$-coefficient is polynomial in the
fields and a finite Laurent polynomial in $\hbar$. A scalar
specialization of $A$, together with a number $h_0\ne0$, therefore
defines a coefficientwise map sending $\hbar$ to $h_0$.
The lower Laurent bound may depend on the $t$-degree; no uniform bound
is needed for this map. Its output is still a formal series in $t$.
The proposition supplies no evaluation at a nonzero numerical $t$.
If odd field generators are present, a scalar algebra map sends them
to zero. The same coefficient argument applies to the even interaction.

This statement differs from specializing an arbitrary Laurent series
in $\hbar$. In $\mathbb C((\hbar))$, the element $\hbar-h_0$ is
invertible for every $h_0\ne0$, since its constant term is nonzero.
The same unit argument excludes a unital $\mathbb C$-algebra map
$\mathbb C((\hbar))\to\mathbb C$ with $\hbar\mapsto h_0$.
The coefficient ring is part of the comparison.

\subsection{A convergent Gaussian logarithm and a divergent interaction}
\label{sec:da-gaussian}

Fix a real number $b>0$. Let $\langle f\rangle_b$ be the normalized
Gaussian integral of a polynomial $f$ with density
$e^{-x^2/(2b)}$ on $\mathbb R$. Integration by parts gives
\[
 \langle x^{2r}\rangle_b=(2r-1)!!\,b^r,\qquad
 \langle x^{2r+1}\rangle_b=0,
 \qquad (-1)!!=1.
\]
Indeed, the integral of the derivative of
$x^{2r-1}e^{-x^2/(2b)}$ is zero and gives the recurrence.
Every polynomial times this Gaussian tends to zero at infinity.
For polynomials, the same functional is
\[
 \langle f\rangle_b=
 \left.\exp\!\left(\frac b2\partial_x^2\right)f\right|_{x=0}.
\]
Both sides have the displayed values on every monomial, and the
operator series terminates on each polynomial.
For a coefficient parameter $b$, the polynomial operator
$T_b=\exp(b\partial_x^2/2)$ is called Gaussian contraction; each
polynomial involves only finitely many derivatives. Extend it
coefficientwise to $\mathbb C[x,b][\hbar,\hbar^{-1}][[t]]$.
The numerical integral above is its specialization at a real $b>0$,
followed by evaluation at $x=0$. This is an equality of linear
functionals on polynomials, not of algebra homomorphisms.

For a real polynomial $J\in\mathbb R[x]$, define the formal interaction weight
\[
 Z_J(t)=\sum_{n\geq0}\frac{\langle J^n\rangle_b}{n!}t^n,
 \qquad W_J(t)=\log Z_J(t)\in t\mathbb R[[t]].
\]
This uses the formal operations above and does not assume that the
integral of $e^{tJ}$ converges.
For $I\in t\mathbb C[x,b][\hbar,\hbar^{-1}][[t]]$, define the
formal effective interaction by
$\mathcal W_b(I)=\hbar\log(T_b(e^{I/\hbar}))$.
The coefficientwise exponential and logarithm are defined because
$I/\hbar$ has positive $t$-order and $T_b(1)=1$.
After fixing $\hbar=h_0\ne0$, taking $I=h_0tJ$, and evaluating
at $x=0$, the weight is $Z_J$ and $\mathcal W_b(I)/h_0$ is $W_J$.

\begin{proposition}[Quadratic and quartic interactions]
\label{prop:da-gaussian}
For $J=x^2/2$, the formal series is the Taylor series of
\[
 Z_J(t)=(1-bt)^{-1/2},\qquad
 W_J(t)=-\tfrac12\log(1-bt)
\]
on $|t|<1/b$, with $Z_J(0)=1$ and $W_J(0)=0$.
The normalized integral defines $Z_J$ on $\operatorname{Re}t<1/b$.
Its continuation around $t=1/b$ changes sign and has exact local
ramification index two.

For $J=x^4$, the formal series $Z_J$ has radius of convergence zero.
The normalized real integral of $e^{tx^4}$ diverges for $t>0$ and
is finite for $t\leq0$.
\end{proposition}
\begin{proof}
For the quadratic interaction, the exponent in the integral is
$-(1-bt)x^2/(2b)$. On a compact subset of
$\operatorname{Re}t<1/b$, this integrand and each of its derivatives
in $t$ are bounded by an integrable polynomial times a fixed Gaussian.
Differentiation under the integral is therefore valid.
For real $t<1/b$, a real change of variable in the normalized
Gaussian integral gives $(1-bt)^{-1/2}$.
Both functions are holomorphic on this half-plane, so the equality
holds there. The logarithm is the branch with value zero at $t=0$;
$1-bt$ has positive real part throughout this domain.

For $|t|<1/b$, the series of absolute values of the integrand is
$e^{|t|x^2/2}$. It is Gaussian-integrable. Termwise integration gives
exactly the formal coefficients. The singularity at $t=1/b$ makes
the Taylor radius exactly $1/b$.
The logarithmic derivative is
\[
 \frac{Z_J'}{Z_J}=\frac{b}{2(1-bt)}
                  =-\frac1{2(t-1/b)}.
\]
Its integral about a small positively oriented loop is $-\pi i$,
so the multiplier is $-1$. On the local coordinate
$t=1/b+u^2$, one has $Z_J=c/u$ with $c^2=-1/b$.
This is a simple pole on the two-sheeted local surface. An unramified
meromorphic function has an integral order and cannot have the
base exponent $-1/2$, so two is minimal.

For the quartic interaction, the coefficient of $t^n$ is
$a_n=(4n-1)!!\,b^{2n}/n!$. Its consecutive ratio is
\[
 \frac{a_{n+1}}{a_n}
 =\frac{(4n+3)(4n+1)}{n+1}b^2\longrightarrow\infty.
\]
For every nonzero complex $t$, the terms $a_nt^n$ therefore fail
to tend to zero. This proves the zero radius.
For real $t>0$, the exponent $tx^4-x^2/(2b)$ tends to positive
infinity, so the integral diverges. For $t\leq0$, its integrand is
bounded by the original Gaussian. The integral is then finite.
\end{proof}

The quartic logarithm remains an exact formal series. It cannot be
the Taylor series of a holomorphic germ whose exponential is $Z_J$,
because exponentiating such a germ would give a convergent Taylor
series for $Z_J$. Thus a connected formal expansion does not by itself
produce a holomorphic effective interaction. The quadratic example
also shows that convergence near the initial point does not remove
fractional monodromy elsewhere.

\subsection{The coefficient algebra of a prime product}
\label{sec:da-dirichlet}

The variables in an Euler product are related by multiplication of
integers. To retain these relations before numerical evaluation, let
\[
 \mathscr D_{\mathrm{form}}=
 \left\{\sum_{n\geq1}a_n[n]:a_n\in\mathbb C\right\},
 \qquad [r][s]=[rs].
\]
The brackets denote formal basis symbols, not values of a function.
Multiplication is Dirichlet convolution:
\[
 (a*b)_n=\sum_{d\mid n}a_d b_{n/d}.
\]
There are only finitely many divisors of $n$, so every coefficient
exists. The unit is $[1]$, and $a\mapsto a_1$ is an augmentation.
Let $F_N$ consist of the series with $a_n=0$ for $n<N$.
These subspaces give the coefficient topology. Their intersection
is zero, compatible finite coefficient lists give a unique series,
and $F_NF_M\subset F_{NM}$.
In particular, if $a_1=0$, the coefficient of $[n]$ in $a^r$ is
zero for $2^r>n$. Exponential and logarithm consequently exist on
$F_2$ and $1+F_2$, respectively, by the same finite coefficient
formulas as before.

For a positive integer $m$, put
\[
 \mathscr L_m=-\sum_p\sum_{j\geq m}\frac{[p^j]}j,
 \qquad \mathscr D_m=\exp(\mathscr L_m).
\]
Each coefficient of the first sum is finite. In a fixed quotient
by $F_N$, only primes and exponents with $p^j<N$ occur. Applying
the exponential identity for a finite sum in this quotient gives
\begin{equation}
 \mathscr D_m=\prod_p
 \left((1-[p])\exp\!\sum_{j=1}^{m-1}\frac{[p^j]}j\right).
 \label{eq:da-formal-product}
\end{equation}
The product in~\eqref{eq:da-formal-product} is defined coefficientwise. Unique prime
factorization ensures that, for each $[n]$, only primes dividing
$n$ can contribute. The coefficient of $[p^m]$ in $\mathscr D_m$
is $-1/m$, since the corresponding single-prime logarithm begins
with $-[p^m]/m$.

The formal product does not allow arbitrary numerical substitution.
For any fixed $s\in\mathbb C$, the element
$1-2^s[2]$ is invertible in $\mathscr D_{\mathrm{form}}$, with
inverse $\sum_{r\geq0}2^{rs}[2^r]$.
A unital $\mathbb C$-algebra homomorphism sending $[n]$ to $n^{-s}$
would send this unit to zero. No such homomorphism exists on the
whole formal algebra. For example,
$\sum_{r\geq0}r![2^r]$ has a finite coefficient at every integer,
but its proposed numerical sum diverges at every $s$.

\subsection{A Banach algebra on which evaluation exists}
\label{sec:da-banach}

To control numerical evaluation, fix a real number $\sigma>0$ and
define
\[
 \mathscr D_\sigma=
 \left\{a:\|a\|_\sigma=
                 \sum_{n\geq1}|a_n|n^{-\sigma}<\infty\right\}.
\]
Here the coefficients have their usual complex absolute value.
This norm imposes a condition absent from the formal coefficient
topology. The underlying coefficient inclusion
$\mathscr D_\sigma\subset\mathscr D_{\mathrm{form}}$ preserves
addition and multiplication; it is not an assertion that the norm
topology equals the formal topology.

\begin{lemma}[Summable coefficients and holomorphic evaluation]
\label{lem:da-evaluation}
The space $\mathscr D_\sigma$ is a commutative unital Banach algebra.
On $\mathbb H_\sigma=\{s:\operatorname{Re}s>\sigma\}$, the map
\[
 a\longmapsto\operatorname{ev}(a)(s)=\sum_{n\geq1}a_nn^{-s}
\]
is an algebra homomorphism to bounded holomorphic functions, with
$\sup_{\mathbb H_\sigma}|\operatorname{ev}(a)|\leq\|a\|_\sigma$.
It is the unique norm-continuous extension of evaluation on finite
Dirichlet sums. Moreover,
$\operatorname{ev}(a)(s)\to a_1$ as real $s\to+\infty$.
\end{lemma}
\begin{proof}
Absolute summation and $(rs)^{-\sigma}=r^{-\sigma}s^{-\sigma}$ give
\[
 \|a*b\|_\sigma
 \leq\sum_{r,s\geq1}|a_r||b_s|(rs)^{-\sigma}
 =\|a\|_\sigma\|b\|_\sigma.
\]
Multiplication is therefore continuous and $\|[1]\|_\sigma=1$.
For a norm-Cauchy sequence, every coefficient is Cauchy. Call its
coefficientwise limit $a$. Bounds on finite sums pass to the limit,
and then increasing the finite set shows that $a$ is summable.
The same argument applied to differences proves norm convergence
to $a$. This proves completeness.

The evaluation series converges absolutely and uniformly on
$\operatorname{Re}s\geq\sigma$. The product of two such series can
be rearranged absolutely, giving the convolution formula.
On $\operatorname{Re}s\geq\sigma+\delta$, where $\delta>0$,
its derivative series is bounded in absolute sum by
\[
 \sum_n |a_n|(\log n)n^{-\sigma-\delta}
 \leq\left(\sup_{x\geq1}(\log x)x^{-\delta}\right)\|a\|_\sigma.
\]
The supremum is finite. Uniform convergence of the function and
derivative series on each such half-plane permits termwise
differentiation, proving holomorphy. Truncations converge in norm,
which proves the uniqueness assertion. Finally, dominated
convergence in the absolutely summable majorant
$|a_n|n^{-\sigma}$ gives the limit $a_1$.
\end{proof}

The prime-harmonic divergence proved in
Theorem~\ref{thm:euler-fredholm} makes the following convergence
threshold sharp, including its boundary.

\begin{theorem}[Exact coefficient domain of the regularized determinant]
\label{thm:da-coefficient-domain}
For every integer $m\geq1$ and real $\sigma>0$,
\[
 \mathscr L_m\in\mathscr D_\sigma
 \quad\Longleftrightarrow\quad
 \mathscr D_m\in\mathscr D_\sigma
 \quad\Longleftrightarrow\quad m\sigma>1.
\]
On $\operatorname{Re}s>1/m$, their evaluations are exactly
\[
 \operatorname{ev}(\mathscr L_m)(s)=L_m(s),\qquad
 \operatorname{ev}(\mathscr D_m)(s)=\Delta_m(s),
\]
with the logarithm, real prime powers and normalization of
Theorem~\ref{thm:prime-regularization} and
Section~\ref{sec:finite-covers}. In particular,
$\Delta_m(s)\to1$ along the positive real axis.
\end{theorem}
\begin{proof}
If $m\sigma>1$, then
\[
 \|\mathscr L_m\|_\sigma
 =\sum_p\sum_{j\geq m}\frac{p^{-j\sigma}}j
 \leq\frac1{m(1-2^{-\sigma})}\sum_p p^{-m\sigma}
 <\infty.
\]
The last sum is bounded by $\sum_{n\geq2}n^{-m\sigma}$, which
converges by comparison with its integral.
In the Banach algebra, the exponential series converges in norm,
with norm at most $\exp\|\mathscr L_m\|_\sigma$.
Its coefficient at $[n]$ equals the formal exponential coefficient,
since only finitely many powers contribute. Thus its sum is
$\mathscr D_m$.

If $m\sigma\leq1$, the coefficients at $[p^m]$ of both
$\mathscr L_m$ and $\mathscr D_m$ have absolute value $1/m$.
Consequently each proposed norm is at least
\[
 \frac1m\sum_p p^{-m\sigma}\geq\frac1m\sum_p\frac1p=\infty.
\]
This proves both necessary conditions, including the boundary.

Given $\operatorname{Re}s>1/m$, choose
$1/m<\sigma<\operatorname{Re}s$.
The evaluation homomorphism of Lemma~\ref{lem:da-evaluation}
commutes with the norm-convergent exponential series. Hence
\[
 \operatorname{ev}(\mathscr D_m)(s)
 =\exp\left(-\sum_p\sum_{j\geq m}\frac{p^{-js}}j\right).
\]
The prime powers use the real logarithm of each prime, exactly as
in the initial determinant definition. The right side is
$\exp L_m(s)=\Delta_m(s)$ on this half-plane.
The coefficient at $[1]$ of $\mathscr L_m$ is zero and that of
$\mathscr D_m$ is one. The limiting normalization follows from
Lemma~\ref{lem:da-evaluation}.
\end{proof}

The theorem identifies the actual germ used for continuation.
For $m\geq2$, continuing it does not extend evaluation on the whole
formal algebra. It also does not place its original coefficients
in a larger weighted algebra across $m\sigma=1$.
For instance, at $m=2$ the prime-square coefficients already force
divergence at $\sigma=1/2$, while the continued determinant has
local exponent $1/2$ at $s=1/2$.
The continuation is a statement about that holomorphic germ on its
branched surface, with total residue and period character as
constructed in Sections~\ref{sec:prime-continuation}
and~\ref{sec:finite-covers}.

For $n>m$, the formal identity
\[
 \mathscr D_n=
 \exp\left(\sum_p\sum_{j=m}^{n-1}\frac{[p^j]}j\right)
 \mathscr D_m
\]
holds coefficientwise. Every factor belongs to $\mathscr D_\sigma$
when $\sigma>1/m$. Evaluation therefore gives exactly the multiplier
$H_{m,n}$ of Theorem~\ref{thm:order-comparison} on that overlap.
Outside it, the comparison uses the common continuation cover of
that theorem; the formal identity alone specifies no branch.

\subsection{The same diagonal matrix on two different spaces}
\label{sec:da-operators}

The formal recipient can be specified without an analytic norm.
Let $\mathbb C[[x,y]]/\mathbb C$ be the vector space of formal
functions modulo constants. Its Hamiltonian Lie bracket is
$\{f,g\}=\partial_xf\,\partial_yg-\partial_yf\,\partial_xg$
modulo constants. Both derivatives and every output coefficient are
well-defined in the coefficient topology. Fix a bijection between the primes and the nonconstant
monomials $x^ay^b/(a!b!)$. The normalization only rescales the
monomial basis. In these coordinates the formal space is
\[
 \widehat V=\mathbb C^{\{p\text{ prime}\}}
       \cong\mathbb C[[x,y]]/\mathbb C.
\]
Give each copy of $\mathbb C$ the discrete topology and give
$\widehat V$ the product topology. A continuous linear operator on
$\widehat V$ is exactly a row-finite matrix. Indeed, each output
coordinate must be a continuous scalar functional, so it depends
on finitely many input coordinates. Conversely, finite support of
each row makes every finite collection of output coordinates
continuous. The operator topology requires agreement of whole rows
in each prescribed finite collection.

The same finite sequences form a dense subspace of the Hilbert space
$V_2=\ell^2(\{p\})$, now over the usual complex field, with its
declared orthonormal basis $e_p$. For $s\in\mathbb C$, define the
diagonal matrix
\[
 K(s)e_p=p^{-s}e_p.
\]
It is a continuous endomorphism of $\widehat V$ for every $s$.
No summation over primes is required for a fixed output coordinate.

\begin{proposition}[Operator domains and scalar traces]
\label{prop:da-operator-domains}
Write $\sigma=\operatorname{Re}s$. On $V_2$, the diagonal matrix
$K(s)$ is bounded exactly when $\sigma\geq0$, and compact exactly
when $\sigma>0$. For every integer $m\geq1$,
\[
 K(s)\in \mathcal S_m(V_2)\quad\Longleftrightarrow\quad m\sigma>1.
\]
On that half-plane, $s\mapsto K(s)$ is holomorphic in the
$\mathcal S_m$ norm. For $j\geq m$, the ordinary Hilbert trace is
\[
 \operatorname{Tr}K(s)^j=\sum_p p^{-js},\qquad
 L_m(s)=-\sum_{j\geq m}\frac{\operatorname{Tr}K(s)^j}{j},
\]
where the final double series is absolutely and locally uniformly
convergent.
\end{proposition}
\begin{proof}
For $\sigma\ge0$ the diagonal entries have supremum $2^{-\sigma}$;
for $\sigma<0$ they are unbounded. Finite diagonal truncations converge
in operator norm exactly when $\sigma>0$. At $\sigma=0$ the images
of the coordinate unit vectors form an orthonormal sequence, proving
noncompactness. The singular values, Schatten threshold and
norm-holomorphy are those of Theorem~\ref{thm:euler-fredholm}.
The trace formulas and absolute local uniform convergence are the
normalized logarithm of Theorem~\ref{thm:prime-regularization}.
\end{proof}

For $\sigma<0$, the diagonal matrix still defines a densely defined
closed operator on $V_2$, with domain
\[
 \mathcal D(K(s))=
 \left\{v\in\ell^2:\sum_p p^{-2\sigma}|v_p|^2<\infty\right\}.
\]
Finite sequences make the domain dense. If $v^{(n)}\to v$ and
$K(s)v^{(n)}\to w$ in $\ell^2$, convergence in each coordinate gives
$w_p=p^{-s}v_p$. Since $w\in\ell^2$, $v$ lies in the displayed
domain and $K(s)v=w$. This proves closedness. It does not supply
a trace-class power in that region.

The formal and Hilbert realizations agree algebraically on finite
sequences and on every coordinate where both are defined.
Their topologies do not agree. The coordinate inclusion
$V_2\to\widehat V$ is not continuous: $n^{-1}e_2$ tends to zero
in Hilbert norm, while its $2$-coordinate never becomes zero in
the discrete field. Likewise, $s\mapsto K(s)$ is not continuous
in the formal operator topology on any open complex domain, since
its first diagonal entry $2^{-s}$ is not locally constant.
These facts prevent a silent change from a formal completion to
a holomorphic family of operators.

\subsection{Trace compatibility and Hamiltonian translation}
\label{sec:da-trace-comparison}

For a continuous finite-rank operator on $\widehat V$, write
$T=\sum_{j=1}^k v_j\otimes\lambda_j$, where each $\lambda_j$ depends
on finitely many coordinates, and define
\[
 \operatorname{tr}_{\mathrm{fin}}T=\sum_{j=1}^k\lambda_j(v_j).
\]
Such an expression exists: choose a basis of the image, then finitely
many coordinates separating its basis vectors. Inverting that finite
coordinate matrix expresses the image coordinates as continuous
functionals. On any finite-dimensional subspace containing the image,
the displayed sum is the ordinary matrix trace of the restriction.
Enlarging this subspace adds a zero quotient operator, so the value is
independent of all choices. For a continuous $B$ and finite-rank $T$,
both $BT$ and $TB$ have trace $\sum_j\lambda_j(Bv_j)$.
On a Hilbert space with inner product linear in the first argument,
the corresponding rank-one operator is $h\mapsto\langle h,g\rangle f$
and its trace is $\langle f,g\rangle$. These definitions agree on
their common finite-coordinate operators.

Let $\epsilon^p(v)=v_p$ and let $P_p=e_p\otimes\epsilon^p$ be the rank-one coordinate
projection. Its finite-rank trace is one in both realizations.
For a finite diagonal sum, both traces equal the sum of its
diagonal entries. This is the precise common finite calculation.

On summable diagonal sequences, the Hilbert trace extends to
\[
 \ell^1(\{p\})\longrightarrow\mathbb C,
 \qquad (a_p)\longmapsto\sum_p a_p,
 \qquad \left|\sum_p a_p\right|\leq\sum_p|a_p|.
\]
Truncations converge in this norm. For the formal operator topology,
the projections $P_p$ tend to zero as the prime index escapes every
finite set, but all their traces equal one. Hence finite-rank trace
has no continuous linear extension to the full formal operator
completion, even when the scalar target has its usual topology.
The Hilbert trace extends by a different convergence condition.

An algebraic cyclic extension is also impossible on all row-finite
operators. Enumerate the basis as $e_0,e_1,\ldots$ and define
$Le_0=0$, $Le_{n+1}=e_n$, $Re_n=e_{n+1}$.
Both matrices are row-finite and
\[
 LR=1,\qquad RL=1-P_0,\qquad [L,R]=P_0.
\]
Any cyclic scalar functional on that algebra must vanish on $P_0$,
contradicting the rank-one trace. On the span of
$e_0,\ldots,e_{N-1}$, the truncated commutator is instead
$P_0-P_{N-1}$. Both finite traces are one and cancel.
The second term vanishes in the formal operator limit, but its
trace does not. This identifies the precise discontinuity.

The prime labeling introduces no Hamiltonian symmetry by itself.
The following computation tests that issue on one translation.

\begin{proposition}[The prime diagonal is not translation-invariant]
\label{prop:da-translation}
For the chosen bijection of primes with normalized monomials,
let $q_n$ label $f_n=y^{n+1}/(n+1)!$, for $n\geq0$.
Let $a_x$ be the formal Hamiltonian action
$a_x f=\{x,f\}=\partial_y f$ modulo constants.
For every real $s>0$ and every $n\geq1$,
\[
 [a_x,K(s)]f_n
 =(q_n^{-s}-q_{n-1}^{-s})f_{n-1}\ne0.
\]
Thus $K(s)$ is not an endomorphism commuting with the full
Hamiltonian action on the formal space.
\end{proposition}
\begin{proof}
The factorial normalization gives $a_xf_n=f_{n-1}$ for $n\geq1$.
Applying the two compositions gives $q_n^{-s}f_{n-1}$ and
$q_{n-1}^{-s}f_{n-1}$. The primes $q_n$ and $q_{n-1}$ are
distinct, and the function $r\mapsto r^{-s}$ is strictly decreasing
on positive reals when $s>0$. Their difference is nonzero.
\end{proof}

The comparison also has an explicit realization on coefficient functions.
Introduce one commuting variable $u_p$ for each basis vector and put
\[
 A_{\mathrm{coeff}}=\prod_\alpha\mathbb C u^\alpha,
 \qquad O_v=\sum_p v_pu_p\quad(v\in\widehat V),
\]
where $\alpha$ runs over finitely supported families of nonnegative
integers. The product is well-defined because a fixed multi-index has
only finitely many decompositions $\alpha=\beta+\gamma$.
Use the product topology with discrete coefficients. The formulas
\[
 \mathcal E(u^\alpha)=\left(\sum_p\alpha_p\right)u^\alpha,
 \qquad
 \delta_{K(s)}(u^\alpha)=
       \left(\sum_p\alpha_pp^{-s}\right)u^\alpha
\]
define continuous derivations of this algebra. Each coefficient sum is
finite, and additivity of the displayed weights under
$\alpha=\beta+\gamma$ proves the product rule.
They satisfy
\[
 \mathcal E(1)=0,\qquad \mathcal E(O_v)=O_v,
 \qquad \delta_{K(s)}(O_v)=O_{K(s)v}.
\]
Thus the Euler derivation is the derivation associated with the identity
kernel, while the prime diagonal assigns distinct weights to the chosen
linear coordinates. These are actual maps on the same coefficient
algebra. Their construction uses no infinite scalar trace. The
translation commutator above proves that the prime assignment does not
identify an operator commuting with the Hamiltonian action.

We have constructed the coefficient evaluation map, identified its
exact domain for $\mathscr D_m$, and compared the resulting
determinant with actual trace-class powers on the prime Hilbert
space. The continued scalar germ retains the same normalization and
the same minimal surfaces and finite-cover criteria of
Sections~\ref{sec:prime-continuation} and~\ref{sec:finite-covers}, including the single-valued
case $m=1$. These constructions supply neither a cyclic trace on the
full formal operator algebra nor an equivariant arithmetic matter
representation. An identification with quantum current amplitudes
would require a specified representation and trace ideal, compatible
contractions, and convergence of the associated spatial integrals.
Those maps are additional mathematical data.

\section{Curve zeta functions in Hall operator exchange}
\label{intake:curve:section}

For a curve over a finite field, the local factors also occur as
coefficients in an operator exchange. The underlying algebra now comes
from extensions of torsion sheaves, rather than the divisor monoid.
A torsion sheaf on a smooth projective curve is the direct sum of its
parts supported at closed points. This decomposition turns local Hall
operators into global operators. Their exchange coefficient counts the
points over every finite extension of the ground field.

Let $X$ be a smooth projective geometrically connected curve over
$\mathbb F_q$. For a closed point $x$, put $d_x=[\mathbb F_q(x):\mathbb F_q]$
and $Q_x=q^{d_x}$. The completion of the local ring is
$\mathbb F_{Q_x}[[t]]$, after a choice of parameter and coefficient field.
Its finite-length modules are nilpotent modules over $\mathbb F_{Q_x}[t]$.

\subsection{Local Hall operators and their normalization}

For a prime power $Q$, let $H_Q$ be the complex Hall algebra of these
nilpotent modules. Its basis element $u_M$ denotes the isomorphism class
of $M$. The coefficient of $u_C$ in $u_Au_B$ counts submodules
$B'\subset C$ with $B'\simeq B$ and $C/B'\simeq A$.
Give it the positive Hermitian form, conjugate-linear in the first slot,
\[
 \langle u_M,u_N\rangle=\frac{\delta_{MN}}{|\operatorname{Aut}M|}.
\]
The Green coproduct is adjoint to multiplication for this form.
In this category it is cocommutative and multiplicative.
These are the classical Hall bialgebra identities, on this untwisted
category with zero Euler form. The polynomial presentation below is
Steinitz--Hall's theorem
\cite[Theorem~2.6 and equation~(2.9)]{SchiffmannHall}.

Write $u_{1^r}$ for the zero endomorphism of $\mathbb F_Q^r$ and set
\[
 e_r=Q^{r(r-1)/2}u_{1^r},\qquad
 E(z)=\sum_{r\ge0}e_rz^r,\qquad e_0=1.
\]
Subspace counting gives
$\Delta u_{1^r}=\sum_{i+j=r}Q^{-ij}u_{1^i}\otimes u_{1^j}$.
Thus $\Delta E(z)=E(z)\otimes E(z)$. The elements $e_r$, $r>0$,
freely generate $H_Q$. One can see the triangular presentation directly:
for a Jordan partition $\lambda$, use the successive dimensions of
$\ker t^i/\ker t^{i-1}$. The product of the corresponding semisimple
classes counts filtrations lowered by $t$. Its coefficient on $u_\lambda$
is one, since equality forces the kernel filtration. Every other term
has larger kernel dimensions. The resulting triangular matrix is
invertible in each finite dimension.

Define $p_r$ by the formal identity
\[
 \log E(z)=\sum_{r\ge1}\frac{(-1)^{r-1}}r p_rz^r.
\]
They are primitive polynomial generators. Their norms can be computed
without a choice of an orthonormal basis. Since
$|\operatorname{GL}_r(\mathbb F_Q)|=Q^{r(r-1)/2}\prod_{j=1}^r(Q^j-1)$,
\[
 \langle E(u),E(v)\rangle
 =\sum_{r\ge0}\frac{Q^{r(r-1)/2}(\bar uv)^r}{\prod_{j=1}^r(Q^j-1)}
 =\prod_{j\ge1}(1-Q^{-j}\bar uv)^{-1}.
\]
The last equality follows by comparing coefficients in
$ (1-t)P(Qt)=P(t)$, with constant coefficient one.
Adjunction and primitivity make distinct primitive degrees orthogonal
and give factorial pairings for their powers. Taking logarithms therefore gives
\[
 \langle p_r,p_s\rangle=\delta_{rs}\frac{r}{Q^r-1}.
\]
Consequently the following operators are adjoint on the polynomial domain:
\[
 a_{-r}=\sqrt{Q^r-1}\,M_{p_r},\qquad
 a_r=\frac{r}{\sqrt{Q^r-1}}\frac{\partial}{\partial p_r},\qquad
 [a_r,a_{-s}]=r\delta_{rs}I.
\]
Here $M_{p_r}$ means multiplication by $p_r$, and the square root is positive.
No assertion that these unbounded operators act everywhere on the Hilbert
completion is needed.

\subsection{The global exchange}

The torsion Hall algebra of $X$ is the algebraic tensor product
$H_X^{\mathrm{tor}}=\bigotimes_{x\in|X|}^{\mathrm{alg}}H_{Q_x}$,
with vacuum $1$ in all but finitely many factors. Distinct supports have
no extensions or homomorphisms between them. Their Hall products and
pairings therefore factor. Give a length-$r$ local module degree $rd_x$.
There are finitely many closed points of bounded degree, so each
homogeneous piece of this algebra is finite-dimensional.
Let $a_{x,\pm r}$ be the preceding operators on the $x$ factor, and define
\[
 A_{-n}=\sum_{d\mid n}d\sum_{d_x=d}a_{x,-n/d},\qquad
 A_n=A_{-n}^{\dagger}\qquad(n>0).
\]
These are finite sums on the common polynomial domain. A closed point
of degree $d$ contributes $d$ geometric points over $\mathbb F_{q^n}$
exactly when $d\mid n$. Thus, writing $N_n=\#X(\mathbb F_{q^n})$,
\[
 [A_n,A_{-m}]
 =\delta_{nm}\sum_{d\mid n}d^2\frac nd\#\{x:d_x=d\}I
 =\delta_{nm}nN_nI.
\]
All commutators between two positive or two negative modes vanish.

Put $s=w^{-1}$, with $z,s$ independent formal variables, and set
\[
 \Gamma_-(z)=\exp\left(\sum_{n>0}\frac{A_{-n}}n z^n\right),\qquad
 \Gamma_+(w)=\exp\left(\sum_{n>0}\frac{A_n}n w^{-n}\right).
\]
Each coefficient is a finite sum of polynomial-domain operators.
Products are identities in
$\operatorname{End}_{\mathbb C}(H_X^{\mathrm{tor}})[[z,s]]$.
For the formal zeta function
\[
 Z_X(t)=\exp\left(\sum_{n>0}\frac{N_n}n t^n\right)
       =\prod_{x\in|X|}(1-t^{d_x})^{-1},
\]
the logarithm of the Euler product proves the second equality coefficientwise.
The central commutator of the two exponents equals $\log Z_X(zs)$.
Expanding the exponentials, or differentiating their formal conjugation
identity, therefore proves
\[
 \Gamma_+(w)\Gamma_-(z)
 =Z_X(z/w)\Gamma_-(z)\Gamma_+(w).
\]
This is an operator identity before taking a trace. Repeated reordering
uses the product of the central factors for all exchanged pairs.
The statement specifies a formal expansion at $z/w=0$; an analytic
operator identity would require additional domain and convergence estimates.

For $X=\mathbb P^1$, counting gives $N_n=q^n+1$, hence
$Z_X(t)=((1-t)(1-qt))^{-1}$.
For the smooth projective cubic
$E/\mathbb F_2: y^2+y=x^3+x+1$, the first four point counts are
$1,5,13,25$. The trace criterion for $y^2+y=b$ gives two solutions
when $\operatorname{Tr}_{\mathbb F_{2^n}/\mathbb F_2}(b)=0$, and none otherwise.
For $n=1,2,3,4$, it gives respectively $0,2,6,12$ admissible $x$ values,
with one additional point at infinity. Thus
\[
 Z_E(t)=1+t+3t^2+7t^3+15t^4+O(t^5).
\]
These coefficients determine the corresponding operator exchange to degree four.
The full rational formula
$Z_E(t)=(1-2t+2t^2)/((1-t)(1-2t))$
uses the classical elliptic-curve zeta theorem in addition to these finite counts.

\subsection{Torsion acting on a line bundle}

The full Hall algebra of coherent sheaves also contains the classes of
line bundles. Fix a rational point $x$ and a line bundle $L$.
There are two middle sheaves in an extension of $\mathcal O_x$ by $L$:
the split sheaf $L\oplus\mathcal O_x$, and $L(x)$.
Indeed $\operatorname{Ext}^1(\mathcal O_x,L)$ is one-dimensional,
and all nonzero classes have the latter middle sheaf.
In the split middle sheaf, the relevant copies of $L$ are the graphs of
the $q$ maps $L\to\mathcal O_x$. In $L(x)$ the relevant subobject is
uniquely $L(x)(-x)=L$. Conversely, an extension of $L$ by $\mathcal O_x$
splits because $\operatorname{Ext}^1(L,\mathcal O_x)=0$.
Its unique torsion subobject is $\mathcal O_x$. Therefore
\[
 u_{\mathcal O_x}u_L=q\,u_{L\oplus\mathcal O_x}+u_{L(x)},\qquad
 u_Lu_{\mathcal O_x}=u_{L\oplus\mathcal O_x}.
\]
Let $v=\sqrt q$ and define the twisted product
$u_A*u_B=v^{\chi(A,B)}u_Au_B$, where
$\chi(A,B)=\dim\operatorname{Hom}(A,B)-\dim\operatorname{Ext}^1(A,B)$.
Here $\chi(\mathcal O_x,L)=-1$ and $\chi(L,\mathcal O_x)=1$.
The ordinary commutator for this twisted product consequently satisfies
\[
 [u_{\mathcal O_x},u_L]_* =v^{-1}u_{L(x)}.
\]
On $\mathbb P^1$, all rational-point modifications of $\mathcal O(n)$
are isomorphic to $\mathcal O(n+1)$. Summing over the $q+1$ rational points gives
\[
 \left[\sum_{x\in\mathbb P^1(\mathbb F_q)}u_{\mathcal O_x},
 u_{\mathcal O(n)}\right]_*
 =(v+v^{-1})u_{\mathcal O(n+1)}.
\]
This mixed Hall operation retains the elementary modification of the line
bundle. Identifying a full quantum loop algebra additionally requires
its remaining generators and relations. Neither this mixed operation nor
the torsion exchange constructs a comparison with a global arithmetic
trace category or with the full Weil functional.

\section{Deninger's determinant package}

The oscillator and prime diagonal give the scalar identity of
Corollary~\ref{cor:completed-euler}. An arithmetic cohomological
realization must produce its operators independently, with the same
normalization. The determinant prescription extends to the following
algebraic spectral setting; the required continuation is an explicit
hypothesis.

\begin{definition}[Zeta-regularized determinant]
\label{def:det}
Let \(V=\bigoplus_\alpha V_\alpha\) be a countable algebraic direct sum of
finite-dimensional \(T\)-invariant complex vector spaces.  Eigenvalues are
counted with algebraic multiplicity.  Use
\[
 \lambda^{-z}=\exp\!\bigl(-z(\log|\lambda|+i\arg\lambda)\bigr),
 \qquad -\pi<\arg\lambda\leq\pi.
\]
If \(0\) is an eigenvalue, set \(\det_\infty(T\mid V)=0\).  Otherwise
suppose
\[
 \zeta_T(z)=\sum_{\lambda\in\operatorname{Spec}_{\mathrm{alg}}T}
 m_T(\lambda)\lambda^{-z}
\]
converges absolutely in a right half-plane, extends meromorphically to
\(\operatorname{Re}z>-\varepsilon\) for some \(\varepsilon>0\), and is
holomorphic at zero.  Then
\[
 \det_\infty(T\mid V)=\exp(-\zeta_T'(0)).
\]
For a real space, this definition is applied after complexification.
\end{definition}

When zero is not an eigenvalue and the spectral zeta hypotheses hold,
for \(c>0\) compatible branches give
\[
 \det_\infty(cT)=c^{\zeta_T(0)}\det_\infty(T).
\]
Thus a factor such as \(2\pi\) has a scale anomaly and cannot be suppressed
without computing \(\zeta_T(0)\).

\begin{theorem}[Conditional Deninger determinant formula]
\label{thm:deninger}
Assume that real vector spaces \(H^i\), concentrated in degrees
\(i=0,1,2\), and endomorphisms \(\Theta_i\) are given, and that each
complexification is a countable algebraic direct sum of finite-dimensional
\(\Theta_i\)-invariant subspaces as in Definition~\ref{def:det}.  For
\(s\notin\operatorname{Spec}_{\mathrm{alg}}(\Theta_i)\), apply
Definition~\ref{def:det} to the pencil
\[
 D_i(s)=\det_\infty\left(
 \frac{s-\Theta_i}{2\pi}\,\middle|\,H^i\right).
\]
Assume that each locally defined determinant is the restriction of a
single-valued meromorphic function \(D_i\) on \(\mathbb C\), with the zero
value of Definition~\ref{def:det} at spectral parameters, and that
\[
 (H^0,\Theta_0)=(\mathbb R,0),
 \qquad (H^2,\Theta_2)=(\mathbb R,1),
\]
and that Deninger's cohomological identity holds in his normalization:
\[
 \widehat\zeta_D(s)=\frac{D_1(s)}{D_0(s)D_2(s)},
 \qquad \widehat\zeta_D(s)=2^{-1/2}\Lambda(s).
\]
Then
\[
 D_0(s)=\frac{s}{2\pi},\qquad
 D_2(s)=\frac{s-1}{2\pi},
\]
and
\[
 \Lambda(s)=\sqrt2\,\frac{D_1(s)}{D_0(s)D_2(s)},
 \qquad
 \xi_{\mathbb R}(s)=2\sqrt2\,\pi^2D_1(s).
\]
\end{theorem}

\begin{proof}
The determinants in degrees zero and two are ordinary one-dimensional
determinants.  Substitution into the assumed alternating determinant formula
gives
\[
 D_1(s)=2^{-1/2}\Lambda(s)\frac{s(s-1)}{4\pi^2}
 =\frac{\xi_{\mathbb R}(s)}{2\sqrt2\,\pi^2}.
\]
The assumed factor $2^{-1/2}$ is the normalization in
\cite[\S3, equation~(3)]{Deninger1998}, and it agrees with
Corollary~\ref{cor:completed-euler}.
\end{proof}

The theorem is an implication from the displayed package.  It does not
construct \(H^1\), \(\Theta_1\), or their determinant.  In particular,
\(\det_\infty(\Theta_1)\) is a scalar, whereas
\(\xi_{\mathbb R}(s)\) is a function; the parameter belongs in the pencil
\((s-\Theta_1)/(2\pi)\).

\begin{corollary}[Conditional Hilbert--P\'olya consequence]
\label{cor:hp}
In addition to Theorem~\ref{thm:deninger}, suppose \(H^1_{\mathbb C}\) has a
Hilbert completion \(\mathscr H^1\) on which \(\Theta_1\) has a densely
defined closed extension with compact resolvent.  Assume, with algebraic
multiplicities, the exact multiset equality
\[
 \operatorname{Spec}(\Theta_1)
 =\{\rho:\xi_{\mathbb R}(\rho)=0\},
 \qquad
 m_{\mathrm{alg}}(\rho,\Theta_1)
 =\operatorname{ord}_{s=\rho}\xi_{\mathbb R}(s),
\]
and assume that there is no additional point, continuous, or residual
spectrum.  Then
\[
 \operatorname{Spec}(\Theta_1)\subset\tfrac12+i\mathbb R
\]
is equivalent to the Riemann Hypothesis.

If, for a positive-definite Hilbert inner product,
\[
 \mathcal D(\Theta_1^*)=\mathcal D(\Theta_1),
 \qquad \Theta_1^*=1-\Theta_1,
\]
then \(H_{\mathrm{HP}}=-i(\Theta_1-\tfrac12)\) is self-adjoint.  Under the
same exact spectral comparison, this adjoint identity implies the Riemann
Hypothesis.
\end{corollary}

\begin{proof}
The first equivalence is the stated equality of multisets rewritten in
spectral language.  For the last assertion, put
\(A=\Theta_1-\tfrac12\).  The domain equality and adjoint identity say
\(A^*=-A\), so \(-iA\) is self-adjoint and
\(\operatorname{Spec}(A)\subset i\mathbb R\).
\end{proof}

For an unbounded operator, a formal adjoint formula checked only on a common
dense invariant domain proves only skew-symmetry unless essential
skew-adjointness is also shown. A compatible nondegenerate indefinite
pairing can yield reflected spectral symmetry without critical-line
confinement. In multiplicative language the eigenvalues
of \(e^{\Theta_1}\), not of \(\Theta_1\), would have modulus
\(e^{1/2}\) under the conclusion of Corollary~\ref{cor:hp}.

\section{Weil's full criterion}
\label{sec:weil-full}

A positive realization of the zeros must recover a full arithmetic
quadratic form. The criterion used below is the classical explicit
formula and positivity theorem of Weil; its general contour and
zero-separation theorem are imported in the stated test class. The
normalization calculation records every local term. The subsequent
narrow-window estimates are direct inequalities for this full form.

\begin{definition}[Weil test class]
\label{def:weil-class}
Let \(\mathcal W\) consist of complex-valued functions
\(f:\mathbb R\to\mathbb C\) which are continuously differentiable away
from finitely many points, have finite one-sided limits there for both
\(f\) and \(f'\), and satisfy the convention
\[
 f(a)=\tfrac12\bigl(f(a+)+f(a-)\bigr)
\]
at every exceptional point \(a\).  Suppose also that, for some \(b>0\),
\[
 f(u),f'(u)=O(e^{-(1/2+b)|u|}).
\]
These hypotheses are Weil's conditions (A)--(B)
\cite[p.~259]{Weil1952}.
For \(f\in\mathcal W\), set
\[
 \widetilde f(u)=\overline{f(-u)},\qquad g=f*\widetilde f,
\]
and fix the Fourier convention
\[
 h(t)=\widehat g(t)=\int_{\mathbb R}g(u)e^{itu}\,du,
 \qquad
 g(u)=\frac1{2\pi}\int_{\mathbb R}h(t)e^{-itu}\,dt.
\]
\end{definition}

For every \(0<b'<b\), the convolution \(g\) satisfies the same conditions
with \(b'\) in place of \(b\).  Piecewise integration by parts, uniformly
for \(|\eta|\leq\tfrac12\), gives
\[
 \widehat f(t+i\eta)=O((1+|t|)^{-1}).
\]
The Fourier transform extends to the strip needed below, and
\[
 h(z)=\widehat f(z)\overline{\widehat f(\overline z)},
 \qquad h(t)=|\widehat f(t)|^2\quad(t\in\mathbb R).
\]
Equivalently, for the centered Mellin transform
\[
 \Phi_f(s)=\int_{\mathbb R}f(u)e^{(s-1/2)u}\,du,
\]
one has
\[
 h\!\left(-i(s-\tfrac12)\right)
 =\Phi_f(s)\overline{\Phi_f(1-\overline s)}.
\]

Write $\psiGamma=\Gamma'/\Gamma$ for the digamma function.
The von Mangoldt function is $\vmLambda(n)=\log p$ when $n=p^a$
with $a\ge1$, and zero otherwise.

\begin{theorem}[Weil explicit formula and positivity criterion]
\label{thm:weil}
Let \(f\in\mathcal W\), \(g=f*\widetilde f\), and
\(h=\widehat g\).  Write a nontrivial zero as
\(\rho=\beta_\rho+i\gamma_\rho\), with multiplicity \(m(\rho)\), and set
\[
 \sum_\rho^{*}h\!\left(-i(\rho-\tfrac12)\right)
 :=\lim_{\substack{T\to\infty\\
                    T\notin\{|\gamma_\rho|\}}}
 \sum_{|\gamma_\rho|<T}m(\rho)
 h\!\left(-i(\rho-\tfrac12)\right).
\]
This limit is in fact absolute for the stated convolution-square class.
Then
\[
\begin{aligned}
 \sum_\rho^{*}h\!\left(-i(\rho-\tfrac12)\right)
 ={}&h(i/2)+h(-i/2)\\
 &+\frac1{2\pi}\int_{\mathbb R}h(t)
 \left[\operatorname{Re}\psiGamma\!\left(\frac14+\frac{it}{2}\right)
 -\log\pi\right]dt\\
 &-\sum_{n\geq2}\frac{\vmLambda(n)}{\sqrt n}
 \bigl(g(\log n)+g(-\log n)\bigr),
\end{aligned}
\]
The Riemann Hypothesis
holds if and only if this full common value is nonnegative for every
\(f\in\mathcal W\).
\end{theorem}

\begin{proof}[Normalization and consequences of the classical theorem]
Moving the logarithmic-derivative contour across the two poles of
\(\Lambda(s)\) gives \(h(i/2)+h(-i/2)\).  Fourier inversion applied to
\[
 -\frac{\zeta'}{\zeta}(s)
 =\sum_{n\geq2}\frac{\vmLambda(n)}{n^s}
\]
gives the two prime-power terms.  Finally,
\[
 \frac{\GammaR'}{\GammaR}(s)
 =-\frac12\log\pi+\frac12\psiGamma(s/2),
\]
so the two sides of the critical line give the kernel
\[
 2\operatorname{Re}\frac{\GammaR'}{\GammaR}
 \left(\frac12+it\right)
 =\operatorname{Re}\psiGamma\!\left(\frac14+\frac{it}{2}\right)-\log\pi.
\]
The decay in Definition~\ref{def:weil-class} justifies the contour passage,
the prime sum, and the archimedean integral.  The estimate preceding the
theorem gives
\(h(-i(\rho-\tfrac12))=O((1+|\gamma_\rho|)^{-2})\); together with
the standard bound \(N(T)=O(T\log T)\), it also proves the asserted absolute
convergence.  This is Weil's formula~(11)
in the present Fourier normalization
\cite[p.~262, formula~(11)]{Weil1952}.

If the Riemann Hypothesis holds, write
\(\rho=\tfrac12+i\gamma\); every zero contributes
\(h(\gamma)=|\widehat f(\gamma)|^2\geq0\).  Conversely, Weil's
Gaussian-polynomial construction associated with an off-line zero produces
an \(f\in\mathcal W\) for which the zero-side sum is negative
\cite[pp.~262--263, unnumbered lemma immediately following
formula~(11)]{Weil1952}.
\end{proof}

When \(g\) is real and even, the last line of the formula becomes
\[
 -2\sum_p\sum_{m\geq1}\frac{\log p}{p^{m/2}}g(m\log p).
\]
Both the polar terms and this prime sum are essential to the criterion.

\begin{proposition}[The gamma term is not positive by itself]
\label{prop:gamma-negative}
Set
\[
 k(t)=\operatorname{Re}\psiGamma\!\left(\frac14+\frac{it}{2}\right)-\log\pi.
\]
There are functions \(f\in C_c^\infty(\mathbb R)\subset\mathcal W\) for
which
\[
 \frac1{2\pi}\int_{\mathbb R}|\widehat f(t)|^2k(t)\,dt<0.
\]
\end{proposition}

\begin{proof}
The special value
\[
 k(0)=-\gamma-\frac\pi2-3\log2-\log\pi
\]
is negative.  Choose nonzero \(\varphi\in C_c^\infty(\mathbb R)\) and put
\(f_R(u)=R^{-1/2}\varphi(u/R)\).  Then
\[
 |\widehat f_R(t)|^2=R|\widehat\varphi(Rt)|^2
\]
and therefore
\[
 \frac1{2\pi}\int_{\mathbb R}|\widehat f_R(t)|^2k(t)\,dt
 =\frac1{2\pi}\int_{\mathbb R}|\widehat\varphi(v)|^2k(v/R)\,dv.
\]
Since \(k(t)=O(\log(2+|t|))\) and \(\widehat\varphi\) is rapidly
decreasing, dominated convergence makes the last expression tend to
\[
 \frac{k(0)}{2\pi}
 \int_{\mathbb R}|\widehat\varphi(v)|^2\,dv<0.
\]
It is therefore negative for all sufficiently large \(R\).
\end{proof}

Proposition~\ref{prop:gamma-negative} shows that the isolated gamma term is
not nonnegative on all Weil convolution squares and therefore cannot replace
the full form in Theorem~\ref{thm:weil}.  The theorem-level criterion uses
the entire test class and every local term.

\section{Positive sectors of the Weil form}
\label{intake:weil:section}

The broad dilations in Proposition~\ref{prop:gamma-negative}
concentrate the Fourier mass near zero. Narrow support instead forces
enough Fourier mass into the region where the archimedean kernel is positive.
The prime-power terms of the full Weil form couple separated intervals.
A sufficiently small width makes the archimedean contribution dominate these couplings.
This gives a positive space of functions with infinitely many independent directions,
while retaining the full arithmetic functional.

\subsection{Conventions and the full form}
For $f\in C_c^\infty(\mathbb R)$ use
\[
 \widehat f(t)=\int_\mathbb R f(u)e^{itu}\,du,\qquad
 f^*(u)=\overline{f(-u)},\qquad
 (f*g)(u)=\int_\mathbb R f(v)g(u-v)\,dv.
\]
Let $\vmLambda(n)$ be the von Mangoldt function: it is $\log p$ for $n=p^a$ and zero otherwise. Put
\[
 k_\infty(t)=\operatorname{Re}\psiGamma(1/4+it/2)-\log\pi,
\]
where $\psiGamma=\Gamma'/\Gamma$. Define the full arithmetic functional by
\[
\begin{split}
 L_\zeta(h)={}&\widehat h(i/2)+\widehat h(-i/2)
 +\frac1{2\pi}\int_\mathbb R\widehat h(t)k_\infty(t)\,dt\\
 &-\sum_{n\ge2}\frac{\vmLambda(n)}{\sqrt n}
          \bigl(h(\log n)+h(-\log n)\bigr).
\end{split}
\]
This is the full functional in Theorem~\ref{thm:weil}, restricted to
compact smooth tests, with exactly the same Fourier normalization.
For compact smooth $h$ the prime sum is finite and the Fourier integral converges. Define
\[
 B_\zeta(f,g)=L_\zeta(f^**g),\qquad Q_\zeta(f)=B_\zeta(f,f).
\]
The form is conjugate-linear in its first argument, linear in its second and Hermitian. The estimates below use this displayed functional directly.

\subsection{A coercive estimate on one narrow interval}
The convergent digamma series gives
\[
 k_\infty(t)=k_\infty(0)+\sum_{n\ge0}
 \left(\frac1{n+1/4}
 -\frac{n+1/4}{(n+1/4)^2+t^2/4}\right).
\]
Every summand is nonnegative and increasing with $|t|$. Also
\[
 k_\infty(0)=-\gamma-\frac\pi2-3\log2-\log\pi>-6.
\]
For example $\gamma<1$, $\pi<22/7$, $\log2<7/10$ and $\log\pi<6/5$ suffice. The last two exponential inequalities follow from their finite positive Taylor lower bounds. Finally $k_\infty(T)\to+\infty$: for $n+1/4\le T/2$, the corresponding summand in the displayed digamma series is at least $1/(2(n+1/4))$, whose partial sums diverge.

Let $f$ have support in an interval of length $L$. Cauchy--Schwarz gives $|\widehat f(t)|^2\le L\|f\|_2^2$, and Plancherel gives
\[
 \frac1{2\pi}\int_{-T}^T|\widehat f(t)|^2dt
       \le\frac{LT}{\pi}\|f\|_2^2,\qquad
 \frac1{2\pi}\int_\mathbb R|\widehat f(t)|^2dt=\|f\|_2^2.
\]
Let $T>0$ and $M\ge-6$ satisfy $k_\infty(T)\ge M$.
The archimedean contribution is then at least
\[
 \left(M-(M+6)\frac{LT}{\pi}\right)\|f\|_2^2.
\]
When $L<\log2$, the support of $f^**f$ contains no nonzero prime-power logarithm, so its prime contribution is zero. The polar term is bounded below by $-2Le^{L/2}\|f\|_2^2$: translate the interval to be centred at zero and apply Cauchy--Schwarz to $\int f(u)e^{u/2}du$ and $\int f(u)e^{-u/2}du$. The translation factors cancel between these two integrals. Consequently
\[
 Q_\zeta(f)\ge c(L,T,M)\|f\|_2^2,
 \qquad c(L,T,M)=M-(M+6)\frac{LT}{\pi}-2Le^{L/2}.
\]
The estimate holds uniformly over all smooth functions in the interval.

\subsection{The cross-window kernel retains the arithmetic primes}
Away from zero, the distribution given by the sum of the polar and archimedean contributions in the displayed definition of $L_\zeta$ is the smooth even kernel
\[
 K(u)=2\cosh(u/2)-\frac{e^{-|u|/2}}{1-e^{-2|u|}}\qquad(u\ne0).
\]
To verify it, use the integral expression
\[
 \psiGamma(a+it/2)=-\gamma+
 \int_0^\infty\frac{e^{-v}-e^{-(a+it/2)v}}{1-e^{-v}}\,dv.
\]
Test on a smooth function supported away from zero. Terms constant in $t$ vanish on that test. Fourier inversion of the cosine term, followed by $v=2|u|$, gives the negative second summand in the displayed formula for $K$. The two polar evaluations give $2\cosh(u/2)$. This argument is on separated supports; it does not ignore a delta term at zero in the diagonal form.

Fix distinct positions $a_1,\ldots,a_N$. Let $I_L=(-L/2,L/2)$ and
\[
 F(u)=\sum_{j=1}^N f_j(u-a_j),\qquad f_j\in C_c^\infty(I_L).
\]
Choose $L$ so small that these supports are disjoint and each difference interval contains a prime-power logarithm only if its centre equals that logarithm. This is possible because there are only finitely many prime powers in a bounded interval. Put
\[
 w_{ij}=\begin{cases}
 \vmLambda(n)/\sqrt n,&|a_i-a_j|=\log n\text{ for an integer }n\ge2,\\
 0,&\text{otherwise}.
 \end{cases}
\]
For a resonant pair its prime term is exactly $-w_{ij}\int f_j(v)\overline{f_i(v)}dv$. The smooth part is the kernel $K$ evaluated at $a_j-a_i+v-u$. If $M_{ij}$ bounds $|K|$ there, its absolute value is at most $LM_{ij}\|f_i\|_2\|f_j\|_2$. Thus
\[
 Q_\zeta(F)\ge
 \left(c(L,T,M)-\max_i\sum_{j\ne i}(w_{ij}+LM_{ij})\right)
       \sum_i\|f_i\|_2^2.
\]
Apply $2ab\le a^2+b^2$ to each unordered pair to obtain this bound.

\begin{theorem}[Positive arithmetic windows of arbitrary finite multiplicity]\label{intake:weil:thm:any-windows}
For any fixed finite set of distinct real positions $a_1,\ldots,a_N$, there exists $L>0$ such that the full Weil form is positive definite on the infinite-dimensional space
\[
 \bigoplus_{j=1}^N C_c^\infty(a_j+I_L).
\]
It has a uniform strictly positive lower bound by the $L^2$ norm. All prime-power terms in the displayed definition of $L_\zeta$ are retained.
\end{theorem}
\begin{proof}
The finite row sums of the $w_{ij}$ are fixed. Choose $M$ larger than their maximum plus two, and then $T$ with $k_\infty(T)\ge M$. For sufficiently small $L$, the separation requirements hold, $LT$ and $Le^{L/2}$ tend to zero, and the $M_{ij}$ remain bounded because their arguments stay away from zero. The coefficient in the row-sum estimate is then positive, in fact larger than one after making $L$ smaller. 
\end{proof}

\subsection{An explicit four-window theorem with a prime-power coupling}
Choose
\[
 (a_0,a_1,a_2,a_3)=(0,\log2,\log3,\log4),\qquad L=10^{-3}.
\]
The resonant pairs are $(0,1),(0,2),(0,3),(1,3)$; their weights are respectively
\[
 \frac{\log2}{\sqrt2},\quad\frac{\log3}{\sqrt3},\quad
 \frac{\log2}{2},\quad\frac{\log2}{\sqrt2}.
\]
In particular the $\log4$ term is a genuine prime-power term, not $\log4/2$. Every other logarithmic difference is separated from prime-power logarithms by more than $L$.

\begin{theorem}[A quantitative positive sector]\label{intake:weil:thm:four-windows}
For arbitrary $f_j\in C_c^\infty((-1/2000,1/2000))$, the function
\[
 F(u)=f_0(u)+f_1(u-\log2)+f_2(u-\log3)+f_3(u-\log4)
\]
satisfies
\[
 \boxed{\displaystyle Q_\zeta(F)\ge\frac15\sum_{j=0}^3\|f_j\|_2^2.}
\]
\end{theorem}
\begin{proof}
A finite rational calculation gives
\[
 \sum_{n=0}^{99}\left(\frac1{n+1/4}
 -\frac{n+1/4}{(n+1/4)^2+2500}\right)>8.
\]
For an explicit integer calculation, put
\[
 d_n=(4n+1)((4n+1)^2+40000),\qquad
 b_j=\sum_{n=10j}^{10j+9}\left\lfloor\frac{1600000000}{d_n}\right\rfloor.
\]
The values for $j=0,\ldots,9$ are
\[
 (b_0,\ldots,b_9)=(64860,6514,3300,1958,1248,832,574,411,300,226).
\]
Their sum is $80223$. The rational partial sum is therefore at least
$80223/10000>8$. Since all later summands are nonnegative and $k_\infty(0)>-6$, we have $k_\infty(100)>2$. The local coercivity estimate, $\pi>3$, and $e^{1/2000}<3/2$ give a diagonal lower bound
\[
 2-\frac8{30}-\frac3{1000}=\frac{5191}{3000}.
\]
All cross-kernel arguments have absolute value between $\log(4/3)-10^{-3}$ and $\log4+10^{-3}$. On this interval each of the two positive terms in the displayed formula for $K$ is less than three. For the denominator term, $\log(4/3)>1/4$ and $\log(3/2)<1/2-1/8+1/24<0.498$ imply $1-e^{-2|u|}>1/3$. For the polar term use $2\cosh((\log4+10^{-3})/2)<2(1.001)+1/2<3$. Thus $|K(u)|<3$, and the smooth row cost is less than $9/1000$.

The largest prime row sum is bounded by
\[
 \frac12+\frac{11}{17}+\frac7{20}=\frac{509}{340},
\]
using $\log2<7/10$, $\sqrt2>7/5$, $\log3<11/10$, $\sqrt3>17/10$. Substitution into the row-sum estimate gives
\[
 \frac{5191}{3000}-\frac9{1000}-\frac{509}{340}=\frac{5719}{25500}>\frac15.
\]
This proves the stated four-window inequality for the entire stated infinite-dimensional space.
\end{proof}

This sector is not a convolution algebra. Its support windows are not preserved under repeated convolution, and the permissible window width in Theorem~\ref{intake:weil:thm:any-windows} depends on the whole finite set of positions. Neither theorem supplies a dense convolution-stable exhaustion with one uniform estimate. That missing extension is exactly where the global Weil-positivity problem remains. Global positivity on the full convolution algebra is not a consequence of these estimates.


\subsection{Finite-rank operators from the arithmetic pairing}

Let $\mathscr H$ be the Hilbert completion of this four-window space for
$B_\zeta$, with the inner product conjugate-linear in its first argument.
For $f,g\in\mathscr H$, define the bounded rank-one operator
$\Theta_{f,g}h=B_\zeta(g,h)f$. Direct substitution gives
\[
 \Theta_{f,g}\Theta_{u,v}=B_\zeta(g,u)\Theta_{f,v},\qquad
 \Theta_{f,g}^*=\Theta_{g,f},\qquad
 \operatorname{Tr}(\Theta_{f,g}\Theta_{u,v})
 =B_\zeta(g,u)B_\zeta(v,f).
\]
The trace identity follows from
$\operatorname{Tr}(\Theta_{f,v})=B_\zeta(v,f)$ in an orthonormal basis.
For vectors in the original test space, every displayed contraction has
its prime, archimedean and polar terms given by the full functional above.
These operators represent the selected quadratic pairing. A representation
of convolution would additionally require compatible operators for every
convolution multiplier on a common invariant domain.


For a finite-dimensional flavour subspace $V\subset\mathscr H$, take the
charged fermion space
\[
 \mathcal F_V=\Lambda_{\mathbb C}\!\left(
 \bigoplus_{r\in\mathbb Z_{>0}-1/2}(V\oplus\overline V)t^r\right).
\]
Use its algebraic finite-particle domain and the induced positive Hermitian form.
Exterior multiplication and its adjoint satisfy the canonical
anticommutation relations with contraction $B_\zeta$.
Writing the resulting charged fields as $\psi_f(z)$, linear in $f$,
and $\psi_g^*(w)$, conjugate-linear in $g$,
the vacuum contraction in the expansion $|w|<|z|$ is
\[
 B_\zeta(g,f)\sum_{k\ge0}z^{-k-1}w^k
 =\frac{B_\zeta(g,f)}{z-w}.
\]
For a finite-rank $A$, choose $V$ containing the ranges of $A$ and $A^*$,
and define $J_{\Theta_{f,g}}(z)=:\psi_f(z)\psi_g^*(z):$, extending linearly in $A$.
Normal ordering places creation operators to the left, with fermion signs.
The singular operator product is
\[
 J_A(z)J_B(w)\sim
 \frac{\operatorname{Tr}(AB)}{(z-w)^2}
 +\frac{J_{[A,B]}(w)}{z-w}.
\]
Indeed the two single contractions compose the matrices in their two
orders, with opposite fermion signs. The double contraction sums their
closed index pair and gives $\operatorname{Tr}(AB)$.
Equivalently, the current modes satisfy
\[
 [J_A(m),J_B(n)]=J_{[A,B]}(m+n)
 +m\delta_{m+n,0}\operatorname{Tr}(AB)I.
\]
For $m>0$, exactly $m$ half-integer modes cross the vacuum in the scalar
term. Antisymmetry determines the remaining signs.
For fixed mode and finite-energy vector, only finitely many creation
terms survive. Enlarging $V$ preserves these identities, so they apply
to every finite set of finite-rank operators on $\mathscr H$.
The arithmetic pairing thus occurs as the contraction of singular
currents on this specified positive sector. Its support is still not
stable under convolution.

\section{The arithmetic orientation and finite torsors}
\label{sec:do26-orientation}

The finite arithmetic Chern--Simons invariant uses two maps.  A torsor
pulls a group-cohomology class to an arithmetic curve.  An arithmetic trace
then evaluates that class.  The first map is a construction in the category
of sheaves.  The second uses class field theory.  We separate these maps
and keep the primes dividing the coefficient order in the base scheme.

Let \(F\) be a totally imaginary number field, fix \(n\geq1\),
and put \(X=\operatorname{Spec}\mathcal O_F\).  Throughout this section,
\(H^q(X,-)\) means ordinary sheaf cohomology on the small \'etale site
\(X_{\mathrm{et}}\).  Its objects are schemes \'etale over \(X\), and
its coverings are jointly surjective \'etale families.  Coefficients are
sheaves of abelian groups, with cohomological grading.  In particular,
\(\mathbb G_m(U)=\Gamma(U,\mathcal O_U)^\times\), and
\(\underline{\mathbb Z/n}\) is the constant sheaf.  No completion or
inverse limit of coefficient sheaves occurs.  We impose no compact-support
condition and use no compactified site in this construction.

\subsection{Roots of unity and the power-map defect}

For a field \(L\) of characteristic zero, put
\(\mu_n(L)=\{u\in L^\times:u^n=1\}\).
Suppose \(\mu_n(\overline F)\subset F\), and choose
\(\zeta\in F\) of exact order \(n\).  For \(n=1\), this means
\(\zeta=1\).  The corresponding group isomorphism is
\[
 \zeta_n:\tfrac1n\mathbb Z/\mathbb Z\longrightarrow\mu_n(F),
 \qquad k/n\longmapsto\zeta^k.
\]
Write \(\mu_{n,X}=\ker([n]:\mathbb G_m\to\mathbb G_m)\), where
\([n](u)=u^n\).  This is a sheaf on the specified small site.

\begin{lemma}[The coefficient map at every finite prime]
\label{lem:do26-roots}
The map
\[
 e_\zeta:\underline{\mathbb Z/n}\longrightarrow\mu_{n,X},
 \qquad k\longmapsto\zeta^k,
\]
is an isomorphism.  Moreover,
\[
 \operatorname{Hom}_{X_{\mathrm{et}}}
   (\underline{\mathbb Z/n},\mathbb G_m)=\mu_n(F).
\]
\end{lemma}

\begin{proof}
A root of unity is an algebraic integer, and its inverse is an algebraic
integer.  Thus \(\zeta\in\mathcal O_F^\times\), so the displayed map
is defined on all of \(X\).  At a geometric point over a finite prime,
the stalk of \(\mathbb G_m\) is
\((\mathcal O_{X,\bar x}^{\mathrm{sh}})^\times\).
The strict henselization is a discrete valuation domain of characteristic
zero.  Its fraction field contains \(F\), and
\[
 T^n-1=\prod_{k=0}^{n-1}(T-\zeta^k)
\]
has exactly those \(n\) distinct roots in that field.  The same argument
applies to the geometric generic stalk.  Therefore \(e_\zeta\) is
bijective on every geometric stalk, which proves the first assertion.
A homomorphism from the constant sheaf \(\underline{\mathbb Z}\) to
any abelian sheaf is determined by the image of its global section \(1\).
Factoring through \(\underline{\mathbb Z/n}\) requires that image to be
killed by \(n\).  For \(\mathbb G_m\), these sections are precisely
\(\mu_n(\mathcal O_F)=\mu_n(F)\).
\end{proof}

The stalk in this proof is a strict henselian local ring.  It is not the
geometric special fiber of the group scheme \(\mu_n\).  At a prime
dividing \(n\), that group scheme need not be \'etale.

For example, take \(F=\mathbb Q(i)\), \(n=2\), and the prime
\((1-i)\).  In the strict henselian local ring, normalize the valuation
by \(v(1-i)=1\).  The unit \(i\) has no square root there.  Indeed,
\(u^2=i\) would imply \(\bar u=1\) in the residue field of
characteristic two.  Both \(u-1\) and \(u+1\) would have positive
valuation, whereas
\[
  v\bigl((u-1)(u+1)\bigr)=v(i-1)=1.
\]
This is impossible.  Thus the \'etale power map can fail to be surjective,
even under the root-of-unity hypothesis.

Define its image and cokernel in the category of abelian sheaves:
\[
 K_n=\operatorname{im}([n]:\mathbb G_m\to\mathbb G_m),
 \qquad Q_n=\operatorname{coker}([n]:\mathbb G_m\to\mathbb G_m).
\]
The cokernel is a sheaf cokernel.  Its sections need not equal the quotient
of the groups of sections.

\begin{lemma}[Nonzero defect at every coefficient prime]
\label{lem:coefficient-prime-defect}
Let \(v\) be a finite prime of residue characteristic \(p\mid n\).
The stalk of \(Q_n\) at a geometric point over \(v\) is nonzero.
This assertion holds for every number field, without a root-of-unity
hypothesis.
\end{lemma}

\begin{proof}
Let \(R\) be the strict henselization at that point and choose a
uniformizer \(\varpi\), with valuation \(v(\varpi)=1\).
The unit \(1+\varpi\) is not a \(p\)-th power in \(R\).
Indeed, if \(x^p=1+\varpi\), reduction to the residue field gives
\(\bar x=1\), so \(x=1+a\) with \(v(a)\geq1\). In
\[
 (1+a)^p-1=\sum_{k=1}^{p-1}\binom pk a^k+a^p,
\]
every summand has valuation at least two: \(p\) divides each
intermediate binomial coefficient and \(p\geq2\).
This contradicts \(v(\varpi)=1\). An \(n\)-th power is a
\(p\)-th power, so \(1+\varpi\) represents a nonzero class in
\(R^\times/(R^\times)^n\), the stalk of \(Q_n\).
\end{proof}

\begin{lemma}[Unramified norms on units]
\label{lem:do26-norm}
Let \(R\) be a henselian discrete valuation ring with finite residue
field \(k\).  Let \(R'/R\) be its unramified extension of degree
\(d\geq1\).  Then
\(N_{R'/R}:R'^\times\to R^\times\) is surjective.
\end{lemma}

\begin{proof}
Write \(|k|=q\).  The norm on residue fields sends
\(x\in\mathbb F_{q^d}^\times\) to
\(x^{(q^d-1)/(q-1)}\).  Since that multiplicative group is cyclic,
the image has order \(q-1\).  It is therefore all of \(k^\times\).

Choose an \(R\)-basis of the free module \(R'\).  In those coordinates,
the norm is the determinant of multiplication and hence a polynomial.
Fix \(a\in R^\times\), and choose a residue unit \(\bar b\) with
norm \(\bar a\).  The derivative of the norm at \(\bar b\) is
\[
 h\longmapsto N(\bar b)
       \operatorname{Tr}_{\mathbb F_{q^d}/\mathbb F_q}(h/\bar b).
\]
The trace is a nonzero linear functional because the finite-field
extension is separable.  Some coordinate derivative is therefore nonzero.
Lift all coordinates of \(\bar b\) to \(R\), and hold all but that
coordinate fixed.  Hensel's lemma solves the resulting one-variable norm
equation with the prescribed residue.  The solution is a unit of \(R'\)
with norm \(a\).
\end{proof}

\begin{lemma}[Cohomology of the power-map defect]
\label{lem:do26-defect}
For every number field \(F\) and every \(n\geq1\),
\[
 \operatorname{supp}(Q_n)=\{v:v\mid n\},
 \qquad H^q(X,Q_n)=0\quad(q\geq1).
\]
These conclusions do not require roots of unity in \(F\).
\end{lemma}

\begin{proof}
If \(n\) is invertible at a geometric point, the equation
\(T^n-u=0\) for a unit \(u\) has a simple root after a finite
separable residue extension.  Hensel's lemma lifts that root.  Thus the
stalk of \(Q_n\) vanishes away from the finite set \(V(n)\).
Lemma~\ref{lem:coefficient-prime-defect} proves nonvanishing at
each point of that set; for \(n=1\), both sides are empty.
For its reduced closed immersion \(i:V(n)\hookrightarrow X\), the
adjunction map \(Q_n\to i_*i^*Q_n\) is an isomorphism on every stalk.
Closed direct image is exact on the \'etale site and preserves injectives,
since its inverse-image left adjoint is exact.  Consequently it suffices
to compute cohomology over each finite residue field.

Fix a prime \(v\mid n\).  Put \(R=\mathcal O_{X,v}^{\mathrm h}\)
and let \(R^{\mathrm{sh}}\) be a strict henselization with residue
field \(\bar k\).  The relevant discrete
\(G_k=\operatorname{Gal}(\bar k/k)\)-module is
\[
 M=(R^{\mathrm{sh}})^\times/
                \bigl((R^{\mathrm{sh}})^\times\bigr)^n.
\]
It is killed by \(n\).  Write \(\phi\) for arithmetic Frobenius,
which generates the dense cyclic subgroup of \(G_k\cong\widehat{\mathbb Z}\).
For a discrete torsion module over this group,
\[
 H^1(G_k,M)=M/(\phi-1)M,\qquad H^q(G_k,M)=0\quad(q\geq2).
\]
Here cochains are continuous.  The first formula follows by evaluating
a cocycle on \(\phi\).  An element with orbit length \(r\) and order
dividing \(e\) gives a cocycle on the cyclic quotient of order \(re\).
For the second formula, use the two-periodic cyclic resolution and take
the limit over finite quotients.  If the action factors through a cyclic
quotient of order \(r\), inflation to the quotient of order \(rs\)
has a factor \(s^{\lfloor q/2\rfloor}\) in degree \(q\).
Choosing \(s\) divisible by the exponent kills each class with
\(q\geq2\).  Every continuous cochain has finite image and lies in
such a finite torsion submodule.

We prove \((\phi-1)M=M\).  Represent \(m\in M\) by a unit
\(u\in R_r^\times\), where \(R_r/R\) is finite unramified of
degree \(r\).  Choose the unramified extension \(R_{rn}/R_r\)
of degree \(n\), and put \(\sigma=\phi^r\) on it.  By
Lemma~\ref{lem:do26-norm}, there is a unit \(b\) with
\(N(b)=u\).  Therefore
\[
 z=u b^{-n}\quad\text{satisfies}\quad N(z)=1.
\]
The elementary cyclic form of Hilbert's theorem 90 gives
\(z=\sigma(w)/w\) in the fraction field of \(R_{rn}\).
For completeness, choose \(t\) for which
\[
 w=\sum_{j=0}^{n-1}
    \left(\prod_{\ell=0}^{j-1}\sigma^\ell(z^{-1})\right)\sigma^j(t)
 \ne0.
\]
Such a \(t\) exists by the linear independence of distinct field
automorphisms.  The equation \(N(z)=1\) makes
\(\sigma(w)=zw\) by cyclic reindexing.  Since the extension is
unramified, a uniformizer of \(R_r\) has valuation one upstairs.
Multiplying \(w\) by a suitable power of that uniformizer makes it a
unit and does not change \(\sigma(w)/w\).
In the additive notation of \(M\), we obtain
\[
 [u]=(\phi^r-1)[w]
     =(\phi-1)(1+\phi+\cdots+\phi^{r-1})[w].
\]
This proves the required surjectivity and all higher-cohomology vanishings.
\end{proof}

\subsection{The global trace and its normalization}

\begin{lemma}[The unramified local invariant]
\label{lem:do26-local-invariant}
Let \(R\) be a henselian discrete valuation ring with finite residue
field \(k\), fraction field \(L\), and normalized valuation \(v\).
Let \(L^{\mathrm{nr}}=\operatorname{Frac}(R^{\mathrm{sh}})\) and
\(G_k=\operatorname{Gal}(L^{\mathrm{nr}}/L)\).
Valuation, the inverse Bockstein, and evaluation at arithmetic Frobenius
\(\phi\) give an isomorphism
\[
 H^2(G_k,(L^{\mathrm{nr}})^\times)
 \xrightarrow{v_*}H^2(G_k,\mathbb Z)
 \xrightarrow{\beta^{-1}}H^1(G_k,\mathbb Q/\mathbb Z)
 \xrightarrow{\operatorname{ev}_{\phi}}\mathbb Q/\mathbb Z.
\]
Here \(\beta\) is the boundary for
\(0\to\mathbb Z\to\mathbb Q\to\mathbb Q/\mathbb Z\to0\).
\end{lemma}

\begin{proof}
Put \(U=(R^{\mathrm{sh}})^\times\).  On a finite unramified cyclic
extension, the odd positive cohomology of its units is zero by Hilbert's
theorem 90.  Its proof for units follows from the explicit formula and
valuation adjustment in Lemma~\ref{lem:do26-defect}.  The even positive
cohomology is zero by the cyclic resolution and
Lemma~\ref{lem:do26-norm}.  Taking the limit over finite unramified
extensions gives \(H^q(G_k,U)=0\) for every \(q>0\).
The valuation sequence
\(0\to U\to(L^{\mathrm{nr}})^\times\xrightarrow{v}\mathbb Z\to0\)
therefore makes \(v_*\) an isomorphism in degree two.
Positive-degree cohomology of the trivial module \(\mathbb Q\)
vanishes by averaging over finite quotients.  Thus \(\beta\) is an
isomorphism.  Finally, every continuous character of
\(\widehat{\mathbb Z}\) is determined by its value on \(\phi\),
and every element of \(\mathbb Q/\mathbb Z\) occurs.
\end{proof}

The preceding local calculation isolates the global arithmetic input.
We use the following class-field trace theorem on the ordinary small site:
\begin{equation}
\label{eq:do26-trace-input}
 H^2(X,\mathbb G_m)=0,\qquad
 \operatorname{tr}_F:H^3(X,\mathbb G_m)
                   \xrightarrow{\sim}\mathbb Q/\mathbb Z.
\end{equation}
The field is totally imaginary in this statement.  These are the
class-field computations in Mazur~\cite[\S2,
(2.1)--(2.3), pp.~537--538]{Mazur1973}.  We specify how their normalization enters.

Write \(\operatorname{Br}(L)=H^2(L,\mathbb G_m)\) for the
cohomological Brauer group of a field \(L\).
Let \(j:\operatorname{Spec}F\to X\) be the generic-point morphism and
let \(i_v:\operatorname{Spec}k_v\to X\) run through the finite primes.
The divisor sequence is
\[
 0\longrightarrow\mathbb G_m\longrightarrow j_*\mathbb G_{m,F}
   \xrightarrow{(v)_v}\bigoplus_v i_{v*}\underline{\mathbb Z}
   \longrightarrow0,
\]
where \(v\) takes a local uniformizer to \(+1\).
Local class field theory gives \(R^qj_*\mathbb G_{m,F}=0\) for
\(q>0\).  The resulting cohomology sequence contains
\begin{equation}
\label{eq:do26-residue}
\begin{split}
0\longrightarrow H^2(X,\mathbb G_m)&\longrightarrow\operatorname{Br}(F)
 \xrightarrow{\partial}\bigoplus_v\mathbb Q/\mathbb Z\\
 &\xrightarrow{\delta}H^3(X,\mathbb G_m)
 \longrightarrow H^3(F,\mathbb G_m)\longrightarrow0.
\end{split}
\end{equation}
The copy of \(\mathbb Q/\mathbb Z\) at \(v\) is
\(H^2(k_v,\underline{\mathbb Z})\), identified using arithmetic
Frobenius and the boundary of
\(0\to\mathbb Z\to\mathbb Q\to\mathbb Q/\mathbb Z\to0\).
The invariant on unramified Brauer classes is
the map of Lemma~\ref{lem:do26-local-invariant}.  Its sign can be
computed without a convention for central simple algebras.  For an
unramified cyclic extension of degree \(d\), with arithmetic Frobenius
generator \(\sigma\) and uniformizer \(\varpi\), the cocycle
\[
 \alpha(\sigma^i,\sigma^j)=
          \varpi^{\lfloor(i+j)/d\rfloor},\qquad 0\leq i,j<d,
\]
has invariant \(1/d\).  Indeed, its valuation is the Bockstein of the
character \(\sigma\mapsto1/d\), using the lift
\(\sigma^i\mapsto i/d\) to \(\mathbb Q\).

The substantive class-field assertions are that \(\partial\) agrees
with these local invariants, that
\[
 0\longrightarrow\operatorname{Br}(F)
  \xrightarrow{\partial}\bigoplus_v\mathbb Q/\mathbb Z
  \xrightarrow{\sum_v}\mathbb Q/\mathbb Z\longrightarrow0
\]
is exact, and that \(H^3(F,\mathbb G_m)=0\).
They imply~\eqref{eq:do26-trace-input}, with the normalization
\begin{equation}
\label{eq:do26-trace-normalization}
 \operatorname{tr}_F\bigl(\delta((a_v)_v)\bigr)=\sum_v a_v.
\end{equation}
The connecting homomorphism in~\eqref{eq:do26-residue} uses the usual
cochain differential: lift a cocycle, apply the differential, and take its
class in the kernel.  Formula~\eqref{eq:do26-trace-normalization} fixes
the sign.  The local unit argument does not prove these global class-field
assertions.  They remain the arithmetic hypotheses whenever the trace
theorem is not invoked.

\begin{proposition}[The degree-three orientation]
\label{prop:do26-orientation}
Assume the trace theorem~\eqref{eq:do26-trace-input}, with
normalization~\eqref{eq:do26-trace-normalization}.  Inclusion of roots of
unity induces an isomorphism
\[
 H^3(X,\mu_{n,X})\xrightarrow{\sim}
       H^3(X,\mathbb G_m)[n].
\]
Consequently
\begin{equation}
\label{eq:do26-orientation}
 \operatorname{inv}_{F,n,\zeta}
 =\operatorname{tr}_F\circ H^3(\mu_{n,X}\hookrightarrow\mathbb G_m)
                 \circ H^3(e_\zeta)
 :H^3(X,\underline{\mathbb Z/n})
       \xrightarrow{\sim}\tfrac1n\mathbb Z/\mathbb Z.
\end{equation}
If \(u\in(\mathbb Z/n)^\times\), then
\(\operatorname{inv}_{F,n,\zeta^u}
   =u\operatorname{inv}_{F,n,\zeta}\).
\end{proposition}

\begin{proof}
The definition of \(K_n\) gives two short exact sequences:
\[
 0\to\mu_{n,X}\to\mathbb G_m\xrightarrow{[n]}K_n\to0,
 \qquad
 0\to K_n\xrightarrow{\iota}\mathbb G_m\to Q_n\to0.
\]
Lemma~\ref{lem:do26-defect} and the second sequence show that
\(H^r(\iota)\) is an isomorphism for \(r\geq2\).  In particular,
\(H^2(X,K_n)=0\).  The first sequence therefore identifies
\(H^3(X,\mu_{n,X})\) with the kernel of
\(H^3(X,\mathbb G_m)\to H^3(X,K_n)\).
After composing with \(H^3(\iota)\), this map is multiplication by
\(n\).  The trace identifies its kernel with
\((\mathbb Q/\mathbb Z)[n]=\tfrac1n\mathbb Z/\mathbb Z\).
Lemma~\ref{lem:do26-roots} supplies the remaining isomorphism.
Finally, \(e_{\zeta^u}=e_\zeta\circ[u]\).  Functoriality of
cohomology gives the stated multiplier, with no inversion of \(u\).
\end{proof}

\begin{corollary}[The top-degree duality pairing]
\label{cor:do26-duality}
Under the same trace hypothesis, the Yoneda evaluation pairing
\[
 H^3(X,\underline{\mathbb Z/n})\times
 \operatorname{Hom}_X(\underline{\mathbb Z/n},\mathbb G_m)
 \longrightarrow\mathbb Q/\mathbb Z,
 \qquad(x,f)\longmapsto\operatorname{tr}_F(f_*x),
\]
is perfect.  If \(f(1)=\zeta^a\), its value is
\(a\operatorname{inv}_{F,n,\zeta}(x)\).
\end{corollary}

\begin{proof}
The formula follows from \(f=(\mu_{n,X}\hookrightarrow\mathbb G_m)
\circ e_\zeta\circ[a]\).  Proposition~\ref{prop:do26-orientation}
identifies the first factor with \(\tfrac1n\mathbb Z/\mathbb Z\).
Lemma~\ref{lem:do26-roots} identifies the second with the cyclic group
generated by \(\zeta\).  The resulting pairing \((t,a)\mapsto at\)
identifies each finite group with the full character group of the other.
For \(n=1\), both groups are zero and this remains valid.
\end{proof}

This proves the required degree-three part from the trace theorem.  It
does not assume perfectness of the full Artin--Verdier pairing in every
degree.  Such perfectness is the stronger theorem in
Mazur~\cite[\S2, (2.4), pp.~538--539]{Mazur1973}.

\subsection{The classifying pullback}

Let \(A\) be a finite discrete group.  A right \(A\)-torsor on
\(X_{\mathrm{et}}\) is a sheaf \(P\) with a right action, locally
isomorphic to \(\underline A\) acting on itself.  For a finite constant
group it is represented by a finite \'etale scheme over \(X\).
Fix a geometric generic point \(b:\operatorname{Spec}\overline F\to X\)
and put \(\pi=\pi_1^{\mathrm{et}}(X,b)\).
The finite-cover equivalence identifies finite \'etale schemes with finite
continuous left \(\pi\)-sets
\cite[Theorem 58.6.2]{do26-StacksFundamental}.

Given a continuous homomorphism \(\rho:\pi\to A\), give the set
\(A\) the action \(\gamma\cdot a=\rho(\gamma)a\).
This action commutes with right multiplication.  The finite-cover
equivalence therefore produces a right torsor \(P_\rho\).
Conversely, a point \(p\in P_b\) determines \(\rho\) by
\(\gamma p=p\rho(\gamma)\).  Replacing \(p\) by \(pa\) replaces
\(\rho\) by \(a^{-1}\rho a\).  The equivalence preserves morphisms,
so it also identifies torsor isomorphisms with these changes of point.
Thus
\[
 \{\text{finite \'etale right }A\text{-torsors on }X\}/\cong
 =\operatorname{Hom}_{\mathrm{cont}}(\pi,A)/A.
\]
No surjectivity of \(\rho\) is required.  A torsor need not be connected.

Let \(BA\) denote the classifying topos of left \(A\)-sets.
For a left \(A\)-set \(M\), define the associated sheaf
\[
 T_P(M)=P\times^A\underline M,
 \qquad (pa,m)\sim(p,am),
\]
where the quotient is formed as a sheaf quotient.  Locally on a cover
trivializing \(P\), this is the constant-sheaf functor on the underlying
set of \(M\).  It preserves finite limits and all colimits.  Its right
adjoint sends a sheaf \(\mathcal F\) to
\(\operatorname{Hom}_X(P,\mathcal F)\), with action
\((a\cdot f)(p)=f(pa)\).  Hence \(T_P\) is the inverse image of a
geometric morphism \(f_P:X_{\mathrm{et}}\to BA\).

\begin{proposition}[Cohomology pulled back by a torsor]
\label{prop:do26-pullback}
Let \(M\) be a discrete abelian group with trivial \(A\)-action.
For each \(q\geq0\), a right \(A\)-torsor \(P\) determines a
homomorphism
\[
 f_P^*:H^q(A,M)\longrightarrow H^q(X,\underline M).
\]
It depends only on the isomorphism class of \(P\).
\end{proposition}

\begin{proof}
On abelian group objects the associated-sheaf functor
\[
 T_P:\mathbb Z[A]\text{-}\mathrm{Mod}
       \longrightarrow\mathrm{Ab}(X_{\mathrm{et}})
\]
is exact.  Exactness can be checked after trivializing \(P\), where
it is the exact constant-sheaf functor.  Since the action on \(M\) is
trivial, \([p,m]\mapsto m\) identifies \(T_P(M)\) canonically with
\(\underline M\).  Also \(T_P(\mathbb Z)=\underline{\mathbb Z}\).

Use the Yoneda descriptions
\[
 H^q(A,M)=\operatorname{Ext}^q_{\mathbb Z[A]}(\mathbb Z,M),
 \qquad
 H^q(X,\underline M)=
 \operatorname{Ext}^q_{X}(\underline{\mathbb Z},\underline M).
\]
Apply \(T_P\) to an exact sequence representing a Yoneda extension.
Exactness gives an extension on \(X\) with the stated endpoints.
The functor preserves equivalences of extensions and their Baer sums.
It therefore gives the displayed homomorphism.  In degree zero it sends
an invariant element to its constant section.
A torsor isomorphism \(P\to P'\) induces the natural isomorphism
\([p,m]\mapsto[p',m]\) of these exact functors.  On the trivial
endpoints this is the identity, so the resulting extension classes agree.
\end{proof}

For \(P=P_\rho\), write this homomorphism as \(\rho^*\).
The notation includes the passage from group cohomology to sheaf
cohomology.  It asserts no isomorphism
\(H^q(\pi,M)\cong H^q(X,\underline M)\), and requires no
\(K(\pi,1)\) property for \(X\).

\subsection{Conjugation, cocycles, and the invariant}

We make the conjugation identity explicit.  For trivial coefficients,
inhomogeneous cochains have differential
\[
\begin{split}
 (dc)(g_1,\ldots,g_{q+1})={}&c(g_2,\ldots,g_{q+1})\\
 &+\sum_{i=1}^{q}(-1)^i
       c(g_1,\ldots,g_i g_{i+1},\ldots,g_{q+1})\\
 &+(-1)^{q+1}c(g_1,\ldots,g_q).
\end{split}
\]
For a profinite group and a finite coefficient group, cochains are
continuous maps to the discrete coefficient group.  For finite \(A\),
all maps are continuous.

\begin{lemma}[An explicit conjugation homotopy]
\label{lem:do26-conjugation}
Let \(c\in Z^3(A,M)\), with trivial \(A\)-action on \(M\), and
let \(a\in A\).  Set \(\operatorname{Ad}_a(g)=aga^{-1}\) and
\[
\begin{split}
 h_a(c)(g,h)={}&c(a^{-1},\operatorname{Ad}_a g,
                         \operatorname{Ad}_a h)\\
              &-c(g,a^{-1},\operatorname{Ad}_a h)
                +c(g,h,a^{-1}).
\end{split}
\]
Then \(d h_a(c)=c\circ\operatorname{Ad}_a-c\).
\end{lemma}

\begin{proof}
Use homogeneous bar chains, with differential
\[
 \partial(g_0,\ldots,g_q)=
 \sum_{j=0}^{q}(-1)^j(g_0,\ldots,\widehat g_j,\ldots,g_q).
\]
Here \(q\geq1\), and \(\partial_0=0\).
The chains carry the diagonal left \(A\)-action.  Right translation
\(R(g_0,\ldots,g_q)=(g_0a^{-1},\ldots,g_qa^{-1})\) commutes
with this action.  Define
\[
 K_q(g_0,\ldots,g_q)=
 \sum_{i=0}^{q}(-1)^i
 (g_0,\ldots,g_i,g_i a^{-1},\ldots,g_q a^{-1}).
\]
Deleting a vertex away from the repeated index cancels against the
corresponding term of \(K_{q-1}\partial\).  At the repeated index,
successive terms cancel.  The two endpoints give
\[
 \partial K_q+K_{q-1}\partial=R-\operatorname{id}.
\]
The convention in degree zero has \(K_{-1}=0\).

Let \(\widetilde c\) be the homogeneous cochain associated to \(c\):
\[
 \widetilde c(g_0,\ldots,g_q)=
 c(g_0^{-1}g_1,\ldots,g_{q-1}^{-1}g_q).
\]
It is left invariant.  This invariance identifies \(R^*\widetilde c\)
with the homogeneous cochain of \(c\circ\operatorname{Ad}_a\).
Dualizing the chain identity gives
\(d(\widetilde c K_2)=R^*\widetilde c-\widetilde c\), since
\(d\widetilde c=0\).  Evaluating \(K_2\) at \((1,g,gh)\)
gives exactly the three displayed terms of \(h_a(c)\).
\end{proof}

\begin{theorem}[The finite arithmetic Chern--Simons function]
\label{thm:do26-kim}\label{thm:kim}
Let \(F\), \(X\), \(n\), and \(\zeta\) satisfy the hypotheses of
Lemma~\ref{lem:do26-roots}.  Use the class-field trace
\eqref{eq:do26-trace-input} with normalization
\eqref{eq:do26-trace-normalization}.  For a finite group \(A\) and a
class \(c\in H^3(A,\mathbb Z/n)\) with trivial coefficient action,
\[
 \operatorname{CS}_{c,\zeta}:
 \operatorname{Hom}_{\mathrm{cont}}(\pi_1^{\mathrm{et}}(X,b),A)/A
   \longrightarrow\tfrac1n\mathbb Z/\mathbb Z,
 \qquad [\rho]\longmapsto
       \operatorname{inv}_{F,n,\zeta}(\rho^*c),
\]
is well defined. For \(u\in(\mathbb Z/n)^\times\), replacing
\(\zeta\) by \(\zeta^u\) multiplies the function by \(u\).
\end{theorem}

\begin{proof}
Proposition~\ref{prop:do26-pullback} constructs \(\rho^*c\) in the
required cohomology group.  It is independent of the cocycle representing
\(c\), because the construction is defined on extension classes.
If \(\rho'=a\rho a^{-1}\), the map of fibers \(x\mapsto ax\)
is \(\pi\)-equivariant and commutes with the right \(A\)-action.
It therefore induces an isomorphism \(P_\rho\to P_{\rho'}\).
The two pullbacks agree by Proposition~\ref{prop:do26-pullback}.
At the group-cochain level, Lemma~\ref{lem:do26-conjugation} gives
the same invariance by an explicit coboundary.
Proposition~\ref{prop:do26-orientation} gives the value and its
dependence on the root.  A change of geometric base point gives an
equivalence of finite-cover fiber functors, unique up to inner
conjugation, so it does not alter the function on torsor isomorphism
classes.  This is the finite construction of
Kim~\cite[\S1, pp.~1--3]{Kim2016}.
\end{proof}

For \(A=1\), positive-degree group cohomology vanishes, so this function
is zero.  It is also zero for \(n=1\).  Neither conclusion requires
\(H^3(X,\underline{\mathbb Z/n})\) to vanish for other values of \(n\).

\subsection{The ramification boundary}

The group \(\pi_1^{\mathrm{et}}(X,b)\) classifies extensions
unramified at every finite prime, including primes above \(n\).
For a finite set \(S\) of finite primes, put
\(X_S=\operatorname{Spec}\mathcal O_F[S^{-1}]\).
A representation of \(\pi_1^{\mathrm{et}}(X_S,b)\) gives the
pullback of Proposition~\ref{prop:do26-pullback} on \(X_S\).
An extension to a torsor on \(X\) requires trivial inertia at every
prime of \(S\).  To see this criterion, normalize \(X\) in the
finite generic algebra of the torsor.  Over a discrete valuation ring,
trivial inertia is equivalent to an unramified extension.  Its residue
extensions are separable because finite fields are perfect.  The
normalization is then finite \'etale at those primes and at all primes
of \(X_S\), and hence supplies the extension.

Kim's boundary construction takes \(S\) to contain every prime above
\(n\). For \(v\in S\), let \(F_v\) be the completion of \(F\)
and let \(G_{F_v}=\operatorname{Gal}(\overline F_v/F_v)\). Choose
embeddings that give restriction maps
\(G_{F_v}\to\pi_1^{\mathrm{et}}(X_S,b)\). Given a local
representation \(\rho_v\) and a cocycle \(c\), the set
\[
 \mathcal H_v(\rho_v)=
 \{\beta\in C^2_{\mathrm{cont}}(G_{F_v},\mathbb Z/n):
                      d\beta=\rho_v^*c\}/B^2
\]
where \(B^2=dC^1_{\mathrm{cont}}(G_{F_v},\mathbb Z/n)\),
is either empty or a torsor under
\(H^2(G_{F_v},\mathbb Z/n)\).  Indeed, the difference of two
solutions is a two-cocycle, and quotienting by coboundaries gives a
free transitive action.  Nonemptiness follows from the local vanishing
\(H^3(F_v,\mathbb Z/n)=0\), proved with the invariant
identification in Lemma~\ref{lem:do27-local}. The chosen root
identifies the acting group with \(\tfrac1n\mathbb Z/\mathbb Z\).
A finite product of these torsors can then be pushed out by the
sum of local invariants, as in
Kim~\cite[\S2, pp.~4--6]{Kim2016}.
Those local primitives and their compatibility are additional data.
The construction on \(X\) above provides no scalar evaluation of an
arbitrarily ramified representation on \(X_S\).

\subsection{Infinite places}

The totally imaginary hypothesis excludes real decomposition groups.
To discuss other fields, let \(F\) now be an arbitrary number field,
put \(X=\operatorname{Spec}\mathcal O_F\), and let \(X_\infty\)
be the set of real embeddings and conjugate pairs of complex embeddings.
The notation \(\overline{\operatorname{Spec}\mathcal O_F}\) alone
does not select a cohomology theory.

The Artin--Verdier site \(\widetilde X_{\mathrm{et}}\) has objects
\((Y,M)\), where \(Y\to X\) is separated and \'etale,
\(M\subset Y_\infty\), and a complex point of \(M\) does not map
to a real point of \(X_\infty\). A morphism carries the chosen
infinite subset into the chosen infinite subset. A family covers when
it jointly covers both the schemes and those subsets. Coefficients on
this site mean specified sheaves, together with their infinite-place
specialization data; they are not determined by a topological
compactification \cite[\S1]{Bienenfeld1987}.

At a real place the decomposition group is \(C_2\). For the trivial
module \(\mathbb Z/2\), its two-periodic resolution has zero
cochain differentials, so
\[
 H^q(C_2,\mathbb Z/2)=\mathbb Z/2\qquad(q\geq0).
\]
For odd-order coefficients, the higher groups vanish by averaging.
Thus ordinary real-place cohomology has no finite
\(2\)-cohomological dimension. A three-dimensional duality statement
with real places must specify its framework: ordinary small-site
cohomology with the \(2\)-primary part excluded; modified cohomology
with Tate cohomology at the real places; or derived sheaf cohomology on
an explicitly defined extended Artin--Verdier site, with a constructible
torsion sheaf, its extension, and its infinite-place specialization
maps. The latter two are different constructions. The trace used in
Proposition~\ref{prop:do26-orientation} has the stated totally imaginary
ordinary small-site domain \cite[\S2]{Mazur1973}.

\section{The arithmetic trace on a finite set of local cochains}
\label{sec:do27-trace}

An arithmetic class of degree three can be evaluated through local
classes of degree two. The passage uses a gluing complex and a choice
of primitives. Its sign is forced by the differential. We construct
this passage and compare it with the arithmetic orientation.

Retain the totally imaginary field \(F\), the ordinary small site of
\(X=\operatorname{Spec}\mathcal O_F\), the positive integer \(n\),
and the chosen primitive root \(\zeta\) from
Section~\ref{sec:do26-orientation}. Put \(M=\underline{\mathbb Z/n}\).
Let \(S\) be a nonempty finite set of finite primes containing every
prime above \(n\), and write
\[
 U=X\setminus S,\qquad
 R_v=\mathcal O_{X,v}^{h},\quad D_v=\operatorname{Spec}R_v,
 \quad K_v=\operatorname{Frac}R_v.
\]
The completion of \(K_v\) is \(F_v\), with valuation ring
\(\widehat R_v\). All sites remain ordinary small \'etale sites.
Every scheme map below gives a morphism of these sites by base change.
Cohomological degrees increase under the differential. Coefficient
sheaves and Galois modules are discrete, and products indexed by \(S\)
are finite. The complexes are algebraic complexes of abelian groups.
Sheaf tensor products are over \(\underline{\mathbb Z}\); none is
completed. We write multiplicative abelian groups additively in complexes.

Use the coefficient homomorphism
\[
 \epsilon=\iota e_\zeta:M\longrightarrow\mathbb G_m,
 \qquad a\longmapsto\zeta^a,
\]
where \(\iota:\mu_{n,X}\hookrightarrow\mathbb G_m\).
Lemma~\ref{lem:do26-roots} identifies its kernel-valued factor with
the constant sheaf, including at coefficient primes. The nonzero
power-map cokernel in Lemma~\ref{lem:coefficient-prime-defect} remains
part of the coefficient sequence on \(X\).

\subsection{The arithmetic input and the local groups}

We use the class-field trace and divisor normalization of
\eqref{eq:do26-trace-input}--\eqref{eq:do26-trace-normalization}.
For the generic point of a henselian discrete valuation ring with finite
residue field, the same local class-field input gives
\(R^{>0}j_*\mathbb G_m=0\).
Its stated scope includes henselian, as well as complete, rings
\cite[\S1, (1.1), p.~525]{Mazur1973}. In particular, with
\(\delta_S\) the divisor connecting map for \(S\),
\begin{equation}
\label{eq:do27-input}
 \operatorname{tr}_F\delta_S((t_v)_{v\in S})=\sum_{v\in S}t_v.
\end{equation}
The connecting map lifts a cocycle and applies the differential.
This sign and the arithmetic Frobenius normalization are fixed before
the gluing complex is formed.

\begin{lemma}[The groups needed at a boundary prime]
\label{lem:do27-local}
For \(v\in S\), one has
\[
 H^q(D_v,\mathbb G_m)=0\ (q>0),\qquad
 H^q(D_v,M)=0\ (q\geq2).
\]
The same statements hold for \(\operatorname{Spec}\widehat R_v\).
On either fraction field \(L=K_v,F_v\),
\[
 H^1(L,\mathbb G_m)=0,\quad
 \operatorname{inv}_v:H^2(L,\mathbb G_m)\xrightarrow{\sim}\mathbb Q/\mathbb Z,
 \quad H^q(L,\mathbb G_m)=0\ (q\geq3).
\]
Consequently \(\epsilon\) gives
\[
 H^2(L,M)\xrightarrow{\sim}\tfrac1n\mathbb Z/\mathbb Z,
 \qquad H^3(L,M)=0.
\]
Completion preserves the displayed invariants.
\end{lemma}

\begin{proof}
Let \(R\) be either valuation ring, with residue field \(k\).
A strictly henselian local scheme has exact global sections on its
small \'etale site. Indeed, any \'etale covering has a member with a
point above the closed point. Hensel's lemma supplies a section of
that member. This lifts each section through a sheaf epimorphism.
Finite \'etale neighborhoods of the closed point of
\(\operatorname{Spec}R\) are cofinal among its \'etale neighborhoods.
Descent along their strict henselian limit therefore computes
\[
 H^q(R,\mathcal F)=H^q_{\mathrm{cont}}
       (G_k,\mathcal F(R^{\mathrm{sh}})).
\]
For \(\mathbb G_m\), the module is
\(E=(R^{\mathrm{sh}})^\times\). The unit-norm and cyclic-resolution
arguments of Lemmas~\ref{lem:do26-norm} and
\ref{lem:do26-local-invariant} give \(H^{>0}(G_k,E)=0\).
They apply to either valuation ring and to unramified extensions of
arbitrary degree, including degrees divisible by the residue
characteristic. For \(M\), the cyclic-resolution and inflation
argument in Lemma~\ref{lem:do26-defect} gives
\(H^q(G_k,M)=0\) for \(q\geq2\), also at coefficient primes.

The local direct-image input identifies fraction-field cohomology
with the cohomology of \(G_k\) acting on
\(L^{\mathrm{nr},\times}\). The valuation sequence is
\[
 0\longrightarrow E\longrightarrow L^{\mathrm{nr},\times}
   \xrightarrow{v}\mathbb Z\longrightarrow0.
\]
Valuation is surjective on invariants, because the base field has a
uniformizer. Thus degree one vanishes. In degree two the sequence
identifies the group with \(H^2(G_k,\mathbb Z)\).
The sequence \(0\to\mathbb Z\to\mathbb Q\to\mathbb Q/\mathbb Z\to0\)
then identifies it with \(H^1(G_k,\mathbb Q/\mathbb Z)\).
Positive cohomology of \(\mathbb Q\) vanishes by averaging on finite
quotients. Finite torsion cohomology of \(G_k\) vanishes above degree
one, so the higher groups in the assertion vanish.
Evaluation at arithmetic Frobenius defines \(\operatorname{inv}_v\).

The fraction field has characteristic zero. Its \'etale Kummer sequence
is therefore exact. The degree-two assertion follows from the vanishing
of \(H^1(L,\mathbb G_m)\). Multiplication by \(n\) on
\(H^2(L,\mathbb G_m)=\mathbb Q/\mathbb Z\) is surjective, so the
same sequence gives \(H^3(L,M)=0\). The chosen root identifies
\(M\) with \(\mu_n\) on this field.

For an unramified cyclic extension of degree \(d\), let
\(\sigma\) be arithmetic Frobenius and \(\varpi\) a uniformizer.
The cocycle of Lemma~\ref{lem:do26-local-invariant} is
\begin{equation}
\label{eq:do27-carry}
 \alpha(\sigma^i,\sigma^j)=
 \varpi^{\lfloor(i+j)/d\rfloor},\qquad 0\leq i,j<d,
 \qquad \operatorname{inv}_v[\alpha]=1/d.
\end{equation}
Completion preserves the unramified extension, arithmetic Frobenius,
and the valuation of \(\varpi\). These classes exhaust
\(\mathbb Q/\mathbb Z\), proving compatibility of the invariants.
\end{proof}

\begin{lemma}[Deleting a nonempty finite set]
\label{lem:do27-open}
One has \(H^3(U,\mathbb G_m)=0\). For every
\(a\in H^2(U,\mathbb G_m)\),
\[
 \sum_{v\in S}\operatorname{inv}_v(a|_{K_v})=0.
\]
\end{lemma}

\begin{proof}
For \(j_U:U\hookrightarrow X\), the local direct-image input gives
\(R^{>0}j_{U*}\mathbb G_m=0\). The divisor sequence is
\[
 0\to\mathbb G_{m,X}\to j_{U*}\mathbb G_{m,U}
    \xrightarrow{(v)_{v\in S}}
       \bigoplus_{v\in S}i_{v*}\mathbb Z\to0.
\]
Since \(H^2(k_v,\mathbb Z)=\mathbb Q/\mathbb Z\) and
\(H^3(k_v,\mathbb Z)=0\), its long exact sequence contains
\[
 H^2(U,\mathbb G_m)\longrightarrow
 \bigoplus_{v\in S}\mathbb Q/\mathbb Z
 \xrightarrow{\delta_S}H^3(X,\mathbb G_m)
 \longrightarrow H^3(U,\mathbb G_m)\longrightarrow0.
\]
The first arrow is the local invariant map by the local input and
valuation computation. Naturality from the full divisor sequence gives
\(\operatorname{tr}_F\delta_S=\sum_{v\in S}\).
This sum is surjective because \(S\neq\varnothing\). The assertions
now follow from exactness.
\end{proof}

\subsection{A cochain model on the small site}

For an abelian sheaf \(\mathcal F\), we specify the cochains
to fix the maps and the cup products.
Choose a conservative family of geometric points of \(X\).
Let \(p^*\) be the family of stalk functors and \(p_*\) its right
adjoint. Put \(T=p_*p^*\). The unit and counit define the augmented
cosimplicial resolution
\[
 \mathcal F\longrightarrow T\mathcal F
       \rightrightarrows T^2\mathcal F\rightrightarrows T^3\mathcal F
       \longrightarrow\cdots.
\]
Use its unnormalized cochain complex \(I^q\mathcal F=T^{q+1}\mathcal F\),
with differential \(\sum_{i=0}^{q+1}(-1)^id^i\).
For an \'etale \(V\to X\), put
\(C^q(V,\mathcal F)=\Gamma(V,I^q\mathcal F)\).
Every geometric stalk of the augmented resolution has the extra
codegeneracy induced by the counit. Thus it is exact.
The terms are acyclic for sections over any such \(V\).
Indeed, the sections of \(p_*A\) are products of the groups \(A\)
indexed by lifts of geometric points to \(V\). This functor is exact.
Resolve \(A\) by injectives in the product category of abelian groups.
Since \(p^*\) is exact, \(p_*\) preserves injectives, proving the
acyclicity. This is the Godement construction
\cite{do27-Godement}; the argument here applies to the specified site.

Choose affine separated \'etale neighborhoods \(V\to X\) of \(v\)
with fiber \(\operatorname{Spec}k_v\) at \(v\), and with no fiber
at the other points of \(S\). Define
\[
 C(D_v,\mathcal F)=\varinjlim_V C(V,\mathcal F),\qquad
 C(K_v,\mathcal F)=\varinjlim_V C(V\cap U,\mathcal F).
\]
The limits of these schemes are \(D_v\) and \(\operatorname{Spec}K_v\).
The continuity theorem for \'etale cohomology identifies the cohomology
of these complexes with the indicated small-site groups
\cite[Theorem 59.51.3]{do27-Limits}.
Its hypotheses hold: the schemes are quasi-compact and separated,
and the transition morphisms can be taken affine. Filtered colimits
of abelian groups are exact. The same complexes apply to the restriction
of any sheaf on an open neighborhood containing \(v\).

Write \(C_U=C(U,\mathcal F)\),
\(C_D=\prod_{v\in S}C(D_v,\mathcal F)\), and
\(C_K=\prod_{v\in S}C(K_v,\mathcal F)\).
Restrictions define \(r:C_U\to C_K\) and \(s:C_D\to C_K\).
The finite gluing complex is
\begin{equation}
\label{eq:do27-complex}
 \begin{split}
 A_S^q(\mathcal F)&=C_U^q\oplus C_D^q\oplus C_K^{q-1},\\
 D(u,a,b)&=(du,da,ru-sa-db).
 \end{split}
\end{equation}
Complexes are zero in negative degrees. Equivalently,
\(A_S=\operatorname{Cone}(r-s)[-1]\) with the displayed sign
convention. The map from global cochains is
\[
 J_S:C(X,\mathcal F)\longrightarrow A_S(\mathcal F),
 \qquad z\longmapsto(z|_U,(z|_{D_v})_v,0).
\]

\begin{proposition}[Finite \'etale gluing]
\label{prop:do27-gluing}
The map \(J_S\) is a quasi-isomorphism. For every bilinear sheaf
pairing \(\mathcal F\otimes\mathcal G\to\mathcal H\), it respects
the cohomological cup pairing.
\end{proposition}

\begin{proof}
For every term \(I^q\mathcal F\), restriction from \(V\) to an open
subscheme of \(V\) is surjective. It is a projection between the
products indexed by lifts of geometric points. Consequently
\[
 0\longrightarrow C(X,\mathcal F)
 \longrightarrow C_U\oplus C_D
 \xrightarrow{r-s}C_K\longrightarrow0
\]
is degreewise exact. To check its kernel, represent the finitely many
local germs on neighborhoods \(V_v\). Shrink them until equality
with the given section on \(U\) holds on every \(V_v\cap U\).
The family \(U\amalg\coprod_vV_v\to X\) is an \'etale cover.
An overlap for distinct \(v\)'s lies over \(U\). On a self-overlap,
the complement of the diagonal also lies over \(U\): the chosen
closed fiber is a single residue-field isomorphism. Thus the sections
agree on every overlap and descend uniquely to \(X\).
Surjectivity on the right follows by lifting separately from
\(V_v\cap U\) to \(V_v\). Finiteness of \(S\) permits these
choices and refinements simultaneously.

For completeness, let \((u,a,b)\) be a cocycle. Lift \(b\) to
\(\ell\in(C_U\oplus C_D)^{q-1}\) with \((r-s)\ell=b\).
Subtracting \(D(\ell,0)\) leaves \(((u,a)-d\ell,0)\), a
cocycle in the kernel. This proves surjectivity of \(H(J_S)\).
If a kernel cocycle is \(D(\ell,b)\), lift \(b\) one degree lower
to \(m\). Then \(\ell-dm\) lies in the kernel and has the same
differential. This proves injectivity.

The functor \(T\) has a natural pairing: pair sections point by
point. It induces pairings on all cosimplicial degrees. For cochains
\(x\) and \(y\) of degrees \(p\) and \(q\), use
\[
 x\smile y=
 (d^{p+q}\cdots d^{p+1}x)\,\cdot\,((d^0)^py).
\]
The dot denotes the given coefficient pairing. This
Alexander--Whitney formula has
\(d(x\smile y)=dx\smile y+(-1)^px\smile dy\).
The cosimplicial identities cancel the interior cofaces in pairs;
the two remaining groups are exactly the two displayed summands.
It represents the derived cup pairing because the augmentation
restricts it to the coefficient pairing in degree zero.
Restrictions commute with every coface and with the pairing.

For \(x=(u,a,b)\in A_S^p(\mathcal F)\) and
\(y=(u',a',b')\in A_S^q(\mathcal G)\), set
\begin{equation}
\label{eq:do27-cup}
 x\smile y=
 \bigl(u\smile u',\ a\smile a',\
        b\smile ru'+(-1)^p sa\smile b'\bigr).
\end{equation}
Its last component has degree \(p+q-1\). The identity
\[
 ru\smile ru'-sa\smile sa'
 =(ru-sa)\smile ru'+sa\smile(ru'-sa')
\]
and the preceding Leibniz rule give
\(D(x\smile y)=Dx\smile y+(-1)^px\smile Dy\).
For associative coefficient pairings, the three terms in a triple
product are
\[
 b\smile ru'\smile ru'',\quad
 (-1)^p sa\smile b'\smile ru'',\quad
 (-1)^{p+q}sa\smile sa'\smile b''.
\]
They are the same under both associations. Finally \(J_S\) has last
component zero, so it commutes with this cup formula on cochains.
\end{proof}

For \(M\), multiplication in \(\mathbb Z/n\) gives an associative
cup product. For \(\mathbb G_m\) we use only specified bilinear
pairings, such as evaluation
\(M\otimes\underline{\operatorname{Hom}}(M,\mathbb G_m)\to\mathbb G_m\).
Multiplication of units is not a bilinear map
\(\mathbb G_m\otimes\mathbb G_m\to\mathbb G_m\).

\subsection{The degree-three trace and its sign}

\begin{theorem}[Trace by local primitives]
\label{thm:do27-trace}
Let \((u,a,b)\in Z^3 A_S(\mathbb G_m)\). Choose
\[
 h\in C^2(U,\mathbb G_m),\quad
 k_v\in C^2(D_v,\mathbb G_m),\qquad dh=u,\quad dk_v=a_v.
\]
Then
\begin{equation}
\label{eq:do27-trace}
 T_S[(u,a,b)]
   =\sum_{v\in S}\operatorname{inv}_v
           [\,r_vh-s_vk_v-b_v\,]
\end{equation}
is independent of the choices and the representative. It defines a
homomorphism \(H^3 A_S(\mathbb G_m)\to\mathbb Q/\mathbb Z\), and
\[
 T_S\circ H^3(J_S)=\operatorname{tr}_F.
\]
\end{theorem}

\begin{proof}
The primitives exist by Lemmas~\ref{lem:do27-local} and
\ref{lem:do27-open}. The cone cocycle equation gives
\(db=ru-sa\), so every cochain inside brackets is closed.
Changing \(k_v\) changes that bracket by the restriction of a
closed degree-two cochain on \(D_v\). Its class is zero by
Lemma~\ref{lem:do27-local}. Changing \(h\) adds restrictions of a
degree-two global cocycle on \(U\). Its local invariants sum to zero
by Lemma~\ref{lem:do27-open}.

If \((u,a,b)\) changes by \(D(p,q,t)\), use \(h+p\) and \(k+q\).
The bracket changes by \(dt\). Thus the value depends only on the
cohomology class. Add primitives to check additivity.

We now compare connecting maps, including their signs. Put
\(B=C_U\oplus C_D\), \(L=C_K\), and
\(C_X=C(X,\mathbb G_m)\). In the degreewise exact
sequence \(0\to C_X\to B\xrightarrow{r-s}L\to0\), lift a
closed \(t\in L^2\) to \(\ell\in B^2\).
The usual connecting class is represented by \(d\ell\in C_X^3\).
But
\[
 D(\ell,0)=(d\ell,t),\qquad
 [(0,0,t)]=-H^3(J_S)[\delta_{\mathrm{glue}}t].
\]
Hence \((0,0,-t)\) represents the positive connecting class.

Under valuation, \(H^2(K_v,\mathbb G_m)\) identifies with
\(H^2(k_v,\mathbb Z)\). We verify that the gluing boundary gives
the divisor boundary with its positive sign. Resolve the finite divisor
sequence by the same exact Godement functors. Taking sections gives
a degreewise exact sequence of complexes: every term of the
resolution is a product of stalk groups. These complexes also
compute the open and local generic groups, by the higher-direct-image
vanishing in the arithmetic input. Common double resolutions compare
them with the complexes already chosen.

Represent a residue class by a cocycle \(t\) of degree two, and
lift it globally to a meromorphic cochain \(\ell\).
Then \(d\ell=z\) is a \(\mathbb G_m\)-cocycle of degree three,
representing \(\delta_S[t]\). Its restriction to \(U\) supplies
\(h\), with \(dh=z_U\). At \(v\), choose a closed meromorphic
cochain \(\beta_v\) with valuation exactly \(t_v\).
Such a cochain exists: valuation is an isomorphism on degree-two
cohomology, and a discrepancy which is a coboundary can be removed
by lifting its primitive through the degreewise surjective valuation
map. The difference \(\ell|_{D_v}-\beta_v\) has valuation zero,
so it is a unit cochain \(k_v\). It satisfies \(dk_v=z_{D_v}\).
In the local generic complex,
\[
 r_vh-s_vk_v=\beta_v.
\]
Thus \((h,k)\) is a lift of \(\beta\) through \(r-s\), and
its differential is \(J_Sz\). This proves that
\(\delta_{\mathrm{glue}}[\beta]=\delta_S[t]\).
The reasoning takes place in degreewise exact resolutions and so
also proves the statement after their comparison maps. By~\eqref{eq:do27-input},
\[
 \operatorname{tr}_F\delta_{\mathrm{glue}}(\beta)
     =\sum_v\operatorname{inv}_v(\beta_v).
\]

Subtract \(D(h,k,0)\) from the original cocycle. The result is
\((0,0,b-rh+sk)\). The connecting-map computation now gives
exactly~\eqref{eq:do27-trace} and the asserted comparison.
\end{proof}

The sign is already visible over \(\mathbb F_3\). Let the complex
\(B\) have \(B^2=B^3=\mathbb F_3\) and differential the identity.
Let \(L\) be concentrated in degree two, with \(L^2=\mathbb F_3\),
and let \(r:B\to L\) be the identity in degree two. Its kernel
\(K\) is concentrated in degree three. Then
\[
 H^3(\operatorname{Cone}(r)[-1])
   =\mathbb F_3^2/\{(t,t):t\in\mathbb F_3\}.
\]
The positive connecting class of \(1\in L^2\) is represented by
\((1,0)\), so the normalized functional is \([(x,y)]\mapsto x-y\).
It sends the bare cone class \((0,1)\) to \(-1=2\pmod3\).
Reversing this sign sends that same class to \(1\), which contradicts
the positive connecting normalization.

For a single prime and a local class of invariant \(1/3\), the
cone class \((0,0,\alpha)\) has trace \(-1/3\).
Taking \(n=3\), \(F=\mathbb Q(\zeta_3)\), and the prime
above \(3\) realizes this case by~\eqref{eq:do27-carry}. Kummer theory gives a
class \(\gamma\in H^2(K_v,M)\) mapping to \(\alpha\), so the
finite-coefficient cone class also has trace \(-1/3\).
The opposite sign would disagree with the divisor orientation.
This case also retains a coefficient prime in the integer scheme.

\subsection{Passage to completed local fields}

We next replace the henselian fields by the local fields in the
finite adelic product \(\prod_{v\in S}F_v\).
This is a change of local cochain complexes, rather than a completion
of the coefficient group.

\begin{lemma}[Finite coefficients and completion]
\label{lem:do27-completion}
For \(K=K_v\), base change to \(\widehat K=F_v\) gives an
equivalence of finite \'etale field algebras. Thus the induced map of
absolute Galois groups is an isomorphism, after choosing compatible
separable closures. It preserves inertia and arithmetic Frobenius
on the unramified quotient.
\end{lemma}

\begin{proof}
We give the approximation argument. A henselian valuation extends
uniquely to an algebraic closure, so its automorphisms are isometries.
For a separable element \(\alpha\), put
\(c=\min_{\alpha'\neq\alpha}|\alpha-\alpha'|>0\), where
\(\alpha'\) runs through its conjugates. If \(|\alpha-\beta|<c\),
every automorphism fixing \(\beta\) fixes \(\alpha\).
Indeed, an automorphism \(\sigma\) fixing \(\beta\) satisfies
\[
 |\sigma\alpha-\alpha|
 \leq\max\{|\sigma\alpha-\beta|,|\beta-\alpha|\}
 =|\alpha-\beta|<c.
\]
It follows that \(K(\alpha)\subset K(\beta)\) when both elements
are algebraic over the henselian field under consideration.
If \(\alpha\in\widehat K\) is separable algebraic over \(K\),
choose \(\beta\in K\) sufficiently close to it. Apply the same
argument in the algebraic extension containing \(\alpha\).
It gives \(\alpha\in K\). Thus \(K\) is separably closed in
\(\widehat K\).

For a finite Galois extension \(E/K\), this implies linear
disjointness from \(\widehat K\). In fact
\(E\cap\widehat K=K\), and restriction identifies
\(\operatorname{Gal}(E\widehat K/\widehat K)\) with the subgroup
of \(\operatorname{Gal}(E/K)\) whose fixed field is that
intersection. The subgroup is the full group. Applying this to
Galois closures preserves degrees and all embeddings of finite
separable extensions.

Conversely, take \(\widehat E=\widehat K(\alpha)\), a finite
separable extension, and its monic minimal polynomial \(f\).
Scale \(\alpha\) to be integral. Approximate the coefficients of
\(f\) by coefficients in \(K\), obtaining a monic polynomial
\(g\) of the same degree. They can be chosen so that
\(|g(\alpha)|<|g'(\alpha)|^2\), and so that the resulting
Hensel root \(\beta\in\widehat E\) is closer to \(\alpha\)
than all its other conjugates. The estimate
\(|\beta-\alpha|\leq|g(\alpha)/g'(\alpha)|\) makes this explicit.
The preceding isometry argument over \(\widehat K\) gives
\(\widehat K(\alpha)\subset\widehat K(\beta)\).
Both fields lie in \(\widehat E\), so equality holds.
The degree of \(K(\beta)\) is at most \(\deg g\) and at least
\([\widehat E:\widehat K]\). These are the same number.
Thus \(K(\beta)\otimes_K\widehat K=\widehat E\).
Products and the embedding statement prove the equivalence.
Approximating coordinates in a finite field basis shows that each
finite extension is dense in its completed extension. A nonzero element
can therefore be approximated with strictly smaller error than its
absolute value. The approximant has the same valuation and, for units,
the same residue. Thus value groups and residue fields are unchanged.
The unramified subcategories and their Frobenius automorphisms agree.
\end{proof}

For clarity, the cochain map under this equivalence is not implicit.
The small \'etale topos of a field \(L\) is the topos of discrete
continuous \(G_L\)-sets. Finite \'etale objects give its finite
continuous sets, and sheaf descent extends this description to all
objects \cite[Lemmas 59.59.1--59.59.2]{do27-Field}.
Its standard cochains with a discrete module \(N\) are
\(C^q_{\mathrm{cont}}(G_L,N)\), with differential
\[
\begin{split}
 (df)(g_1,\ldots,g_{q+1})={}&g_1f(g_2,\ldots,g_{q+1})\\
 &+\sum_{i=1}^q(-1)^if(g_1,\ldots,g_ig_{i+1},\ldots,g_{q+1})\\
 &+(-1)^{q+1}f(g_1,\ldots,g_q).
\end{split}
\]
For a group map \(\theta:G_{F_v}\to G_{K_v}\) and a compatible
module map \(\lambda:N\to N'\), completion sends
\[
 f\longmapsto\bigl[(g_1,\ldots,g_q)\mapsto
       \lambda f(\theta g_1,\ldots,\theta g_q)\bigr].
\]
For \(M\), \(\lambda\) is the identity. For \(\mathbb G_m\),
it is the chosen inclusion of separable closures on their units.
For the latter coefficients, the map in degree zero is not an
isomorphism. The local invariant comparison concerns degree two.
The group cup formula is
\[
 (f\smile g)(g_1,\ldots,g_{p+q})=
 f(g_1,\ldots,g_p)\cdot
 (g_1\cdots g_p)g(g_{p+1},\ldots,g_{p+q}),
\]
for a specified equivariant bilinear pairing. Substitution shows
that the completion map respects this formula and the differential.
The bar resolution and the Godement resolution are acyclic resolutions
of the same coefficient object. Their common double resolution
provides comparison quasi-isomorphisms compatible with restriction
and the derived cup pairing.

Here is also a direct description before passing to the derived
category. For a map of the local schemes, choose geometric points
with compatible images. The adjunction maps give
\(f^{-1}T\mathcal F\to T'f^{-1}\mathcal F\), by sending a section
to each of its target stalks. Iterate this map in every cosimplicial
degree and apply the coefficient map. Taking sections gives the
cochain maps. The unit and counit identities make them commute with
cofaces, pairings, and the restriction square. Apply these maps to
the filtered neighborhood complexes above.

Use these comparison maps to define
\[
 \widehat A_S(\mathcal F)=
 \operatorname{Cone}\left(
 C(U,\mathcal F)\oplus
 \prod_{v\in S}C(\widehat R_v,\mathcal F)
 \xrightarrow{r-s}\prod_{v\in S}C(F_v,\mathcal F)
 \right)[-1].
\]
There is a cochain map \(\kappa_S:A_S\to\widehat A_S\).
For \(M\), it is a quasi-isomorphism: the field terms agree by
Lemma~\ref{lem:do27-completion}, and the ring terms have the same
residue-field cohomology. The cone long exact sequence proves the
claim. For \(\mathbb G_m\), it induces an isomorphism in degree
three. Indeed, both degree-three groups are the cokernel of
\[
 H^2(U,\mathbb G_m)\longrightarrow
          \prod_{v\in S}H^2(K_v,\mathbb G_m)
 \quad\text{or}\quad
 H^2(U,\mathbb G_m)\longrightarrow
          \prod_{v\in S}H^2(F_v,\mathbb G_m).
\]
Lemmas~\ref{lem:do27-local} and \ref{lem:do27-open} justify this
cokernel description. Local invariant compatibility identifies the
arrows. Formula~\eqref{eq:do27-trace} defines \(\widehat T_S\)
on the completed complex and gives
\[
 \widehat T_S H^3(\kappa_S)=T_S.
\]

\subsection{Orientation, boundary choices, and finite torsors}

Write the already constructed arithmetic orientation as
\[
 \operatorname{or}_{F,n,\zeta}
 :=\operatorname{inv}_{F,n,\zeta}
 =\operatorname{tr}_F\circ H^3(\epsilon):
 H^3(X,M)\xrightarrow{\sim}\tfrac1n\mathbb Z/\mathbb Z.
\]
Proposition~\ref{prop:do26-orientation} proves this isomorphism on
\(X\), including its coefficient primes.
For \(n=1\), all constructions with \(M\) have zero cohomology.
The comparison just proved is the commutative diagram
\begin{equation}
\label{eq:do27-comparison}
\begin{CD}
 H^3(X,M) @>{H^3(\epsilon)}>> H^3(X,\mathbb G_m)\\
 @V{H^3(\kappa_SJ_S)}V{\sim}V @VV{H^3(\kappa_SJ_S)}V\\
 H^3\widehat A_S(M) @>{H^3(\epsilon)}>>
 H^3\widehat A_S(\mathbb G_m) @>{\widehat T_S}>> \mathbb Q/\mathbb Z.
\end{CD}
\end{equation}
Composing either path with the trace gives
\(\operatorname{or}_{F,n,\zeta}\).
For a sheaf homomorphism \(f:M\to\mathbb G_m\), the same diagram
is the cup evaluation of degree three against degree zero.
Proposition~\ref{prop:do27-gluing} proves this compatibility for all
specified cup pairings, not only for \(f=\epsilon\).

To state the ramified boundary precisely, put
\[
 C_{c,S}(U,\mathcal F)=
 \operatorname{Cone}\left(C(U,\mathcal F)
         \xrightarrow r\prod_{v\in S}C(F_v,\mathcal F)\right)[-1],
 \quad D(u,b)=(du,ru-db).
\]
The subscript denotes this displayed finite cone. For \(\mathcal F=M\)
or \(\mathbb G_m\), its degree-three cohomology has a canonical
map to \(H^3(X,\mathcal F)\), obtained from
\((u,b)\mapsto(u,0,b)\) and the gluing comparison.
For \(\mathcal F=M\), this map is an isomorphism: the quotient
complex consists of the local ring terms, whose degrees two and
three vanish. For \(\mathbb G_m\), the same conclusion follows
from their positive-degree vanishing and the degree-three comparison.
On \(M\)-valued cochains the compact trace is
\begin{equation}
\label{eq:do27-compact-trace}
 \operatorname{Tr}_{c,S}[(u,b)]
  =\sum_{v\in S}\operatorname{inv}_v[\,r_vh-\epsilon b_v\,],
 \qquad dh=\epsilon u.
\end{equation}
In particular the value is in \(\tfrac1n\mathbb Z/\mathbb Z\).
This follows by its comparison with \(H^3(X,M)\), rather than an
assertion that each auxiliary \(\mathbb G_m\)-cochain is killed by \(n\).

\subsection{The positive boundary torsor and its dual}
\label{sec:do28-dual}

Put \(\Lambda=\tfrac1n\mathbb Z/\mathbb Z\) and
\[
 B_S=\prod_{v\in S}H^2(F_v,M),\qquad
 \sigma_\zeta:B_S\longrightarrow\Lambda,
 \quad(\eta_v)\longmapsto
       \sum_v\operatorname{inv}_v(\epsilon_*\eta_v).
\]
For \(u\in Z^3C(U,M)\), define
\[
 \mathcal H_S(u)=
 \{b\in\prod_{v\in S}C^2(F_v,M):db=ru\}/
                \prod_{v\in S}dC^1(F_v,M).
\]
It is a nonempty \(B_S\)-torsor by Lemma~\ref{lem:do27-local}.
For two points \(b,b'\), the notation \(b'-b\) means their unique
difference in \(B_S\). All torsor actions below are additive right
actions. The finite torsors carry the discrete topology.

The positive and negative pushouts are the \(\Lambda\)-torsors
\begin{equation}
\label{eq:do28-pushouts}
 L_S^\pm(u)=(\mathcal H_S(u)\times\Lambda)/\sim_\pm,
 \qquad (b+\eta,t)\sim_\pm(b,t\pm\sigma_\zeta\eta).
\end{equation}
Their actions are \([b,t]_\pm+a=[b,t+a]_\pm\).
For any fixed \(b\), the map \(t\mapsto[b,t]_\pm\) is a bijection.
Indeed, freeness of \(\mathcal H_S(u)\) gives a unique difference
between any two representatives. Thus the displayed actions are free
and transitive. The natural map \(b\mapsto[b,0]_+\) has positive
translation \(\sigma_\zeta\).

After the field-cochain identification, \(\mathcal H_S\) is Kim's
primitive torsor, and \(L_S^+\) is its positive sum pushout
\cite[\S2, p.~5, the definitions of \(H(\rho_S)\) and
\(L(\rho_S)\)]{Kim2016}. A scalar functional with negative translation
instead defines an equivariant map on \(L_S^-\).

For a \(\Lambda\)-torsor \(L\), define its opposite \(L^{\mathrm{op}}\)
to have the same underlying set and action \(\ell*a=\ell-a\).
Define its affine dual by
\[
 L^\vee=\{f:L\to\Lambda:f(\ell+a)=f(\ell)+a\},
 \qquad(f+a)(\ell)=f(\ell)+a.
\]
This is a \(\Lambda\)-torsor. Each \(\ell_0\in L\) determines the
map \(f_{\ell_0}(\ell)=\ell-\ell_0\). Every element of \(L^\vee\)
has this form, since every equivariant map has exactly one zero.

\begin{proposition}[Opposite torsor and dual evaluation]
\label{prop:do28-dual}
There are canonical equivariant isomorphisms
\[
 L_S^-(u)\xrightarrow{\iota_S}(L_S^+(u))^{\mathrm{op}}
       \xrightarrow{\sim}(L_S^+(u))^\vee.
\]
They are given by
\begin{equation}
\label{eq:do28-dual-map}
 \iota_S[b,t]_-=[b,-t]_+,
 \qquad
 \Psi_S([b,t]_-)([b',s]_+)
       =s+t+\sigma_\zeta(b'-b).
\end{equation}
In particular
\(E_S(\ell,m)=\Psi_S(m)(\ell)\) is a canonical evaluation pairing.
It satisfies
\(E_S(\ell+a,m+a')=E_S(\ell,m)+a+a'\).
\end{proposition}

\begin{proof}
The negative pushout relation maps under \(\iota_S\) to
\[
 [b+\eta,-t]_+=[b,-t+\sigma_\zeta\eta]_+.
\]
This is the image of \([b,t-\sigma_\zeta\eta]_-\), so the map
is well defined. Adding \(a\) on its source subtracts \(a\) in the
positive torsor, which is the opposite action.
The map \(L^{\mathrm{op}}\to L^\vee\), \(\ell_0\mapsto f_{\ell_0}\),
is equivariant because \(f_{\ell_0-a}=f_{\ell_0}+a\).
For \(\ell_0=[b,-t]_+\), its value at \([b',s]_+\) is the
right side of~\eqref{eq:do28-dual-map}.
This proves the formula and both isomorphisms. The evaluation identity
follows directly from addition in the two coordinates.
\end{proof}

This also gives a literal duality of complex lines. Let
\(\chi(a)=\exp(2\pi ia)\) and define
\(\mathcal L(L)=(L\times\mathbb C)/((\ell+a,z)\sim(\ell,\chi(a)z))\),
with its vector-space structure and usual complex topology defined in
any torsor coordinate.
The formula
\[
 ([\ell,z],[m,w])\longmapsto zw\chi(E_S(\ell,m))
\]
is independent of representatives and is a nonzero bilinear pairing.
It identifies \(\mathcal L(L_S^-)\) with
\(\mathcal L(L_S^+)^*\). These are finite-dimensional associated
lines, with no determinant or analytic trace asserted.

\begin{lemma}[A global primitive with finite coefficients]
\label{lem:do28-global-primitive}
Under the stated hypotheses, \(H^3(U,M)=0\).
For every \(a\in H^2(U,M)\), one has
\(\sigma_\zeta(r a)=0\).
\end{lemma}

\begin{proof}
The sequence in Lemma~\ref{lem:do27-open}, together with
\(H^2(X,\mathbb G_m)=0\), identifies
\[
 H^2(U,\mathbb G_m)
 \xrightarrow{\sim}
 \ker\left(\bigoplus_{v\in S}\mathbb Q/\mathbb Z
                         \xrightarrow{\sum_v}\mathbb Q/\mathbb Z\right).
\]
Choose \(v_0\in S\). Projection to the other coordinates identifies
this kernel with \((\mathbb Q/\mathbb Z)^{|S|-1}\), the omitted
coordinate being the negative sum. Hence multiplication by \(n\)
on \(H^2(U,\mathbb G_m)\) is surjective.
Every prime above \(n\) lies in \(S\). The \'etale Kummer sequence
is therefore exact on \(U\), and \(e_\zeta:M\to\mu_n\) is an
isomorphism there. Its long exact sequence, with
\(H^3(U,\mathbb G_m)=0\), gives \(H^3(U,M)=0\).
Finally apply \(\epsilon_*\) to \(a\), and use the zero-sum
statement of Lemma~\ref{lem:do27-open}.
\end{proof}

\begin{theorem}[The global boundary point and the arithmetic trace]
\label{thm:do28-state}
For \(u\in Z^3C(U,M)\), choose \(\beta\in C^2(U,M)\) with
\(d\beta=u\). There is a well-defined point
\begin{equation}
\label{eq:do28-state}
 \lambda_S(u)=[r\beta,0]_+\in L_S^+(u),
\end{equation}
independent of \(\beta\). For every \(b\in\mathcal H_S(u)\),
\begin{equation}
\label{eq:do28-state-trace}
 \tau_u(b):=\operatorname{Tr}_{c,S}[(u,b)]
   =\sigma_\zeta(r\beta-b)
   =\lambda_S(u)-[b,0]_+.
\end{equation}
Consequently the map
\[
 \Theta_u:L_S^-(u)\longrightarrow\Lambda,
 \qquad [b,t]_-\longmapsto\tau_u(b)+t,
\]
is an equivariant isomorphism and satisfies
\(\Theta_u(m)=E_S(\lambda_S(u),m)\).
Its zero is the dual point \(\lambda_S(u)^\vee=\Psi_S^{-1}(f_{\lambda_S(u)})\).
\end{theorem}

\begin{proof}
Lemma~\ref{lem:do28-global-primitive} gives \(\beta\).
Two choices differ by a global two-cocycle \(a\), whose local
invariants sum to zero. The positive pushout relation gives
\([r\beta+r a,0]_+=[r\beta,\sigma_\zeta(r a)]_+=[r\beta,0]_+\).
Thus \(\lambda_S(u)\) is independent of that choice.

Take \(h=\epsilon\beta\) in~\eqref{eq:do27-compact-trace}.
It gives the first equality in~\eqref{eq:do28-state-trace}.
The positive pushout relation gives
\([r\beta,0]_+=[b,\sigma_\zeta(r\beta-b)]_+\), which proves
the torsor difference formula. In particular
\(\tau_u(b+\eta)=\tau_u(b)-\sigma_\zeta\eta\).
This is exactly the relation needed to define \(\Theta_u\).
It is equivariant and therefore is an isomorphism of torsors.
Substitute \(\lambda_S(u)=[b,\tau_u(b)]_+\) in the evaluation
formula to obtain \(\Theta_u(m)=E_S(\lambda_S(u),m)\).
In coordinates its zero is \([b,-\tau_u(b)]_-\), independent of
\(b\). Under \(\Psi_S\) this is the function with zero
\(\lambda_S(u)\), as asserted.
\end{proof}

Kim's global boundary point is \([r\beta,0]_+\)
\cite[\S2, p.~6, the displayed definition of \(\operatorname{CS}_c\)]{Kim2016}.
Theorem~\ref{thm:do28-state} identifies the cone trace with its
difference from the chosen boundary point. Kim also prints the scalar
coordinate \(\sum_v\operatorname{inv}_v(b_v-r_v\beta)\)
\cite[\S5, p.~12]{Kim2016}. For the same cochains and local
invariants that displayed coordinate equals \(-\tau_u(b)\).
Negation is the comparison between these two scalar coordinates.
The divisor normalization remains~\eqref{eq:do27-input}.

\subsection{Gluing and orientation compatibility}
\label{sec:do28-gluing}

For two \(\Lambda\)-torsors define
\(L_1\otimes_\Lambda L_2=(L_1\times L_2)/((\ell_1+a,\ell_2)\sim
(\ell_1,\ell_2+a))\), with addition in either factor.
Write \(\sigma_T\) for the sum of local invariant maps over a
finite subset \(T\), with the root \(\zeta\) fixed.
A finite partition \(S=S_1\amalg S_2\) gives canonical isomorphisms
\begin{equation}
\label{eq:do28-block-gluing}
 L_S^\pm(u)\cong L_{S_1}^\pm(u)\otimes_\Lambda L_{S_2}^\pm(u),
 \qquad
 ([b_1,t_1]_\pm,[b_2,t_2]_\pm)\longmapsto
       [(b_1,b_2),t_1+t_2]_\pm,
\end{equation}
where the arrow displays the inverse to the indicated isomorphism.
The factors mean boundary torsors for restrictions of this same \(u\).
Indeed, \(B_S=B_{S_1}\times B_{S_2}\) and
\(\sigma_S=\sigma_{S_1}+\sigma_{S_2}\). Substitution proves that
the formula respects both pushout relations and the contracted
product relation. It is equivariant and hence bijective.
The evaluation formula gives
\[
 E_S(\ell_1\otimes\ell_2,m_1\otimes m_2)
   =E_{S_1}(\ell_1,m_1)+E_{S_2}(\ell_2,m_2).
\]
It follows that this gluing commutes with the opposite and dual
identifications. It is associative under further finite partitions,
because addition of the displayed coordinates is associative.
For an empty block, use the torsor \(\Lambda\) with point zero as
the tensor unit.

A chosen global primitive \(\beta\) gives block points
\(\lambda_{S_i}(\beta)=[r_{S_i}\beta,0]_+\).
Their tensor is \(\lambda_S(u)\).
Changing \(\beta\) by a global cocycle \(a\) translates those
block points by \(\sigma_{S_i}(r a)\). Their sum is zero by
Lemma~\ref{lem:do28-global-primitive}. Thus the glued point is
independent of \(\beta\), even when the individual block points
are not.

There is also a direct gluing formula for two degree-three cocycles.
Suppose \(u_1,u_2\in Z^3C(U,M)\) satisfy \(r u_1=r u_2\) as
cochains. Their boundary torsors have the same definition, so write
\(L^+\) for this common torsor. Then
\begin{equation}
\label{eq:do28-difference-gluing}
 E_S(\lambda_S(u_1),\lambda_S(u_2)^\vee)
   =\operatorname{Tr}_{c,S}[(u_1-u_2,0)].
\end{equation}
To prove the formula, take a common boundary primitive \(b\).
The left side is \(\lambda_S(u_1)-\lambda_S(u_2)\), which equals
\(\tau_{u_1}(b)-\tau_{u_2}(b)\) by
Theorem~\ref{thm:do28-state}. Subtraction of the two relative
cocycles gives \((u_1-u_2,0)\), proving the equality.
Through the finite \'etale gluing isomorphism, the right side is the
arithmetic orientation of the corresponding class in \(H^3(X,M)\).
Thus dual evaluation has exactly the difference sign required by
that orientation.

Let \(A\) be a finite group, and let \(P_X\) be a finite \'etale
right \(A\)-torsor on \(X\). Fix \(c\in H^3(A,\mathbb Z/n)\)
with trivial coefficient action. The exact associated-sheaf functor
\(N\mapsto P_X\times^A\underline N\) acts on Yoneda extensions
and gives \(f_{P_X}^*c\in H^3(X,M)\).
Choose a cocycle \(z\) representing this class and put \(u=z_U\).
At each deleted prime choose \(a_v\) with
\(da_v=z|_{\widehat R_v}\). Local ring cohomology makes the
class of \(a_v\) modulo coboundaries unique. Put \(b^0=sa\).
Subtracting \(D(0,a,0)\) from \((z_U,z_D,0)\) gives
\((z_U,0,b^0)\). The established trace comparison therefore gives
\begin{equation}
\label{eq:do28-unramified-orientation}
 \operatorname{or}_{F,n,\zeta}(f_{P_X}^*c)
   =\lambda_S(u)-[b^0,0]_+
   =E_S(\lambda_S(u),[b^0,0]_-).
\end{equation}
Notice that \([b^0,0]_-\) is the dual of the positive boundary
point \([b^0,0]_+\). It need not be the dual of the global point
\(\lambda_S(u)\). This distinction permits a nonzero orientation.
For a torsor on \(U\) that does not extend to \(X\), the global
point \(\lambda_S(u)\) still exists. The scalar requires a boundary
point, and its translation remains \(-\sigma_\zeta\).

\subsection{Functoriality of the dual comparison}
\label{sec:do28-functoriality}

A boundary cocycle \(v\in\prod_S Z^3C(F_v,M)\) has its primitive
torsor \(\mathcal H(v)\). A boundary two-cochain \(q\) satisfying
\(dq=v'-v\) gives a transport
\[
 F_q:\mathcal H(v)\longrightarrow\mathcal H(v'),
 \quad[b]\longmapsto[b+q],\qquad
 F_q^\pm[b,t]_\pm=[b+q,t]_\pm.
\]
Replacing \(q\) by a coboundary leaves these maps unchanged.
Composition satisfies \(F_{q'}F_q=F_{q+q'}\).
Differences of two primitive points are unchanged by translation.
Thus~\eqref{eq:do28-dual-map} gives the exact identities
\begin{equation}
\label{eq:do28-functor-dual}
 \iota F_q^-=(F_q^+)^{\mathrm{op}}\iota,
 \qquad
 \Psi F_q^-(m)=\Psi(m)\circ(F_q^+)^{-1}.
\end{equation}
The inverse in the second formula specifies the variance of the dual.
In particular
\(E_{v'}(F_q^+\ell,F_q^-m)=E_v(\ell,m)\).
These identities apply to the local cochain translations used in
Kim's torsor functor \cite[\S2, pp.~5--6]{Kim2016}.
They require no choice of a torsor origin.

For global cocycles \(u,u'\), let \(u'-u=dz\).
The compatible boundary transport is \(q=rz\).
One has
\[
 (u',b+rz)-(u,b)=D(z,0),\quad
 F_{rz}^+\lambda_S(u)=\lambda_S(u'),\quad
 \Theta_{u'}F_{rz}^-=\Theta_u.
\]
The first identity proves trace compatibility. For the second, use
\(\beta+z\) as a primitive for \(u'\).
If another cochain \(z'\) has the same differential, then
\(z'-z\) is a global two-cocycle. Reciprocity gives
\(\sigma_\zeta(r(z'-z))=0\), so the two induced maps on
\(L_S^+\) and \(L_S^-\) agree.

This gives a precise functor on cocycles: its objects are the global
three-cocycles, and its arrows are two-cochains \(z\) modulo
\(dC^1(U,M)\), with \(dz=u'-u\). Composition is addition.
The constructions \(L_S^\pm\), \(\Psi_S\), \(\lambda_S\), and
\(\Theta\) respect this composition by the displayed identities.
Every automorphism acts trivially on the pushed-out torsors, since it
is a global two-cocycle. Different choices of representatives therefore
give canonically compatible torsor maps.

The associated-sheaf functor of a finite torsor is exact and preserves
tensor pairings, as can be checked on a cover trivializing the torsor.
An isomorphism of torsors gives a natural isomorphism of these functors
which is the identity on the trivial coefficient endpoints.
It therefore gives equal pulled-back Yoneda classes. Choose a two-cochain
\(z\) whose differential is the difference of the chosen
three-cocycles. The preceding independence from \(z\) makes the
maps on \(L_S^\pm\) independent of this choice. For a composite
of torsor isomorphisms, the sum of the two cochains has differential
equal to the difference for the composite. Thus the maps compose correctly.
The global point and its dual evaluation are natural under these maps.
Together with the cup comparison in Proposition~\ref{prop:do27-gluing},
this proves functorial compatibility with the orientation pairing.

We specify how this statement compares actual group cochains with
sheaf cochains. A finite continuous group cochain factors through a
finite quotient. A finite collection of cochains and their differential
identities therefore factors through one finite quotient \(Q\) of
\(\pi_1(U)\). The corresponding finite \'etale torsor gives the
exact functor \(T_Q\). Apply it to the coinduced bar resolution of
\(M\), and resolve its terms by the Godement functors. The resulting
first-quadrant double complex is a resolution of \(M\).
Its total differential is \(d_h+(-1)^p d_v\), with \(p\) the
horizontal degree. An invariant bar cochain maps to its associated
global section and then to Godement degree zero. The augmentation
makes its other differential zero, so this is a cochain map.
Passing through larger finite quotients gives the continuous-cochain
comparison. The same construction on each local field is the
field-site equivalence, so it is a quasi-isomorphism there.
All constructions use restriction of the finite torsor, so they give
compatible global and local maps through the common double resolutions.
No global \(K(\pi,1)\) assertion is needed.

For a square presented with a cochain homotopy, write its comparison
maps as \(F_U,F_\partial\) and fix a degree-minus-one map \(H\)
with
\[
 r'F_U-F_\partial r=dH+Hd.
\]
The induced primitive map is explicitly
\[
 (u,b)\longmapsto(F_Uu,F_\partial b+Hu).
\]
Its boundary equation follows by applying \(d\).
If \(d\beta=u\), then
\[
 F_\partial r\beta+Hu-r'F_U\beta=-dH\beta.
\]
Thus it carries \([r\beta,0]_+\) to the positive point represented
by \(r'F_U\beta\). It also intertwines \(\Psi\) and \(\Theta\),
since the local comparisons preserve the invariant maps.
This proves that Kim's global point, when represented by a continuous
group primitive, maps to~\eqref{eq:do28-state}. The argument needs
a comparison map, not a global cohomology isomorphism.

Finally replace \(\zeta\) by \(\zeta^a\), where
\(a\in(\mathbb Z/n)^\times\). Then
\(\sigma_{\zeta^a}=a\sigma_\zeta\) and
\(\tau_{u,\zeta^a}=a\tau_{u,\zeta}\).
The maps between the pushed-out torsors are
\[
 [b,t]_{\pm,\zeta}\longmapsto[b,at]_{\pm,\zeta^a}.
\]
They are equivariant for the automorphism \([a]\) of \(\Lambda\).
They carry the global points to the global points and multiply the
evaluation pairing by \(a\). The opposite action commutes with
\([a]\). Hence the dual comparison has the same primitive-root
normalization as the arithmetic orientation.

For a concrete sign test, take \(n=3\) and a boundary class \(\eta\)
with \(\sigma_\zeta\eta=1/3\). The positive boundary point changes
by \(+1/3\), its opposite changes by \(-1/3\), and
\(\tau_u(b+\eta)=\tau_u(b)-1/3\).
In the ramified example \(F=\mathbb Q(\zeta_9)\), with \(S\)
containing the primes above \(3\) and \(7\), take the torsor
\(y^3=7\) on \(U\). Its lift \(y^9=7\) is a \(C_9\)-torsor.
The ninth roots of unity lie in the unramified degree-three extension
of \(\mathbb Q_7\): \(9\mid7^3-1\), and Hensel's lemma lifts
their simple residue roots. Hence each prime over \(7\) in \(F\)
has \(v(7)=1\). Both displayed polynomials are Eisenstein there,
so these torsors are ramified at those primes. Use representatives
\(0,\ldots,8\) for \(C_9\), and write \(\bar i\) for the
residue of \(i\) in \(\{0,1,2\}\). On \(C_3\), let
\[
 c(i,j,k)=i\left\lfloor\frac{j+k}{3}\right\rfloor\pmod3.
\]
Its pullback to \(C_9\) has the primitive
\(B(i,j)=\bar i\lfloor j/3\rfloor\pmod3\).
Indeed, expanding \(dB\) cancels the quotients of \(j\) and \(k\),
leaving \(\bar i\lfloor(\bar j+\bar k)/3\rfloor\).
Pull back \(B\) to a global primitive \(\beta\).
For \(b=r\beta\), the global point is \([b,0]_+\) and the trace
is zero. Replacing one local primitive by \(b+\eta\), with invariant
\(1/3\), leaves the global point fixed and changes its coordinate
to \(-1/3\). This tests the opposite-torsor comparison on a fixed
ramified arithmetic torsor.

\subsection{Restriction and refinement of the support}

Let \(S\subset T\) be nonempty finite sets as above, and put
\(U_T=X\setminus T\). There is a cochain map
\(R_{ST}:A_S(\mathcal F)\to A_T(\mathcal F)\).
Restrict the global component from \(U\) to \(U_T\).
Keep the old local components. At \(w\in T\setminus S\), set
\[
 a_w=u|_{D_w},\qquad b_w=0.
\]
The square of restriction maps commutes, so this formula commutes
with \(D\). Formula~\eqref{eq:do27-cup} shows that it preserves
the cup pairing. Moreover
\[
 R_{ST}J_S=J_T,\qquad R_{TV}R_{ST}=R_{SV}\quad(S\subset T\subset V).
\]
The same formulas give maps for completed local terms.

The trace is compatible on cohomology. To see this directly, let
\((u,a,b)\) be a degree-three \(\mathbb G_m\)-cocycle and choose
\(dh=u\). At a new prime choose the local primitive
\(k_w=h|_{D_w}\). Its contribution to~\eqref{eq:do27-trace} is
\(r_wh-s_wk_w=0\). The old contributions are unchanged. Hence
\[
 T_T H^3(R_{ST})=T_S,
 \qquad \widehat T_T H^3(\widehat R_{ST})=\widehat T_S.
\]

For a degree-three class in the finite compact cone, first map it
to the gluing complex and then restrict. A new local ring component
is \(u|_{\widehat R_w}\). Choose \(a_w\) with
\(da_w=u|_{\widehat R_w}\) and remove it as above.
The new boundary primitive is \(b_w=s_wa_w\).
Its ambiguity is zero in cohomology because
\(H^2(\widehat R_w,M)=0\). This gives a canonical refinement map
in degree three, and it preserves~\eqref{eq:do27-compact-trace}.
For a torsor unramified at the new prime, these are precisely the
primitives supplied by its local ring extension.

The same refinement acts on the positive and negative boundary torsors.
For \(w\in T\setminus S\), the local primitive class
\([a_w]\), with \(da_w=u|_{\widehat R_w}\), is unique modulo
coboundaries. Its generic restriction is a distinguished boundary
point. Define
\[
 \mathcal R_{ST}^\pm[b,t]_\pm
       =[(b,(s_wa_w)_{w\in T\setminus S}),t]_\pm.
\]
If \(a_w\) changes by a closed ring cochain, its class in degree
two is zero. Thus the formula is independent of the chosen primitives.
On the old acting group, the new sum map is \(\sigma_S\), so the
formula respects both pushout relations. It is an equivariant
isomorphism of \(\Lambda\)-torsors.

For a global primitive \(\beta\) on \(U\), choose
\(a_w=\beta|_{\widehat R_w}\). This proves
\[
 \mathcal R_{ST}^+\lambda_S(u)=\lambda_T(u|_{U_T}),\qquad
 \Theta_{u|_{U_T}}\mathcal R_{ST}^-=\Theta_u.
\]
Moreover
\[
 \Psi_T\mathcal R_{ST}^-(m)
    =\Psi_S(m)\circ(\mathcal R_{ST}^+)^{-1}.
\]
The last identity follows by subtracting two refined primitive tuples:
their new components have difference zero. Hence the evaluation pairing
is unchanged. Refinements compose, since adjoining these unique local
ring primitive classes in either order gives the same tuple.
For \(u'=u+dz\), choose the new ring primitive \(a_w+z|_{\widehat R_w}\).
This also proves commutation with the global cocycle transports.
Thus support refinement commutes with the orientation functor, its
positive boundary point, and its opposite and dual torsors.

Changing an \'etale neighborhood representative refines the filtered
system defining \(C(D_v)\) and \(C(K_v)\). Equality in those
colimits means equality after a common refinement. Every formula
uses only finitely many components and finitely many cofaces.
It therefore descends to these colimits and commutes with such
refinement. A change of the cochain resolution gives the same
cohomological maps through the common double resolution and the
fixed connecting-map convention.

Local unramified primitive classes, the global boundary point, and a
prepared state have different carriers. The local primitives belong
to quotients of degree-two cochains on the local rings. The global
point \(\lambda_S(u)\) belongs to the finite
\(\Lambda\)-torsor \(L_S^+(u)\) and determines the nonzero vector
\([\lambda_S(u),1]\) in its associated complex line. Interpreting this
vector as a prepared state requires a specified quantum state space and
a preparation map compatible with that theory's gluing operations.
The finite torsor construction supplies no such quantum comparison.

The finite adelic cochain trace therefore supplies the oriented scalar
for an unramified finite torsor and the difference coordinate for a
ramified torsor with a chosen boundary point. A comparison
with integration of Schwartz--Bruhat functions, a Mellin transform,
or a determinant requires a map from the arithmetic coefficient
complex to that analytic carrier. No such map is part of the
finite gluing theorem. A continuous arithmetic Chern--Simons theory
would additionally require its continuous gauge group, coefficient
complex, topology, and compatible gluing and state maps. The finite
sums above imply no vanishing statement for an infinite product of
cohomology groups.

\section{Arithmetic state and determinant comparisons}
The finite boundary point is invariant under the changes proved above.
It is an element of a torsor, with an affine evaluation pairing, and
gives the canonical associated-line vector $[\lambda_S(u),1]$.
A prepared quantum state would additionally require a specified space of
observables and a functional on it. The arithmetic cohomology, the
positive Weil sectors and the analytic determinants therefore leave
precise comparison problems.

\begin{openproblem}[From the classical invariant to a chiral determinant]
\label{problem:acs}
A comparison with Section~\ref{sec:arithmetic-comparison} and
Theorem~\ref{thm:deninger} requires the following constructions.

\begin{enumerate}
\item Fix one cohomological framework: \(X_{\mathrm{et}}\) under the totally
imaginary hypotheses of Theorem~\ref{thm:kim}; ordinary small-site
cohomology with odd-order coefficients when real places are present;
explicitly defined modified cohomology with Tate-cohomological real terms;
or ordinary derived cohomology on \(\widetilde X_{\mathrm{et}}\) with a
specified constructible torsion sheaf, an extension to the site, and
infinite-place specialization maps.  Specify separately any ramification
set and its boundary torsors.

\item For a continuous or non-finite gauge object, define an arithmetic
field groupoid, an action functional, measures or finite sums, state spaces,
gluing maps, and anomaly data.  Assemble them into a symmetric monoidal
functor on a specified arithmetic bordism category.  Kim's finite classical
function does not provide these objects.

\item Construct complexes \(C_{\mathrm{Ar}}\), \(C_{\mathrm{Den}}\), and
\(C_{\mathrm{ACS}}\) in one named derived category and chain maps
\[
 C_{\mathrm{Ar}}\longrightarrow C_{\mathrm{Den}}
 \longrightarrow C_{\mathrm{ACS}}
\]
with stated degrees, completions, quasi-isomorphism hypotheses, and
intertwining relations for the relevant operators.  The first complex must
in particular solve Open Problem~\ref{problem:bar}, and the middle complex
must supply the data assumed in Theorem~\ref{thm:deninger}.

\item To compare a partition object with \(\xi_{\mathbb R}(s)\), construct
the complex parameter \(s\), a holomorphic or meromorphic family of
partition objects, and a proof of the determinant identity.  A single
closed-theory partition number cannot equal a nonconstant function of an
undeclared parameter.

\item A Wilson surface requires higher gauge data: at minimum a crossed
module or \(2\)-group, a \(2\)-connection or arithmetic replacement, and a
surface-holonomy/transgression construction.  It is not defined by an
ordinary finite \(A\)-torsor or by Theorem~\ref{thm:kim}.

\item Any extension to a completed motivic \(L\)-function must begin with
a specified motive, its local factors and archimedean factor, its weight,
and a no-cancellation determinant comparison.  A genus label alone supplies
none of these data.
\end{enumerate}
\end{openproblem}

\section{The comparisons and their remaining extensions}

The common local factor determines a scalar after the measure or
operator normalization is fixed. The rational Tate construction
assembles these scalars and proves their functional equation. The
equilibrium construction specifies positive states, with a phase
simplex above the critical inverse temperature and uniqueness at
and below it. Its common Toeplitz restriction leaves the polynomial
augmentation distinct from both a thermal state and the vacuum.

For the prime determinants, the continued objects are normalized germs
on explicit Riemann surfaces. The local exponents are total divisor
contributions, and descent is controlled by periods on all closed
curves. On a smaller positive half-plane, several orders have a
connected simultaneous cover determined by the vector of their
period characters. The reciprocal-real poles generate its full image;
no assumption on coincidences or simplicity of nonreal zeta zeros is
needed. The degree-twenty-four surface for orders two and three is
an explicit case where a connected component of a normalized fibre
product carries more information than either coordinate degree.
For higher orders the unbounded reciprocal-prime ramification
obstructs a finite realization on the whole positive half-plane and
near zero. It establishes no natural-boundary theorem along the
entire imaginary axis.

The formal coefficient algebra and the analytic determinant agree
through an actual evaluation homomorphism on a weighted Dirichlet
Banach algebra. The sharp threshold is the same as the Schatten
threshold, while scalar continuation leaves this coefficient domain
unchanged. The row-finite formal operator algebra admits neither a
continuous extension of finite-rank trace nor an algebraic cyclic
extension with its rank-one normalization. These conclusions concern
the specified recipients. They leave open comparison maps using a
different coefficient category, representation or trace ideal.

The integral divisor bar retains multiplication, augmentation,
coalgebra maps and every finite prime-set inclusion. The torsion
Hall construction supplies a separate geometric product and its
normalized adjoints. The arithmetic finite-boundary construction
supplies a degree-three trace, a global torsor point and finite
gluing, with the positive divisor orientation. A bare cone boundary
has the negative connecting sign. These facts distinguish a
cohomological boundary point from local primitives and from a
prepared state.

Three further constructions would change the arithmetic consequences.
First, the map $\tau_S$ of Open Problem~\ref{problem:bar} must satisfy
its Maurer--Cartan equation, converge in the declared filtration and
commute with enlargement of the set of places. This would turn the
integral divisor model into an arithmetic twisted complex. Second,
an independently defined positive representation must recover the
full Weil functional on a dense admissible convolution domain;
the positive narrow-window subspaces do not give that map. Third,
an arithmetic cohomology theory must produce the operators and exact
spectral multiplicities in Theorem~\ref{thm:deninger} and
Corollary~\ref{cor:hp}. Its degree-one determinant would then equal
$\xi_{\mathbb R}(s)/(2\sqrt2\pi^2)$ with the normalization proved here.
Finite arithmetic gluing supplies no continuous gauge measure or
determinant map; those additional state and comparison data are
specified in Open Problem~\ref{problem:acs}.

\begin{thebibliography}{99}

\bibitem{Bienenfeld1987}
M.~Bienenfeld,
\emph{An \'etale cohomology duality theorem for number fields with a real
embedding}, Trans. Amer. Math. Soc. 303 (1987), no.~1, 71--96.

\bibitem{BostConnes1995}
J.-B.~Bost and A.~Connes,
\emph{Hecke algebras, type III factors and phase transitions with
spontaneous symmetry breaking in number theory}, Selecta Math. 1 (1995),
no.~3, 411--457.

\bibitem{Deninger1998}
C.~Deninger,
\emph{Some analogies between number theory and dynamical systems on
foliated spaces}, in Proceedings of the ICM 1998, Vol.~I, 163--186.

\bibitem{Kim2016}
M.~Kim,
\emph{Arithmetic Chern--Simons theory I}, arXiv:1510.05818v4 (2016).

\bibitem{LurieHA}
J.~Lurie,
\emph{Higher Algebra}, 2017, \S4.4.2,
\texttt{https://www.math.ias.edu/{\string~}lurie/papers/HA.pdf}.

\bibitem{Mazur1973}
B.~Mazur,
\emph{Notes on \'etale cohomology of number fields}, Ann. Sci. \'Ecole
Norm. Sup. (4) 6 (1973), 521--552.

\bibitem{Tate1950}
J.~T.~Tate,
\emph{Fourier analysis in number fields and Hecke's zeta-functions},
Ph.D. thesis, Princeton University, 1950.

\bibitem{Weil1952}
A.~Weil,
\emph{Sur les ``formules explicites'' de la th\'eorie des nombres
premiers}, Meddelanden fr\aa n Lunds Univ. Mat. Sem., supplement (1952),
252--265.

\bibitem{Riemann1859}
B.~Riemann,
\emph{Ueber die Anzahl der Primzahlen unter einer gegebenen Gr\"osse},
Monatsberichte der Berliner Akademie, November 1859, pp.~671--680.
German transcription by D.~R.~Wilkins, p.~3,
\url{https://www.maths.tcd.ie/pub/HistMath/People/Riemann/Zeta/Zeta.pdf}.

\bibitem{do27-Limits}
The Stacks Project Authors, \emph{Colimits}, Theorem 59.51.3,
\href{https://stacks.math.columbia.edu/tag/03Q4}{Tag 03Q4}.

\bibitem{do27-Godement}
The Stacks Project Authors, \emph{Godement resolution}, Section 20.30,
\href{https://stacks.math.columbia.edu/tag/0FKR}{Tag 0FKR}.

\bibitem{do27-Field}
The Stacks Project Authors, \emph{Cohomology of a point},
Lemmas 59.59.1--59.59.2,
\href{https://stacks.math.columbia.edu/tag/03QQ}{Tag 03QQ}.

\bibitem{do26-StacksFundamental}
The Stacks Project Authors,
\emph{Fundamental groups}, Theorem 58.6.2,
\href{https://stacks.math.columbia.edu/tag/0BQ8}{Tag 0BQ8}.

\bibitem{SchiffmannHall}
O.~Schiffmann, \emph{Lectures on Hall algebras},
arXiv:math/0611617, Theorem~2.6 and equation~(2.9).

\end{thebibliography}

\end{document}
